\documentclass[11pt,a4paper]{article}
\usepackage[utf8]{inputenc}
\usepackage[T1]{fontenc}
\usepackage[czech]{babel}
\usepackage[margin=2.2cm]{geometry}
\usepackage{amsmath,amssymb,amsthm}
\usepackage{parskip}
\usepackage{enumitem}
\setlist{nosep,leftmargin=1.4em}
\usepackage{booktabs}
\usepackage{array}
\usepackage{xcolor}
\definecolor{linkblue}{RGB}{20,60,150}
\definecolor{defbg}{RGB}{240,244,252}
\definecolor{defframe}{RGB}{100,130,190}
\definecolor{markcol}{RGB}{176,84,0}
\definecolor{markbg}{RGB}{253,246,236}
\usepackage[most]{tcolorbox}
\newtcolorbox{defbox}[1][]{breakable,colback=defbg,colframe=defframe,boxrule=0.6pt,arc=1.5mm,left=2mm,right=2mm,top=1mm,bottom=1mm,title={#1},fonttitle=\bfseries}
\newtcolorbox{poznbox}[1][]{breakable,colback=markbg,colframe=markcol,boxrule=0.6pt,arc=1.5mm,left=2mm,right=2mm,top=1mm,bottom=1mm,title={#1},fonttitle=\bfseries}
\newtheorem{veta}{Věta}[section]
\theoremstyle{definition}
\newtheorem{definice}{Definice}[section]
\usepackage[colorlinks=true,linkcolor=linkblue,urlcolor=linkblue]{hyperref}
\usepackage{algorithm}
\usepackage{algpseudocode}

\floatname{algorithm}{Algoritmus}
\renewcommand{\algorithmicrequire}{\textbf{Vstup:}}
\renewcommand{\algorithmicensure}{\textbf{Výstup:}}

% Značka pro látku, která není ve slidech NDBI023 / NAIL116.
\newcommand{\mimo}{\textcolor{markcol}{\textsf{\footnotesize[$\blacklozenge$~nad rámec slidů]}}}
\newcommand{\mimoq}[1]{\textcolor{markcol}{\textsf{\footnotesize[$\blacklozenge$~nad rámec slidů: #1]}}}

\title{\textbf{Dobývání znalostí}\\[2mm]
\large Učební text k~okruhu státní závěrečné zkoušky}
\author{SZZ Umělá inteligence, MFF UK}
\date{srpen 2026}

\begin{document}
\maketitle

\begin{center}
\begin{minipage}{0.92\textwidth}
\small
Text pokrývá okruh \emph{Dobývání znalostí} magisterské SZZ oboru Umělá inteligence.
Členění je záměrně přímočaré: kapitoly odpovídají \textbf{jedna ku jedné} větám oficiálního
znění okruhu, takže se u~zkoušky dá postupovat rovnou podle nich. Přehled dává mapovací
tabulka níže.

Výklad vychází ze slidů dvou přednášek I.~Mrázové na MFF UK, totiž \textbf{NDBI023 Dobývání
znalostí} a~\textbf{NAIL116 Sociální sítě a~jejich analýza}. Tam, kde slidy samy nestačí, je
doplněn o~standardní látku z~učebnic: Han, Kamber, Pei: \emph{Data Mining --- Concepts and
Techniques}, Tan, Steinbach, Kumar: \emph{Introduction to Data Mining} a~Aggarwal:
\emph{Data Mining --- The Textbook}; pro sítě pak Easley, Kleinberg: \emph{Networks, Crowds
and Markets}, Newman: \emph{Networks} a~Zafarani, Abbasi, Liu: \emph{Social Media Mining}.
Vedle definic a~pseudokódů obsahuje text také odvození, číselné příklady a~srovnávací
tabulky. Každou kapitolu uzavírá shrnutí pro ústní zkoušku.
\end{minipage}
\end{center}

\begin{center}
\begin{minipage}{0.92\textwidth}
\begin{poznbox}[Rozsah textu a~značení]
\small
Rozsah jednotlivých kapitol odpovídá váze látky. Nejvíce prostoru dostávají metody dobývání
znalostí (kap.~\ref{sec:metody}) a~analýza sociálních sítí (kap.~\ref{sec:sna}); právě těmto
dvěma tématům se obě přednášky věnují nejpodrobněji a~právě je okruh jmenuje nejkonkrétněji.

Značka \mimo{} upozorňuje na látku, která \textbf{není ve slidech} ani jednoho z~obou kurzů.
Slidy jsou přitom k~dispozici kompletní, takže nejde o~mezeru v~podkladech, ale o~látku, která
se v~nich opravdu neprobírá. V~textu přesto zůstává, a~to ze dvou důvodů: buď ji
\emph{jmenuje oficiální znění okruhu}, nebo bez ní nedává okolní výklad smysl.

Značka je proto zároveň nabídkou ke~škrtnutí: co označeno není, je ve slidech doloženo.
Následuje-li za značkou dvojtečka, upřesňuje, čeho přesně se výhrada týká. Typicky jde
o~případ, kdy souhrnná tabulka nebo rámec jsou moje, ale jednotlivé položky v~nich ze slidů
pocházejí.

Ze slidů NDBI023 pochází analýza dat a~shlukování, asociační pravidla, rozhodovací stromy
a~ensembly, SVM spolu s~vyhodnocováním klasifikátorů, bayesovské modely a~ELM. NAIL116 k~tomu
přidává základní pojmy, modelování vlivu, homofilii a~afiliaci, analýzu vazeb a~komplexní sítě
včetně detekce komunit a~predikce vazeb. Nejnápadnější mezery jsou dvě:
\textbf{proces KDD a~metodologie} (kap.~\ref{sec:paradigmata}) a~\textbf{charakteristická
a~diskriminační pravidla} (kap.~\ref{sec:pravidla}). Okruh je jmenuje výslovně, ale
\emph{ani jeden z~kurzů je neprobírá}, a~jsou tedy zpracovány podle Han--Kamber.
\end{poznbox}
\end{minipage}
\end{center}

\vspace{2mm}

\section*{Mapování okruhu na kapitoly}
\addcontentsline{toc}{section}{Mapování okruhu na kapitoly}

\begin{center}
\small
\begin{tabular}{@{}>{\raggedright\arraybackslash}p{8.6cm}>{\raggedright\arraybackslash}p{7.4cm}@{}}
\toprule
\textbf{Věta oficiálního znění okruhu} & \textbf{Kapitola v~tomto textu} \\
\midrule
Základní paradigmata dobývání znalostí.
  & \ref{sec:paradigmata}~Základní paradigmata dobývání znalostí \\
Příprava dat; výběr atributů a~metody pro analýzu jejich relevance.
  & \ref{sec:priprava}~Příprava dat, výběr atributů a~analýza jejich relevance \\
Metody pro dobývání znalostí; asociační pravidla, přístupy založené na principu učení
s~učitelem a~klastrová analýza.
  & \ref{sec:metody}~Metody pro dobývání znalostí (\ref{sec:metody}.1~asociační pravidla,
    \ref{sec:metody}.2--\ref{sec:metody}.3~učení s~učitelem, \ref{sec:metody}.4~klastrová
    analýza) \\
Metody pro extrakci charakteristických diskriminačních pravidel a~měření jejich
zajímavosti.
  & \ref{sec:pravidla}~Charakteristická a~diskriminační pravidla a~jejich zajímavost \\
Reprezentace, vyhodnocování a~vizualizace získaných znalostí.
  & \ref{sec:reprez}~Reprezentace, vyhodnocování a~vizualizace znalostí \\
Modely pro analýzu sociálních sítí; míry centrality, detekce komunit.
  & \ref{sec:sna}~Modely pro analýzu sociálních sítí, centralita a~detekce komunit \\
Praktické využití technik pro dobývání znalostí a~analýzu sociálních sítí.
  & \ref{sec:praxe}~Praktické využití technik \\
\bottomrule
\end{tabular}
\end{center}

\vspace{1mm}
\noindent\small\textit{Poznámka k~terminologii.} Text používá české termíny; anglický
ekvivalent je uveden kurzívou v~závorce při prvním výskytu, protože obě přednášky i~celá
literatura jsou anglicky.\normalsize

\vspace{2mm}
\tableofcontents
\newpage

%=====================================================================
\section{Základní paradigmata dobývání znalostí}
\label{sec:paradigmata}
%=====================================================================

Tato kapitola vymezuje obor, popisuje jeho proces a~zavádí základní dělení úloh a~dat.
Je to rámec, do něhož se zasazují všechny kapitoly následující. Bez něj by nebylo zřejmé,
proč si příprava dat zaslouží samostatnou kapitolu a~proč se metody dělí zrovna na
asociační pravidla, učení s~učitelem a~shlukování.

\subsection{Motivace a~vymezení oboru}

Organizace shromažďují data v~objemech, které daleko přesahují kapacitu člověka je vůbec
prohlédnout. Klasické dotazování, tedy SQL a~OLAP, si s~tím poradí jen zčásti: umí
spolehlivě odpovědět na otázku, kterou už umíme položit. \emph{Dobývání znalostí} míří
jinam: odpovídá na otázky, které jsme položit neuměli, protože hledá
v~datech struktury, o~jejichž existenci jsme předem nevěděli.

Dobře to ilustruje jeden příklad. Systém prošel 20 milionů výkonů účtovaných zdravotní
pojišťovně, 214 tisíc z~nich označil za podivné a~14 tisíc nejpodezřelejších předal
k~ruční kontrole. Revizní lékaři jsou schopni prohlédnout řádově stovky případů; z~těch,
které jim systém předložil, jich okamžitě 8 tisíc označili za podvodné. Stroj tady experta
nenahradil, jen \emph{zaostřil} jeho pozornost.

\begin{defbox}[Dobývání znalostí z~databází]
\emph{Dobývání znalostí z~databází} (\emph{Knowledge Discovery in Databases}, KDD) je
netriviální proces získávání implicitních, dosud neznámých a~potenciálně užitečných
informací z~dat.

Klíčová jsou všechna čtyři slova:
\begin{itemize}
\item \textbf{Netriviální} vylučuje prostý dotaz i~agregaci: kdyby výsledek stačilo
      z~databáze vyčíst, o~dobývání znalostí by nešlo.
\item \textbf{Implicitní} znamená, že znalost není v~datech nikde explicitně zapsaná,
      leží v~nich jen jako struktura.
\item \textbf{Dosud neznámá} vyřazuje závěry, které každý zná předem; potvrzená
      samozřejmost je bezcenná.
\item \textbf{Potenciálně užitečná} požaduje, aby nález vedl k~akci nebo k~rozhodnutí.
\end{itemize}
\end{defbox}

Termín \emph{dobývání dat} (\emph{data mining}, DM) se v~úzkém pojetí vztahuje pouze
k~jedinému kroku procesu KDD, totiž k~vlastní aplikaci analytických algoritmů. Běžná praxe
tento rozdíl nedodržuje a~oba pojmy používá zaměnitelně. Jako samostatný obor se dobývání
znalostí konstituovalo v~90.~letech 20.~století, a~to na průsečíku tří starších disciplín:

\begin{itemize}
\item \textbf{Umělá inteligence} dodala především metody strojového učení, tedy indukci
      pravidel, stromy a~neuronové sítě.
\item \textbf{Databázové systémy} se starají o~ukládání velkých datových souborů
      a~efektivní přístup k~nim; odtud pocházejí datové sklady i~techniky vyhledávání
      informací.
\item \textbf{Statistika} přispěla modelováním a~analýzou závislostí nalezených v~datech
      a~testováním hypotéz.
\end{itemize}

Těmi třemi ale výčet nekončí. Čtvrtou složkou, v~praxi ovšem rozhodující, je nutnost
využít zpracovaná data k~podpoře (strategického) rozhodování v~podniku. Bez ní se dobývání
znalostí redukuje na akademické cvičení.

\textbf{Vztah k~datovým skladům a~OLAP.} \mimo{} Typickým zdrojem dat pro KDD bývá
\emph{datový sklad} (\emph{data warehouse}), tedy integrovaná, subjektově orientovaná,
časově variantní a~stálá databáze konsolidovaná z~provozních systémů. Data se v~něm
organizují do \emph{datové kostky}, kterou tvoří \emph{dimenze} (čas, produkt, region,
\dots) $\times$ \emph{míry} (tržby, počty kusů). V~relační implementaci kostce odpovídá
\emph{hvězdicové schéma}, v~němž je tabulka faktů obklopena tabulkami dimenzí; fyzicky se
data ukládají buď multidimenzionálně (MOLAP), nebo relačně (ROLAP).

Nad kostkou pracuje \textbf{OLAP} (\emph{On-Line Analytical Processing}) a~umožňuje
agregační operace \emph{roll-up} a~\emph{drill-down} a~řezy \emph{slice \& dice}. I~on ale
odpovídá jen na otázky, které si analytik umí \emph{předem položit}. Dobývání znalostí jde
o~krok dál a~hledá vzory, o~nichž předem nevíme, že v~datech jsou. Souvislost generalizace
v~OLAP s~indukcí orientovanou na atributy viz kap.~\ref{sec:pravidla}.

\subsection{Proces KDD}

\mimoq{pětifázový rámec jako celek; náplň jednotlivých fází ve slidech je}
Dobývání znalostí není přímočará posloupnost kroků, ale \textbf{interaktivní
a~iterativní} proces: výsledky jednoho kroku běžně vedou k~návratu ke krokům předchozím.
Klasické schéma (Fayyad et al.) rozděluje celý postup do pěti fází:

\begin{center}
\small
\begin{tabular}{@{}p{3.3cm}p{4.2cm}p{7.6cm}@{}}
\toprule
\textbf{Fáze} & \textbf{Vstup $\to$ výstup} & \textbf{Náplň} \\
\midrule
Výběr (\emph{selection}) & původní data $\to$ vybraná data & výběr relevantních záznamů
  a~atributů z~datového skladu \\
Předzpracování (\emph{preprocessing}) & vybraná data $\to$ předzpracovaná data & čištění,
  chybějící hodnoty, šum, odlehlé hodnoty, integrace zdrojů \\
Transformace (\emph{transformation}) & předzpracovaná data $\to$ transformovaná data &
  normalizace, diskretizace, odvozené atributy, redukce dimenze \\
Vlastní dobývání (\emph{data mining}) & transformovaná data $\to$ vzory & aplikace
  algoritmů (stromy, asociace, klastry, \dots) \\
Interpretace (\emph{interpretation}) & vzory $\to$ znalost & vyhodnocení z~pohledu
  koncového uživatele, vizualizace, nasazení \\
\bottomrule
\end{tabular}
\end{center}

Cílem prvních tří fází, souhrnně nazývaných \emph{příprava dat} (kap.~\ref{sec:priprava}),
je z~dat uložených ve složité struktuře vybudovat \textbf{jednu tabulku} obsahující
relevantní údaje o~zkoumaných objektech, tedy o~klientech banky, o~zákaznících a~podobně.
Zní to jako pouhá administrativa, ale právě tento krok typicky spotřebuje 60--80~\%
celkového času projektu.

\textbf{Sedm kroků očima manažera.} \mimo{} Na tentýž proces se lze podívat i~organizačně;
impulsem je pak \uv{aktuální problém} a~výstupem \uv{znalost k~jeho řešení}.
\emph{(1)}~Nejprve se sestaví tým odborníků, v~němž jsou zastoupeni experti na oblast, na
data a~na metody KDD; žádná z~těchto rolí není zastupitelná. \emph{(2)}~Pak se specifikuje
problém, což znamená převést obchodní zadání do kontextu dobývání znalostí: \uv{zvýšit
zisk} $\to$ \uv{predikovat pravděpodobnost odchodu zákazníka}.
\emph{(3)}~Následuje sběr všech dostupných dat včetně externích, která popisují prostředí
(roční období, reklamní kampaně, počasí); tento krok může vést až k~přeformulování
problému. \emph{(4)}~Teprve potom se volí metoda či kombinace metod a~k~nim metody
vizualizační.

\emph{(5)}~Data se předzpracují. \emph{(6)}~Vlastní dolování probíhá opakovaně, protože
výsledky předchozích běhů ovlivňují parametry běhů následujících. \emph{(7)}~Nakonec se
výsledky interpretují. Část z~nich bývá nezajímavá nebo triviální, proto se hodí shrnout
je do analytické zprávy; výstupem ovšem může být i~provedení konkrétní akce.

\subsection{Typy úloh dobývání znalostí}

\mimo{} Jakou znalost od dobývání vlastně čekáme, není samozřejmé: jednou má být přesná,
jindy srozumitelná a~potřetí především nová. Úlohy se proto dělí především podle toho,
jaký vztah má nalezená znalost k~cílovému konceptu:

\begin{center}
\small
\begin{tabular}{@{}p{2.9cm}p{6.0cm}p{6.4cm}@{}}
\toprule
\textbf{Typ úlohy} & \textbf{Cíl} & \textbf{Charakter výsledku} \\
\midrule
Klasifikace / predikce & zařadit nové případy, resp.\ odhadnout budoucí vývoj entity &
  preferujeme \emph{přesnost pokrytí} před jednoduchostí; hodně znalosti, obtížně
  interpretovatelné \\
Popis (charakterizace, profil) & najít dominantní strukturu či vztahy v~datech &
  požadujeme \emph{srozumitelnost} při plném pokrytí konceptu; méně znalosti a~méně přesné \\
Hledání \uv{nugetů} & najít zajímavou (novou, překvapivou) znalost & znalost \emph{nemusí}
  pokrývat celý koncept; hodnotí se překvapivost \\
\bottomrule
\end{tabular}
\end{center}

Rozdíl je nejlépe vidět na příkladech. Predikcí je předpověď počasí, ceny akcií nebo
spotřeby elektřiny. Popisem je třeba profil bonitního klienta. A~nugetem je neobvyklá
kombinace zboží v~košíku či podezřelý vzorec transakcí.

\begin{defbox}[Deskriptivní vs.\ prediktivní úlohy]
\textbf{Deskriptivní} (\emph{descriptive}) úlohy charakterizují obecné vlastnosti dat.
Patří k~nim charakterizace a~diskriminace, asociační a~korelační analýza, shlukování,
detekce odlehlých hodnot a~analýza vývoje.

\textbf{Prediktivní} (\emph{predictive}) úlohy naproti tomu provádějí nad daty inferenci
za účelem predikce. Jsou to klasifikace s~diskrétním cílovým atributem, regrese se
spojitým cílovým atributem a~predikce časových řad.

Ortogonálně k~tomu se úlohy třídí podle dostupnosti cílové proměnné. Při \emph{učení
s~učitelem} (\emph{supervised}) cílová proměnná je, což je případ klasifikace a~regrese.
Při \emph{učení bez učitele} (\emph{unsupervised}) není, a~sem spadá shlukování, asociační
pravidla a~detekce odlehlých hodnot. \emph{Poloučené} úlohy (\emph{semi-supervised}) leží
mezi obojím: malá část dat je označkovaná a~velká nikoli.
\end{defbox}

Právě tímto dělením se řídí i~kap.~\ref{sec:metody}: okruh v~ní jmenuje jednoho
reprezentanta učení bez učitele s~transakčními daty, totiž asociační pravidla, dále učení
s~učitelem a~konečně učení bez učitele s~metrickými daty, tedy klastrovou analýzu.

\subsection{Typy dat}

Metodu nelze volit bez ohledu na to, jak data vypadají. Třídí se ze dvou hledisek. Prvním
z~nich je \textbf{typ atributu}, tedy škála, na níž se hodnoty měří; ta určuje, jaké
operace jsou nad hodnotami vůbec smysluplné, a~tím i~to, jaké metody lze použít:

\begin{center}
\small
\begin{tabular}{@{}p{3.3cm}p{2.2cm}p{5.2cm}p{4.0cm}@{}}
\toprule
\textbf{Škála} & \textbf{Operace} & \textbf{Příklad} & \textbf{Statistika} \\
\midrule
Nominální (\emph{nominal}) & $=$, $\neq$ & barva, druh ovoce & modus, četnosti \\
Ordinální (\emph{ordinal}) & $=$, $<$ & známka, \{nízký, střední, vysoký\} & medián,
  kvantily \\
Intervalová (\emph{interval-scaled}) & $+$, $-$ & teplota ve $^{\circ}$C, datum & průměr,
  směr.\ odchylka \\
Poměrová (\emph{ratio-scaled}) & $\times$, $/$ & příjem, hmotnost, četnost & geometrický
  průměr \\
\bottomrule
\end{tabular}
\end{center}

U~intervalové škály je klíčové, že \emph{intervaly mají v~celém rozsahu stejnou
důležitost}: rozdíl mezi věkem 10 a~20 let je týž jako rozdíl mezi 40 a~50 lety. Poměrová
škála má navíc absolutní nulu, takže u~ní dává smysl i~poměr dvou hodnot. Poměrové
atributy s~exponenciálním chováním, jako je počet mikroorganismů $Ae^{Bt}$, se před dalším
zpracováním logaritmují a~pak se s~nimi zachází jako s~intervalovými.

Pro numerické metody se nečíselné škály převádějí. \emph{Nominální} atribut o~$v$
hodnotách se rozepíše na $v$ binárních atributů, z~nichž odpovídající kategorie dostane
hodnotu 1 a~ostatní 0. \emph{Ordinální} atribut se běžně zpracovává jako intervalový.

Druhým hlediskem je \textbf{struktura datového souboru}, tedy to, co vlastně tvoří jeden
záznam:
\begin{itemize}
\item \textbf{Tabulková, čili relační data} mají tvar objekty $\times$ atributy a~jsou
      standardním vstupem většiny metod.
\item \textbf{Transakční data} mají záznam proměnné délky, typickým příkladem je nákupní
      košík. Vstupují do asociačních pravidel.
\item \textbf{Temporální data} (\emph{temporal}) jsou třeba časové řady kurzů akcií.
      Typickou úlohou nad nimi je predikce budoucích hodnot a~řada se kvůli ní
      transformuje do tabulky posuvným oknem: vstupy $y(t_0),y(t_1),y(t_2),y(t_3)$ $\to$
      výstup $y(t_4)$; $y(t_1),\dots,y(t_4)\to y(t_5)$ atd.
\item \textbf{Prostorová data} (\emph{spatial}) pocházejí z~geografických informačních
      systémů. Vyznačují se implicitní relací sousedství, tedy tím, že hodnoty atributů
      sousedních objektů se od sebe příliš neliší; užitečná je zejména pro interpretaci.
\item \textbf{Strukturní data} (\emph{structural}) popisují objekt jeho stavbou; jsou to
      třeba chemické sloučeniny zapsané řetězci SMILES nebo soubory faktů v~Prologu.
      Úlohou pak bývá klasifikace látek podle struktury.
\item \textbf{Textová a~webová data, multimédia, grafy a~sítě} vyžadují specifickou
      reprezentaci, ať už vektorový model, nebo grafové modely (viz kap.~\ref{sec:sna}).
\end{itemize}

\subsection{Metodologie dobývání znalostí}
\label{ssec:metodologie}

\mimo{} Metodologie má poskytnout jednotný rámec a~vést aplikace bez ohledu na obor
podnikání, a~umožnit tak, aby se zkušenosti z~úspěšných projektů daly sdílet a~přenášet.
Jednotlivé metodologie vznikají dvěma cestami. Buď u~velkých výrobců BI nástrojů, odkud
pochází SEMMA od SAS Institute a~5A od SPSS, nebo jako neutrální, softwarově nezávislé
standardy; zástupcem druhé cesty je CRISP-DM.

\textbf{SEMMA} (SAS Enterprise Miner) má pět fází. \emph{Sample} vybere data pro
modelování, tedy vzorkuje, imputuje a~dělí na trénovací, testovací a~validační množinu.
\emph{Explore} data vizuálně prozkoumá a~zredukuje; sem patří odlehlé hodnoty,
vícerozměrné histogramy, shlukování a~Kohonenovy mapy. \emph{Modify} připraví objekty,
hodnoty a~proměnné, tedy vybírá, tvoří a~transformuje. \emph{Model} aplikuje vlastní
metody: stromy, regresi a~neuronové sítě. \emph{Assess} nakonec vyhodnotí spolehlivost
a~užitečnost výsledků, přičemž platí, že \emph{uživatel jim musí rozumět}.

\textbf{CRISP-DM} (\emph{CRoss-Industry Standard Process for Data Mining}) je otevřený
model procesu, neutrální oborově, nástrojově i~aplikačně. Vývoj financovala iniciativa
ESPRIT a~konsorcium tvořily ISL (později SPSS, dnes IBM), Teradata, NCR, Daimler AG
a~pojišťovna OHRA. Fází je šest a~jejich \textbf{pořadí není striktní}: výsledky jedné
fáze určují volbu dalších kroků a~některé kroky je nutné opakovat. Proces se proto
zobrazuje jako cyklus, kolem něhož vede ještě vnější cyklus nasazení a~monitoringu.

\begin{enumerate}
\item \textbf{Porozumění problému} (\emph{Business Understanding}) stanoví obchodní cíle
      a~kritéria úspěchu. Provede se inventura zdrojů, tedy lidí, informací, softwaru
      a~výpočetní kapacity, sepíšou se předpoklady a~omezení, rizika, náklady a~přínosy.
      Z~toho se odvodí cíle dobývání znalostí a~sestaví plán projektu.
\item \textbf{Porozumění datům} (\emph{Data Understanding}) začíná sběrem počátečních dat
      a~jejich popisem (množství, vlastnosti, formát). Následuje průzkum a~vizualizace,
      tedy rozdělení klíčových atributů, vztahy dvojic, subpopulace a~jednoduché
      statistiky. Fáze končí ověřením kvality dat.
\item \textbf{Příprava dat} (\emph{Data Preparation}) zahrnuje výběr, čištění, konstrukci
      odvozených atributů, integraci, agregaci a~formátování. Tyto kroky se provádějí
      iterativně a~v~různém pořadí.
\item \textbf{Modelování} (\emph{Modeling}) začíná volbou techniky, a~to s~ohledem na její
      předpoklady, což si může vynutit další úpravy dat. Pak se navrhne test, sestaví
      model s~různými parametry a~model se posoudí.
\item \textbf{Vyhodnocení} (\emph{Evaluation}) poměřuje výsledky s~\emph{obchodními}
      kritérii úspěchu, schvaluje modely, reviduje celý proces a~určuje další kroky.
\item \textbf{Nasazení} (\emph{Deployment}) prezentuje výsledky ve formě, kterou uživatel
      snadno interpretuje a~nasadí. Patří sem plán nasazení, monitoring a~údržba
      a~závěrečná zpráva s~revizí projektu, v~níž se dokumentují obtíže a~slepé uličky.
\end{enumerate}

\textbf{ASUM-DM} (IBM, 2015) reaguje na hlavní výtku vůči CRISP-DM, totiž na malou podporu
projektového řízení. Jeho fáze se jmenují \emph{Analyze}, \emph{Design}, \emph{Configure
\& Build}, \emph{Deploy} a~\emph{Operate \& Optimize}; průřezově přes ně vede
\emph{Project Management}.

Všechny tři metodologie popisují tentýž postup, liší se však původem, mírou neutrality
a~tím, kolik pozornosti věnují obchodní stránce věci a~řízení projektu:

\begin{center}
\small
\begin{tabular}{@{}llll@{}}
\toprule
& \textbf{SEMMA} & \textbf{CRISP-DM} & \textbf{ASUM-DM} \\
\midrule
Původ & SAS Institute & ESPRIT / konsorcium & IBM (2015) \\
Neutralita & vázáno na SAS EM & nástrojově neutrální & vázáno na IBM \\
Byznys fáze & chybí & ano (fáze 1 a~5) & ano \\
Projektové řízení & ne & slabé & explicitní, průřezové \\
Počet fází & 5 & 6 & 5 + PM \\
\bottomrule
\end{tabular}
\end{center}

\subsection{Shrnutí ke~zkoušce}

\begin{itemize}
\item \textbf{KDD} je netriviální proces získávání implicitních, dosud neznámých
      a~potenciálně užitečných informací z~dat; probíhá interaktivně a~iterativně.
\item \textbf{Pět fází KDD} jde za sebou takto: výběr $\to$ předzpracování $\to$
      transformace $\to$ dobývání $\to$ interpretace. Samotná příprava dat spotřebuje
      60--80~\% času.
\item \textbf{Kořeny oboru} leží v~umělé inteligenci (strojové učení), v~databázích
      a~ve statistice; čtvrtou složkou je potřeba podpořit rozhodování. \textbf{OLAP}
      odpovídá na otázky položené předem, kdežto KDD hledá vzory, o~nichž nevíme.
\item \textbf{Tři typy úloh} se liší tím, co se od výsledku žádá: klasifikace a~predikce
      staví přesnost nad jednoduchost, popis požaduje srozumitelnost při plném pokrytí
      a~hledání nugetů oceňuje překvapivost, i~když nález koncept nepokrývá.
\item Úlohy se dělí na \textbf{deskriptivní a~prediktivní}, a~nezávisle na tom podle
      dostupnosti cílové proměnné na \textbf{učení s~učitelem, bez učitele a~poloučené}.
\item \textbf{Škály atributů} --- nominální, ordinální, intervalová a~poměrová --- určují,
      jaké operace jsou přípustné, a~tím i~volbu metody. Nominální atribut se převádí na
      $v$ binárních atributů, poměrový s~exponenciálním chováním se logaritmuje.
\item Z~\textbf{metodologií} se čeká znalost SEMMA (Sample, Explore, Modify, Model,
      Assess), CRISP-DM (šest fází, jejichž pořadí není striktní, proces je cyklický)
      a~ASUM-DM, což je CRISP-DM doplněný o~projektové řízení.
\item Nejčastější doplňující otázka zní: \emph{proč pořadí fází CRISP-DM není striktní?}
      Protože výsledky fáze určují další kroky; zjištěná kvalita dat například vede zpět
      k~přeformulování cíle.
\end{itemize}

%=====================================================================
\newpage

\section{Příprava dat, výběr atributů a~analýza jejich relevance}
\label{sec:priprava}
%=====================================================================

Věta okruhu spojuje do jedné kapitoly dvě věci, které spolu souvisejí těsněji, než se na
první pohled zdá. Tou první je \emph{příprava dat}, tedy vše, co je nutné udělat, aby data
byla vůbec zpracovatelná. Tou druhou je \emph{výběr atributů s~analýzou jejich relevance},
tedy otázka, jak z~mnoha atributů vybrat právě ty, které nesou informaci. Druhá část je ve
skutečnosti speciálním případem té první, protože i~ona redukuje dimenzi dat. Okruh ji přesto
jmenuje zvlášť, a~to z~praktického důvodu: právě u~ní se očekává znalost konkrétních kritérií
a~schopnost je spočítat.

\subsection{Proč je příprava dat klíčová}

Předzpracování má \emph{klíčový význam pro úspěch aplikace}. Zároveň je obtížné a~časově
náročné, takže se vyplatí přizvat doménového experta, který o~datech ví to, co ze samotných
čísel vyčíst nelze. Práce má tři cíle:
\begin{itemize}
\item z~dostupných dat vybrat, případně teprve vytvořit, ta, která jsou pro řešenou úlohu
      relevantní;
\item reprezentovat je ve formě vhodné pro zvolený algoritmus;
\item dospět do stavu, kdy je výsledkem \emph{(jedna)} datová tabulka obsahující hodnoty
      atributů objektů.
\end{itemize}

\mimo{} Kvalitu dat posuzujeme podle několika rozměrů: \emph{přesnost} (\emph{accuracy}),
\emph{úplnost} (\emph{completeness}), \emph{konzistence}, \emph{aktuálnost}
(\emph{timeliness}), \emph{důvěryhodnost} a~\emph{interpretovatelnost}. Zásada \uv{garbage
in, garbage out} přitom platí bez výjimky. Sebelepší algoritmus na špatných datech vytvoří
jen sebejistý nesmysl.

\subsection{Čištění dat: chybějící hodnoty}

Pro chybějící hodnoty se nabízejí čtyři základní strategie, seřazené zde od nejjednodušší
k~nejpřesnější.

\begin{enumerate}
\item \textbf{Ignorovat objekty} s~libovolnou chybějící hodnotou. Postup je jednoduchý, ale
      při mnoha atributech tak ztratíme většinu dat. Navíc zavádí zkreslení všude tam, kde
      chybějící hodnoty nejsou náhodné.
\item \textbf{Nahradit novou hodnotou} \uv{nevím} (\emph{I don't know}). Chybějící hodnota se
      tím stane plnohodnotnou kategorií, což se hodí tehdy, když už sama absence nese
      informaci: typickým případem je nevyplněný příjem u~žádosti o~úvěr.
\item \textbf{Nahradit existující hodnotou atributu}. Nabízí se nejčastější hodnota (modus,
      u~numerických atributů průměr či medián), volba úměrná poměru všech hodnot, tedy
      náhodné losování podle empirického rozdělení, anebo prostě libovolná hodnota.
\item \textbf{Doplnit hodnotu pomocí modelu}. Naučíme prediktivní model (regrese, rozhodovací
      strom, $k$-NN), který chybějící atribut předpovídá z~ostatních. Je to volba nejpřesnější
      a~současně výpočetně nejnáročnější a~riskuje se u~ní, že mezi atributy zavlečeme umělou
      závislost.
\end{enumerate}

\mimo{} Terminologicky se rozlišují tři situace. MCAR (\emph{missing completely at random})
znamená zcela náhodné chybění, MAR (\emph{missing at random}) chybění závislé na pozorovaných
atributech a~MNAR (\emph{missing not at random}) chybění závislé na nepozorované hodnotě
samotné. Jen u~MCAR je prosté vypuštění nezkreslené.

Některé metody si s~chybějícími hodnotami poradí samy. C4.5 je při konstrukci stromu
ignoruje, protože kritérium dělení počítá jen ze známých hodnot, a~při klasifikaci je
odhaduje z~ostatních vzorů. CART chybějící hodnoty do rozhodovacího kritéria nezapočítává
a~naivní Bayes umí klasifikovat i~vzory popsané nekompletně (kap.~\ref{sec:metody}).

\subsection{Čištění dat: šum a~odlehlé hodnoty}

\emph{Šum} (\emph{noise}) je náhodná chyba nebo rozptyl v~měřené veličině. Vyhlazuje se čtyřmi
technikami. \textbf{Binning} rozdělí hodnoty do košů a~nahradí je průměrem, mediánem nebo
hranicí koše. \textbf{Regrese} nahradí hodnoty predikcí regresní funkce. \textbf{Shlukování}
označí za odlehlé ty hodnoty, které padnou mimo shluky. A~\textbf{kombinace stroje a~člověka}
předloží podezřelé hodnoty ke kontrole.

\begin{defbox}[Rozpětí, kvartily, mezikvartilové rozpětí, odlehlá hodnota]
\emph{Rozpětí} (\emph{range}) množiny dat je rozdíl nejvyšší a~nejnižší hodnoty.

\emph{Kvartily} $Q_1,Q_2,Q_3$ dělí data na čtyři části: 25~\% pozorování má hodnotu menší
než dolní kvartil $Q_1$, 50~\% pod mediánem $Q_2$ a~75~\% pod horním kvartilem $Q_3$.

\emph{Mezikvartilové rozpětí} (\emph{interquartile range}): $\mathrm{IQR}=Q_3-Q_1$.

\emph{Odlehlá hodnota} (\emph{outlier}) je extrémní hodnota ležící mimo celkový vzorec dat.
Podle běžného pravidla je odlehlá každá hodnota, pro kterou platí
\[ x > Q_3 + 1{,}5\cdot \mathrm{IQR} \qquad\text{nebo}\qquad x < Q_1 - 1{,}5\cdot \mathrm{IQR}. \]
\end{defbox}

\textbf{Kvartily} se pro $n$ vzestupně seřazených hodnot počítají u~diskrétních dat takto. Pro
$Q_1$ vydělíme $n$ čtyřmi: vyjde-li $n/4$ jako celé číslo, bereme střed mezi odpovídajícím
a~následujícím členem, a~nevyjde-li, zaokrouhlíme nahoru a~vezmeme odpovídající člen. Pro
$Q_3$ postupujeme analogicky s~podílem $3n/4$. U~spojitých dat se místo zaokrouhlení
interpoluje.

\begin{defbox}[Příklad --- krabicový graf a~odlehlé hodnoty]
Hladina glukózy v~krvi u~30 žen (mmol/l), stonkový diagram (klíč: $2\mid 1$ znamená 2,1):
\begin{center}\ttfamily
2 | 2 2 3 3 5 7 \hfill (6)\\
3 | 1 2 6 7 7 7 8 8 8 8 9 9 9 \hfill (13)\\
4 | 0 0 0 0 4 5 6 7 8 \hfill (9)\\
5 | 1 5 \hfill (2)
\end{center}
Zde je $n=30$, takže $30/4=7{,}5$ zaokrouhlíme na 8; osmé číslo je 3,2, a~tedy $Q_1=3{,}2$.
Medián vychází $Q_2=3{,}8$. Pro $Q_3$ počítáme $3\cdot 30/4=22{,}5$, zaokrouhlíme na 23,
a~protože 23.\ číslo je 4,0, dostáváme $Q_3=4{,}0$.

Odtud $\mathrm{IQR}=4{,}0-3{,}2=0{,}8$, dolní mez $=3{,}2-1{,}5\cdot 0{,}8=2{,}0$ a~horní mez
$=4{,}0+1{,}5\cdot 0{,}8=5{,}2$.

\textbf{Odlehlá hodnota je 5,5.} Nejvyšší hodnota, která odlehlá není, je 5,1: v~krabicovém
grafu (\emph{box plot}) k~ní vede vous, zatímco odlehlá hodnota se vyznačí křížkem.
\end{defbox}

Násobek 1,5 je pouhá konvence; s~násobkem 1 dostaneme kritérium přísnější. Jak se meze počítají
v~praxi, ukazují dva soubory měření. Pro želví samce s~$Q_1=400$~g a~$Q_3=580$~g je
$\mathrm{IQR}=180$ a~meze vycházejí $400-180=220$ a~$580+180=760$, takže žádná z~hodnot 400,
260, 550 a~640~g není odlehlá a~největší hmotnost, kterou ještě za outlier neoznačíme, je
760~g. U~samic je $Q_1=260$, $Q_3=340$ a~$\mathrm{IQR}=80$, meze tedy vycházejí 180 a~420
a~odlehlé jsou hodnoty 170~g a~440~g.

\subsection{Integrace dat}

\mimo{} Při \emph{integraci} (\emph{data integration}) spojujeme více zdrojů do jedné tabulky
a~narazíme přitom na čtyři typické problémy. \textbf{Problém identity entity} nastává, když
atribut \texttt{cust\_id} v~jedné databázi odpovídá atributu \texttt{customer\_number}
v~druhé; k~jeho řešení jsou potřeba metadata a~párování záznamů. \textbf{Redundance} znamená,
že jeden atribut je odvoditelný z~jiných, například roční příjem $=$ 12$\times$měsíční;
u~numerických atributů ji odhalí korelační analýza, u~kategoriálních $\chi^2$ test (viz
odd.~\ref{ssec:filtr}). \textbf{Konflikty hodnot} vznikají z~různých jednotek, škál
a~kódování. A~\textbf{duplicity} znamenají, že tentýž objekt je v~datech vícekrát, pokaždé
s~drobně odlišnými hodnotami.

\subsection{Transformace a~standardizace atributů}

V~euklidovském prostoru se standardizace atributů doporučuje proto, aby všechny atributy mohly
mít na výpočet vzdáleností \emph{stejný vliv}. Bez ní rozhoduje ten atribut, který má náhodou
největší rozsah. Pro body $x_i=(0{,}1;\,20)$ a~$x_j=(0{,}9;\,720)$ vyjde
\[ \mathrm{dist}(x_i,x_j)=\sqrt{(0{,}9-0{,}1)^2+(720-20)^2}=700{,}000457, \]
takže vzdálenost je zcela ovládnuta rozdílem $720-20=700$ ve druhém atributu a~první atribut
na ni nemá prakticky žádný vliv. Standardizace nutí atributy do společného rozsahu hodnot.

\begin{defbox}[Základní standardizační transformace]
Nechť $f$ je atribut a~$x_{if}$ jeho hodnota u~$i$-tého objektu.

\textbf{Desítkové škálování} (\emph{decimal scaling}) do $[-1,1]$: dělíme hodnoty nejmenší
mocninou 10, která udrží všechny transformované hodnoty v~$[-1,1]$.

\textbf{Min-max normalizace} (\emph{range standardization}) do $[0,1]$:
\[ \mathrm{range}(x_{if})=\frac{x_{if}-\min(f)}{\max(f)-\min(f)}, \]
do $[-1,1]$ pak $\mathrm{range}_{[-1,1]}(x_{if})=2\cdot \mathrm{range}(x_{if})-1$.

\textbf{Standardizace podle průměrné absolutní odchylky} dá nulový průměr a~jednotkovou
střední absolutní odchylku:
\[ m_f=\tfrac1n\big(x_{1f}+\cdots+x_{nf}\big),\qquad
s_f=\tfrac1n\big(|x_{1f}-m_f|+\cdots+|x_{nf}-m_f|\big),\qquad
z(x_{if})=\frac{x_{if}-m_f}{s_f}. \]

\textbf{Standardizace podle směrodatné odchylky} (\emph{std-score}) dá nulový průměr
a~jednotkovou korigovanou výběrovou směrodatnou odchylku:
\[ \mathrm{cor}s_f=\sqrt{\frac{1}{n-1}\sum_{i=1}^n\big(x_{if}-m_f\big)^2},\qquad
\mathrm{std}(x_{if})=\frac{x_{if}-m_f}{\mathrm{cor}s_f}. \]
\end{defbox}

Volba mezi oběma z-skóry se řídí jednoduchou úvahou. Střední absolutní odchylka je k~odlehlým
hodnotám \emph{robustnější} než směrodatná, protože se odchylky neumocňují na druhou, takže
odlehlé hodnoty stlačí ostatní data méně.

\subsection{Diskretizace numerických atributů}

Řada metod pracuje výhradně s~kategoriálním vstupem: patří sem asociační pravidla, ID3, CHAID
i~bayesovská klasifikace s~diskrétními atributy. \emph{Diskretizace} proto rozdělí rozsah
numerického atributu na intervaly. Jednotlivé metody se od sebe liší ve čtyřech ohledech:
\emph{strategií} (dělí intervaly shora, nebo je slučují zdola), \emph{počtem} intervalů,
\emph{typem} intervalů a~\emph{kritériem} kvality, kterým bývá minimální klasifikační chyba,
entropie nebo $\chi^2$ test.

\begin{itemize}
\item \textbf{Ekvidistantní} diskretizace (\emph{equi-width}) rozdělí rozsah na intervaly
      stejné délky. Je nejjednodušší, zato citlivá na odlehlé hodnoty a~na nerovnoměrné
      rozdělení.
\item \textbf{Ekvifrekvenční} diskretizace (\emph{equi-depth}) \mimo{} vytváří intervaly se
      stejným počtem objektů.
\item \textbf{S~využitím informace o~třídě} se hranice volí tak, aby intervaly vyšly co
      nejčistší. Kritériem je buď entropie, a~pak se dělí rekurzivně jako v~C4.5, nebo
      $\chi^2$, a~pak se slučuje zdola jako v~CHAID.
\end{itemize}

\subsubsection{Fuzzy diskretizace}

U~běžné diskretizace jsou hranice intervalů ostré, takže dva objekty s~téměř stejnou hodnotou
mohou skončit v~různých kategoriích. \emph{Fuzzy diskretizace} hranice \uv{rozostří}: jedné
numerické hodnotě pak mohou odpovídat \emph{dvě} diskretizované hodnoty, přičemž součet
příslušných charakteristických funkcí je 1. Vznikají tak \uv{fuzzy}-objekty s~vahou $<1$; váha
fuzzy-objektu je dána \emph{součinem} hodnot charakteristických funkcí všech atributů:
\[ w = \prod_j \mu_j. \]
Za tuto jemnost se ale platí. Hrozí ztráta informace a~\emph{kontradikce}, tedy situace, kdy
dva objekty mají stejné hodnoty diskretizovaných atributů, ale patří do různých tříd.

\begin{defbox}[Procedura \emph{Interval} (fuzzy diskretizace)]
\begin{enumerate}
\item Má-li posloupnost hodnot stejnou třídu, vytvoř z~ní interval s~charakteristickou
      funkcí rovnou 1 na celém intervalu.
\item Pro každý interval $\mathit{Int}_i$ patřící do třídy \uv{?}:
      \begin{itemize}
      \item patří-li sousední intervaly $\mathit{Int}_{i-1}$ a~$\mathit{Int}_{i+1}$ do
            \emph{téže} třídy, spoj je do jednoho intervalu
            $\mathit{Int}_{i-1}\cup \mathit{Int}_i\cup \mathit{Int}_{i+1}$
            (charakteristická funkce $=1$ na celém intervalu);
      \item jinak vytvoř \emph{dva fuzzy-intervaly}: spojením $\mathit{Int}_{i-1}\cup
            \mathit{Int}_i$ s~charakteristickou funkcí rovnou 1 na
            $[L_{i-1},U_{i-1}]$ a~lineárně klesající
            $\frac{U_{i+1}-x}{U_{i+1}-L_{i-1}}$ na přechodové oblasti, a~analogicky
            spojením $\mathit{Int}_i\cup \mathit{Int}_{i+1}$ s~funkcí lineárně rostoucí
            $\frac{x-L_{i-1}}{U_{i+1}-L_{i-1}}$.
      \end{itemize}
\item Zajisti spojité pokrytí celého rozsahu (mezery mezi intervaly se chápou jako
      intervaly třídy \uv{?}).
\end{enumerate}
\end{defbox}

\subsubsection{Seskupování hodnot kategoriálních atributů (KEX)}

Má-li kategoriální atribut příliš mnoho hodnot, vyplatí se je seskupit. \textbf{Cílem} je
vytvořit takové grupy $A(\mathit{Grp})$, aby se podmíněná pravděpodobnost
$P(\mathrm{Class}\mid A(\mathit{Grp}))$ \emph{významně} lišila od apriorní pravděpodobnosti
$P(\mathrm{Class})$, a~vytvořit jich tolik, kolik je tříd, plus jednu navíc pro kód \uv{?}.

\begin{defbox}[Algoritmus seskupování hodnot (KEX)]
\textbf{Hlavní cyklus:}
\begin{enumerate}
\item vytvoř uspořádaný seznam hodnot uvažovaného atributu;
\item pro každou hodnotu: \emph{(a)}~spočítej četnosti výskytu objektů s~touto hodnotou
      v~jednotlivých třídách; \emph{(b)}~přiřaď hodnotě kód třídy procedurou
      \emph{Assign};
\item vytvoř grupy hodnot procedurou \emph{Group}.
\end{enumerate}

\textbf{Procedura \emph{Assign}:} patří-li pro danou hodnotu všechny objekty do téže třídy,
přiřaď hodnotě kód této třídy; jinak --- liší-li se rozdělení tříd pro danou hodnotu
\emph{významně} (podle $\chi^2$ testu) od apriorního rozdělení tříd --- přiřaď kód
nejčastější třídy, a~pokud ne, přiřaď kód \uv{?}.

\textbf{Procedura \emph{Group}:} vytvoř grupu z~hodnot, kterým byl přiřazen stejný kód
třídy.
\end{defbox}

\subsection{Redukce dat}

\subsubsection{Příliš mnoho objektů}

Roste-li počet objektů, dávkové zpracování se stává problematickým. Řešení jsou tři. Buď se
pracuje jen se \textbf{vzorkem} objektů, nebo se data uloží tak, aby k~nim šlo přistupovat bez
nutnosti mít je celá v~operační paměti, anebo se postaví \textbf{více modelů nad podmnožinami}
a~ty se pak zkombinují do ensemblu (kap.~\ref{sec:metody}).

Klíčovou vlastností vybraných dat je jejich \textbf{reprezentativnost}. Vzorek má zachytit
povahu dat co nejlépe, a~to zejména co se týče souhlasu rozdělení hodnot atributů na vzorku
a~na celých datech.

Vlastní kapitolou jsou \textbf{nevyvážené třídy} v~původních datech, které vedou k~tendenci
preferovat majoritní třídu. Řešení jsou dvě: přiřadit různým typům chyb různé
\emph{váhy}, tedy použít cenově citlivé učení, nebo vybírat objekty z~různých tříd
s~\emph{různou pravděpodobností}, což znamená převzorkovat minoritní třídu nebo podvzorkovat
majoritní.

\subsubsection{Příliš mnoho atributů}

Trápí-li nás naopak počet atributů, vedou z~toho dvě principiálně odlišné cesty.

\begin{itemize}
\item \textbf{Redukce transformací} vytvoří z~existujících atributů \emph{méně nových},
      typicky jako jejich lineární kombinace; hlavním zástupcem je analýza hlavních komponent
      (PCA). Nevýhody jsou dvě: je nutné znát hodnoty \emph{všech} původních atributů a~nové
      atributy nemusí mít jasnou interpretaci.
\item \textbf{Redukce selekcí} (\emph{feature selection}) z~existujících atributů vybere jen
      ty nejdůležitější. Interpretace se tím zachovává a~nepotřebné atributy se pak vůbec
      nemusí měřit.
\end{itemize}

\mimo{} PCA hledá ortogonální směry maximálního rozptylu, tedy vlastní vektory kovarianční
matice s~největšími vlastními čísly. Je to transformace \emph{nesupervizovaná}, a~proto může
zahodit směr, který je pro predikci klíčový, ale má malý rozptyl.

\subsection{Analýza relevance atributů: filtry a~wrappery}
\label{ssec:filtr}

Tím se dostáváme k~druhé polovině okruhu. Hledáme atributy, které nejlépe přispívají ke
správné klasifikaci objektů, a~existují na to dva principiálně různé přístupy.

\begin{defbox}[Filtry vs.\ wrappery]
\textbf{Filtrační metody} (\emph{filter}) spočítají pro každý atribut charakteristiku
odrážející jeho přínos ke klasifikaci, atributy podle ní seřadí a~vyberou stanovený počet
nejlepších. Charakteristika je \emph{nezávislá} na následně použitém modelu, a~metoda je proto
rychlá.

\textbf{Obalové metody} (\emph{wrapper}) postaví pomocí algoritmu strojového učení model nad
\emph{podmnožinou} atributů a~ten model vyhodnotí; nakonec se použije nejlepší z~postavených
modelů i~s~odpovídající podmnožinou atributů. Metoda je výpočetně drahá, zato zohledňuje
interakci atributů \emph{s~konkrétním} modelem.

Prostor podmnožin se prohledává dvěma směry:
\begin{itemize}
\item \uv{zdola nahoru} (\emph{bottom-up}, \emph{dopředná selekce}) začíná od modelů nad
      jednotlivými atributy a~postupně přidává další;
\item \uv{shora dolů} (\emph{top-down}, \emph{zpětná eliminace}) začíná od modelu nad
      úplnou množinou atributů a~postupně odstraňuje nepotřebné.
\end{itemize}
\end{defbox}

Filtry mají \textbf{zásadní omezení}: hodnotí každý atribut \emph{samostatně}, takže vzájemný
vliv několika atributů na přesnost klasifikace zachycují jen obtížně. Kritérium sice lze
počítat i~pro celou \emph{množinu} atributů, tedy pro všechny jejich kombinace, ale
kombinatorická explozivnost to prakticky omezuje na hladové dopředné či zpětné prohledávání.
\mimo{} Praktickým důsledkem je, že filtr propustí dvojici silně korelovaných atributů,
protože oba dostanou vysoké skóre, přestože druhý z~nich nepřináší nic nového. A~naopak
zahodí atribut, který je užitečný jen \emph{v~kombinaci} s~jiným, což je typicky případ
XOR-vztahu.

\subsubsection{Kritéria na kontingenčních tabulkách}

Vztah vstupního atributu $A$ s~hodnotami $v_1,\dots,v_R$ a~cílového atributu $C$ se třídami
$v_1,\dots,v_S$ se popíše kontingenční tabulkou o~četnostech $a_{kl}$ kombinace
$A(v_k)\wedge C(v_l)$. Její marginály jsou
\[ r_k=\sum_{l=1}^{S}a_{kl},\qquad s_l=\sum_{k=1}^{R}a_{kl},\qquad
n=\sum_{k=1}^{R}\sum_{l=1}^{S}a_{kl}. \]
Z~tabulky pak odhadneme pravděpodobnosti jako $P(A(v_k)\wedge C(v_l))=a_{kl}/n$,
$P(A(v_k))=r_k/n$ a~$P(C(v_l))=s_l/n$. Očekávaná četnost při nezávislosti je
$e_{kl}=r_k s_l/n$.

\begin{defbox}[$\chi^2$ jako kritérium relevance]
\[ \chi^2=\sum_{k=1}^{R}\sum_{l=1}^{S}\frac{(a_{kl}-e_{kl})^2}{e_{kl}}
   = n\sum_{k=1}^{R}\sum_{l=1}^{S}\frac{\big(a_{kl}-\tfrac{r_ks_l}{n}\big)^2}{r_ks_l}
   \qquad\text{(čím větší, tím lepší).}\]
Pro ruční výpočet se hodí alternativní tvar $\chi^2=\sum_k\sum_l \frac{a_{kl}^2}{e_{kl}}-n$.

\textbf{Jako statistický test:} při platnosti nulové hypotézy o~nezávislosti
\[ H_0:\ P(A=v_k\wedge C=v_l)=P(A=v_k)\,P(C=v_l)\quad \forall k,l \]
má statistika $\chi^2$ rozdělení s~$(R-1)(S-1)$ stupni volnosti. Platí-li
$\chi^2\ge\chi^2_{(R-1)(S-1)}(\alpha)$, nulovou hypotézu na hladině významnosti $\alpha$
zamítáme a~přijímáme alternativní hypotézu o~\emph{závislosti}.
\end{defbox}

\begin{defbox}[Příklad --- $\chi^2$ test na čtyřpolní tabulce]
\begin{center}
\begin{tabular}{@{}llrrr@{}}
\toprule
& & \multicolumn{2}{c}{Poskytnutí úvěru} & \\
& & ANO & NE & $r_k$ \\
\midrule
Výše příjmu & vysoký & 50 & 0 & 50 \\
            & nízký & 30 & 40 & 70 \\
\midrule
$s_l$ & & 80 & 40 & 120 \\
\bottomrule
\end{tabular}
\end{center}

Očekávané četnosti vycházejí $e_{11}=50\cdot 80/120=33{,}3$; $e_{12}=50\cdot 40/120=16{,}7$;
$e_{21}=70\cdot 80/120=46{,}7$ a~$e_{22}=70\cdot 40/120=23{,}3$. Rozdíly proti skutečnosti
jsou tedy $+16{,}7$, $-16{,}7$, $-16{,}7$ a~$+16{,}7$, což dává hodnotu statistiky
$\chi^2=42{,}857$.

Na hladině významnosti $\alpha=0{,}05$ je kritická hodnota $\chi^2_{(1)}(0{,}05)=3{,}84$.
Protože $42{,}857>3{,}84$, platí, že \textbf{mezi výší příjmu a~poskytnutím úvěru je
závislost}.
\end{defbox}

\begin{defbox}[Fisherův exaktní test]
$\chi^2$ test lze použít jen pro dostatečně velké četnosti: musí platit $(r_k s_l)/n\ge 5$
pro všechna $k,l$. Pro čtyřpolní tabulky s~malými četnostmi se proto používá \emph{Fisherův
test}. Pravděpodobnost, že čtyřpolní tabulka má při daných marginálech $r_1,r_2,s_1,s_2$
právě četnosti $a_{kl}$, je
\[ p=\frac{r_1!\,r_2!\,s_1!\,s_2!}{n!\,a_{11}!\,a_{12}!\,a_{21}!\,a_{22}!}
     =\frac{\binom{s_1}{a_{11}}\binom{s_2}{a_{12}}}{\binom{n}{r_1}}. \]
Pravděpodobnosti se pak sečtou přes všechny uvažované hodnoty skutečných četností při daných
marginálech (nechť $m=\min(r_1,s_1)$):
\[ P=\sum_{i=0}^{m}\frac{r_1!\,r_2!\,s_1!\,s_2!}{n!\,(a_{11}-i)!\,(a_{12}+i)!\,(a_{21}+i)!\,(a_{22}-i)!}. \]
Platí-li $P\le\alpha$, nulová hypotéza se na hladině $\alpha$ zamítá.

\textbf{Příklad.} Pro tabulku s~marginály $r=(12,18)$, $s=(6,24)$ a~$n=30$ vychází
$P_0=0{,}031$ a~$P_1=0{,}173$, takže $P=P_0+P_1=0{,}204>0{,}05$ a~\textbf{nulovou hypotézu
o~nezávislosti přijímáme}. (Kontrola: součet $P_0,\dots,P_6$
$=0{,}031+0{,}173+0{,}340+0{,}302+0{,}128+0{,}024+0{,}002=1{,}000$.)
\end{defbox}

\subsubsection{Kritéria založená na entropii}

Druhou rodinu filtračních kritérií dodává teorie informace.

\begin{defbox}[Entropie, vzájemná informace, informační míra závislosti]
\textbf{Entropie} atributu $A$ jako vážený průměr entropií jeho hodnot
(čím \emph{menší}, tím lepší):
\[ H(A)=\sum_{k=1}^{R}\frac{r_k}{n}\,H\big(A(v_k)\big),\qquad
   H\big(A(v_k)\big)=-\sum_{l=1}^{S}\frac{a_{kl}}{r_k}\log_2\frac{a_{kl}}{r_k}. \]

\textbf{Vzájemná informace} (\emph{mutual information}):
\[ \mathrm{MI}(A,C)=\sum_{k=1}^{R}\sum_{l=1}^{S}\frac{a_{kl}}{n}
   \log_2\frac{n\,a_{kl}}{r_k s_l}. \]

\textbf{Informační míra závislosti} (\emph{information measure of dependence}) je
normalizovaná vzájemná informace (čím \emph{větší}, tím lepší):
\[ \mathrm{ID}(A,C)=\frac{\mathrm{MI}(A,C)}{H(C)}
   =\frac{\mathrm{MI}(A,C)}{-\sum_{l=1}^{S}\frac{s_l}{n}\log_2\frac{s_l}{n}}. \]
\end{defbox}

Normalizace v~$\mathrm{ID}$ není kosmetická. Samotná $\mathrm{MI}$ systematicky zvýhodňuje
atributy s~mnoha hodnotami, takže v~limitě vychází jako \uv{nejinformativnější} atribut prostý
identifikátor objektu. Totéž zvýhodnění řeší poměrný informační zisk v~C4.5
(kap.~\ref{sec:metody}).

\subsection{Vztahy mezi atributy: regrese a~korelace}

Kontingenční tabulky popisují vztah \emph{kategoriálních} atributů. U~numerických atributů se
místo nich používá regresní a~korelační analýza, přičemž každá z~nich odpovídá na jinou
otázku. \emph{Korelační analýza} říká, zda mezi dvěma numerickými veličinami lineární vztah
vůbec \emph{existuje}; \emph{lineární regrese} určuje jeho \emph{parametry}.

K~\textbf{regresní analýze} vedou čtyři důvody. \emph{(1)}~Měření výstupu je nákladné, takže
jej chceme predikovat ze snadno dostupných vstupů. \emph{(2)}~Vstupy jsou známé dříve než
výstup, takže z~nich lze výstup odhadnout dopředu. \emph{(3)}~Řízené vstupní hodnoty pomohou
najít správné chování výstupu. \emph{(4)}~Mezi vstupem a~výstupem může být kauzální vztah,
který chceme najít.

\begin{defbox}[Lineární regrese jedné proměnné metodou nejmenších kvadrátů]
Vzory $[x_1,y_1],\dots,[x_p,y_p]$, regresní rovnice $y=\beta x+\alpha$, kvadratická chyba
\[ E=\sum_{i=1}^{p}\varepsilon_i^2=\sum_{i=1}^{p}\big(y_i-\hat y_i\big)^2
   =\sum_{i=1}^{p}\big(y_i-\beta x_i-\alpha\big)^2. \]
Položíme parciální derivace nule:
\[ \frac{\partial E}{\partial\alpha}=-2\sum_i\big(y_i-\beta x_i-\alpha\big)=0,\qquad
   \frac{\partial E}{\partial\beta}=-2\sum_i\big(y_i-\beta x_i-\alpha\big)x_i=0, \]
odkud po úpravě $\alpha=\bar y-\beta\bar x$ a
\[ \beta=\frac{\sum_{i=1}^{p}(x_i-\bar x)(y_i-\bar y)}{\sum_{i=1}^{p}(x_i-\bar x)^2}
   \qquad(\text{tj.\ výběrová kovariance dělená výběrovým rozptylem }x). \]
Predikce: $\hat y=\beta x+\alpha$.
\end{defbox}

\begin{defbox}[Korelační koeficient]
Posouzení \uv{množství} lineární závislosti:
\[ \rho(x,y)=\frac{\sigma_{xy}}{\sqrt{\sigma_x^2\sigma_y^2}}=\frac{\sigma_{xy}}{\sigma_x\sigma_y},
\qquad
\sigma_{xy}=\frac{1}{n-1}\sum_i (x_i-\bar x)(y_i-\bar y), \]
\[ \sigma_x^2=\frac{1}{n-1}\sum_i (x_i-\bar x)^2,\qquad
   \sigma_y^2=\frac{1}{n-1}\sum_i (y_i-\bar y)^2. \]
Platí $\rho\in[-1,1]$, přičemž $|\rho|=1$ znamená přesný lineární vztah a~$\rho=0$ žádný
\emph{lineární} vztah. Nulová korelace ovšem neznamená žádnou závislost: $\rho$ nezachytí
vztahy nelineární.
\end{defbox}

\textbf{Vícerozměrná lineární regrese} předpokládá lineární závislost vysvětlované veličiny
$Y$ na několika vysvětlujících veličinách $X_1,\dots,X_m$, tedy
$\hat y=\beta_0+\beta_1x_1+\cdots+\beta_mx_m$, maticově $\hat Y=X\beta$. Zde je $X$ rozšířená
matice vstupních vzorů, jejíž první sloupec tvoří jedničky, a~$\beta=(\beta_0,\dots,\beta_m)$.
Kvadratická odchylka je $E=(Y-X\beta)^{\!\top}(Y-X\beta)$ a~z~podmínky
$\partial E/\partial\beta=0$ dostáváme normální rovnice $X^{\!\top}\!X\beta=X^{\!\top}Y$,
tedy
\[ \beta=\big(X^{\!\top}\!X\big)^{-1}X^{\!\top}Y. \]
U~složitých praktických úloh je tento výpočet nákladný, a~proto se sahá po aproximativních
(iterativních) řešeních.

\textbf{Nelineární regrese} předpokládá složitější funkční závislost, například kvadratickou
nebo exponenciální. Jejím nejdůležitějším případem je \emph{logistická regrese}, která
modeluje kategoriální (dvouhodnotový) výstup. Protože jde současně o~klasifikátor, je
zpracována v~kap.~\ref{sec:metody}, odd.~\ref{ssec:logreg}.

\subsection{Shrnutí ke~zkoušce}

\begin{itemize}
\item \textbf{Cílem přípravy dat} je vybrat relevantní data a~reprezentovat je ve formě
      vhodné pro zvolený algoritmus; výsledkem je \emph{jedna} tabulka objekty $\times$
      atributy.
\item \textbf{Chybějící hodnoty} se řeší čtyřmi způsoby: ignorovat celý objekt, zavést novou
      hodnotu \uv{nevím}, dosadit existující hodnotu (modus, poměrně podle rozdělení, nebo
      libovolně), anebo hodnotu doplnit modelem. U~zkoušky se chce zdůvodnit, kdy je která
      volba vhodná: nese-li informaci už sama absence, volí se \uv{nevím}.
\item \textbf{Odlehlé hodnoty} se určují z~mezikvartilového rozpětí $\mathrm{IQR}=Q_3-Q_1$;
      odlehlá je každá hodnota mimo interval
      $[Q_1-1{,}5\,\mathrm{IQR},\,Q_3+1{,}5\,\mathrm{IQR}]$. Ke zkoušce patří výpočet kvartilů
      ze vzestupně seřazených dat podle pravidla $n/4$ a~$3n/4$ se zaokrouhlením nahoru
      a~nakreslení box plotu.
\item \textbf{Standardizaci} lze provést desítkovým škálováním, min-max normalizací do
      $[0,1]$ i~do $[-1,1]$, z-skórem podle střední absolutní odchylky (ten je robustnější)
      nebo std-skórem podle směrodatné odchylky. Motivací je, že bez standardizace ovládne
      vzdálenost atribut s~největším rozsahem.
\item \textbf{Diskretizace} může být ekvidistantní, ekvifrekvenční nebo využívající informaci
      o~třídě (entropie, $\chi^2$). \textbf{Fuzzy diskretizace} hranice rozostří: jedné
      hodnotě odpovídají dvě kategorie se součtem charakteristických funkcí 1 a~váha
      fuzzy-objektu je součin $\prod_j\mu_j$; rizikem je ztráta informace a~kontradikce.
      \textbf{KEX} seskupuje hodnoty kategoriálního atributu podle $\chi^2$ testu proti
      apriornímu rozdělení tříd, a~to procedurami \emph{Assign} a~\emph{Group}.
\item \textbf{Redukce dat} má dvě větve. Je-li objektů příliš mnoho, sáhne se po vzorkování
      (a~hlídá se reprezentativnost!), po efektivním uložení dat nebo po více modelech nad
      podmnožinami; nevyvážené třídy se řeší váhami chyb nebo různými pravděpodobnostmi
      výběru. Je-li příliš mnoho atributů, volí se buď \emph{transformace} (PCA dá méně nových
      atributů jako lineární kombinace, ale je nutné znát všechny původní a~chybí
      interpretace), nebo \emph{selekce} (podmnožina původních atributů, u~níž interpretace
      zůstává).
\item \textbf{Filtry a~wrappery} se liší v~tom, co hodnotí. Filtr počítá charakteristiku pro
      každý atribut samostatně a~je nezávislý na modelu, takže je rychlý, ale nezachytí
      interakce; wrapper staví model nad podmnožinou atributů a~vyhodnocuje jej, takže je
      drahý, zato interakce zohledňuje. Prohledávat lze \emph{bottom-up} (přidáváním) nebo
      \emph{top-down} (odebíráním).
\item \textbf{Filtračními kritérii} jsou $\chi^2$ (větší $=$ lepší), entropie $H(A)$ (menší
      $=$ lepší), $\mathrm{MI}(A,C)$ a~$\mathrm{ID}(A,C)=\mathrm{MI}/H(C)$ (větší $=$ lepší).
      Normalizaci v~$\mathrm{ID}$ je nutné umět zdůvodnit: bez ní jsou zvýhodněny atributy
      s~mnoha hodnotami.
\item \textbf{$\chi^2$ test} má $(R-1)(S-1)$ stupňů volnosti a~podmínku $r_ks_l/n\ge 5$; pro
      malé četnosti se místo něj použije \textbf{Fisherův exaktní test}. Výpočet na čtyřpolní
      tabulce se u~zkoušky prochází celý (příklad s~úvěrem dá $\chi^2=42{,}857>3{,}84$, a~tedy
      závislost).
\item \textbf{Regrese a~korelace}: koeficienty se hledají metodou nejmenších kvadrátů, kde
      $\beta$ vyjde jako kovariance dělená rozptylem a~$\alpha=\bar y-\beta\bar x$;
      vícerozměrně je $\beta=(X^{\!\top}\!X)^{-1}X^{\!\top}Y$. Korelační koeficient
      $\rho=\sigma_{xy}/(\sigma_x\sigma_y)$ měří jen závislost \emph{lineární}.
\end{itemize}

%=====================================================================
\newpage

\section{Metody pro dobývání znalostí}
\label{sec:metody}
%=====================================================================

Tohle je nejobsáhlejší věta celého okruhu a~jako jediná jmenuje rovnou tři skupiny metod.
Nejde o~konkurenční přístupy k~téže úloze, nýbrž o~odpovědi na tři různé situace, do kterých
se analytik dostává podle toho, jaká data má před sebou. \emph{Asociační pravidla} se hodí
na transakční data, v~nichž žádný atribut nehraje roli cíle. \emph{Učení s~učitelem}
naopak předpokládá, že cílový atribut existuje a~má se předpovídat. \emph{Klastrová analýza}
pracuje s~metrickými daty, kde cíl opět chybí a~struktura se hledá v~podobnosti samotných
objektů. Podle těchto tří situací je kapitola rozdělena na tři oddíly.

\subsection{Asociační pravidla a~frekventované vzory}
\label{ssec:ar}

\subsubsection{Analýza nákupního košíku}

\emph{Analýza nákupního košíku} (\emph{market basket analysis}, MBA) je archetypální úloha
celé oblasti a~shrne ji jediná otázka: \uv{které položky se v~košíku vyskytují společně?}
Výsledky mají podobu pravidel, takže se dají použít okamžitě a~bez dalšího překladu. Obchodník
podle nich plánuje prodejnu a~její rozvržení, nabízí kupony a~časově omezené slevy nebo
produkty \uv{balíčkuje}.

Základní pojmy jsou dva a~jsou jednoduché. \textbf{Položka} (\emph{item}) je nabízený produkt
nebo služba, \textbf{transakce} obsahuje jednu nebo více položek. Nad nimi se staví
\textbf{frekvenční tabulka}, která udává počet společných výskytů každých dvou položek
v~transakcích; na diagonále pak nese počet transakcí obsahujících danou položku.

\begin{defbox}[Příklad --- frekvenční tabulka]
Uvažujme následující transakce v~obchodě s~potravinami:
\begin{center}
\small
\begin{tabular}{@{}rl@{\hspace{1.2cm}}lrrrr@{}}
\toprule
Zák. & Položky & & chléb & máslo & mléko & káva \\
\midrule
1 & chléb, máslo          & chléb & 4 & 3 & 1 & 2 \\
2 & mléko, chléb, máslo   & máslo & 3 & 4 & 1 & 2 \\
3 & chléb, káva           & mléko & 1 & 1 & 1 & 0 \\
4 & chléb, máslo, káva    & káva  & 2 & 2 & 0 & 3 \\
5 & káva, máslo           &       &   &   &   &   \\
\bottomrule
\end{tabular}
\end{center}
I~z~takto malé tabulky je povaha prodeje vidět okamžitě: \emph{chléb a~máslo se kupují spolu},
kdežto \emph{mléko se nikdy nekupuje s~kávou}.
\end{defbox}

Nalezené vztahy se zapisují jako pravidla ve tvaru \texttt{IF} \emph{podmínka} \texttt{THEN}
\emph{závěr}. Od dobrého asociačního pravidla se přitom žádají dvě vlastnosti:
\begin{itemize}
\item \textbf{snadno pochopitelné} pravidlo je takové, u~kterého lze nalezený vztah snadno
      ověřit;
\item \textbf{použitelné} pravidlo obsahuje užitečnou informaci, která vede k~dalším zásahům;
\end{itemize}
a~naopak se dvěma vlastnostem chceme vyhnout. Pravidlo nemá být \textbf{triviální}, tedy
takové, jehož výsledek už každý zná, ani \textbf{nevysvětlitelné}: pro takové pravidlo
neexistuje vysvětlení a~nevede k~žádné akci.

\subsubsection{Support, confidence a~lift}

Pravidel se v~datech najdou tisíce, a~tak je potřeba měřit, která z~nich stojí za pozornost.
Slouží k~tomu tři míry, a~každá se ptá na něco jiného.

\begin{defbox}[Support, confidence, lift]
\textbf{Support} (podpora) říká, jak často lze pravidlo použít:
\[ \mathrm{sup}(R)=\frac{\text{počet transakcí obsahujících } i \text{ a } j}
   {\text{počet všech transakcí}}\cdot 100\,\%. \]

\textbf{Confidence} (spolehlivost) říká, jak spolehlivé jsou výsledky pravidla:
\[ \mathrm{conf}(R)=\frac{\text{počet transakcí obsahujících } i \text{ a } j}
   {\text{počet transakcí obsahujících } i}\cdot 100\,\%. \]

\textbf{Lift} (zdvih) měří, kolikrát je lepší použít pravidlo pro predikci, než jen
předpokládat jeho výsledek:
\[ \mathrm{lift}(R)=\frac{p(i\wedge j)}{p(i)\,p(j)}. \]
Je-li $\mathrm{lift}<1$, je pravidlo pro predikci \emph{horší než náhodná volba}. V~takovém
případě může dát lepší pravidlo \textbf{negace závěru}, tedy tvar \texttt{IF} \emph{podmínka}
\texttt{THEN NOT} \emph{závěr}.
\end{defbox}

\begin{defbox}[Příklad --- support, confidence, lift]
Počítejme nad daty z~předchozího příkladu, tedy nad 5 transakcemi:

\emph{Pravidlo 1:} \uv{kdo kupuje chléb, kupuje i~máslo.}
$\mathrm{sup}=3/5=60\,\%$, $\mathrm{conf}=3/4=75\,\%$.

\emph{Pravidlo 2:} \uv{kdo kupuje kávu, kupuje i~máslo.}
$\mathrm{sup}=2/5=40\,\%$, $\mathrm{conf}=2/3=66\,\%$.

\emph{Pravidlo 3:} \uv{kdo kupuje mléko, kupuje i~máslo.}
$\mathrm{sup}=1/5=20\,\%$, $\mathrm{conf}=1/1=100\,\%$,
$\mathrm{lift}=\frac{1/5}{(1/5)\cdot(4/5)}=\frac54=1{,}25$.

Pravidlo 3 vypadá na první pohled nejlépe, protože má \emph{stoprocentní} spolehlivost. Jenže
jeho support je jen 20~\%, takže se opírá o~jedinou transakci. A~lift 1,25 dodává, že zdvih
proti náhodě je malý: máslo kupuje 80~\% zákazníků tak jako tak. To je základní past celé
metody, totiž že \emph{vysoká confidence sama o~sobě nestačí}.
\end{defbox}

\subsubsection{Volba položek, taxonomie a~virtuální položky}

Než se dá cokoli dolovat, musí se rozhodnout, co vlastně bude položkou. Transakční data bývají
méně kvalitní než jiné zdroje a~vyžadují rozsáhlejší předzpracování, a~navíc se relevance
položek mění v~čase. Volba úrovně zpracování je proto kompromisem: čím konkrétnější položky,
tím přímo použitelnější výsledky, ale tím rychleji roste i~počet kombinací a~tím obtížněji se
udrží dostatečně vysoký support.

Uspořádanou odpovědí na tento kompromis je \textbf{taxonomie}, tedy hierarchie kategorií.
Zpočátku se používají obecnější položky a~teprve později se generují pravidla pro položky
konkrétní, a~to jen z~těch transakcí, které je obsahují. Aby byly výsledky použitelné, mají se
uvažované položky vyskytovat v~přibližně stejném počtu transakcí. Proto se \emph{vzácné
položky posunou výše} v~taxonomii, kde se vyskytují relativně častěji, kdežto \emph{častější
položky se nechají na nižších úrovních}, aby nejfrekventovanější položky pravidlům
nedominovaly.

\textbf{Virtuální položky} jdou ještě dál a~přesahují rámec tradiční taxonomie: rozostřují
hranici mezi typy produktů (například značky, řekněme Calvin Klein) a~mohou nést informaci
o~transakci samotné. Ta bývá buď anonymní, jako den v~týdnu nebo čas, nebo podepsaná, tedy
informace o~zákazníkovi a~jeho chování. Za tuto volnost se ale platí \emph{redundancí}
v~pravidlech, a~to dvojím způsobem. Buď virtuální položka odpovídá právě jedné položce
z~taxonomie (\uv{IF produkt\_Škoda THEN Škoda}), nebo se v~pravidle vyskytne virtuální
i~obecná položka současně: místo \uv{IF Fabia THEN stan} vyjde \uv{IF produkt\_Škoda AND malé
auto THEN stan}.

Virtuální položky umožňují i~\textbf{srovnání prodejen}. Zavede se virtuální položka, která
označí typ prodejny, kde transakce proběhla, a~neodpovídá tedy žádnému produktu ani službě.
Postup má čtyři kroky: \emph{(1)}~poskytnout transakční data z~nových prodejen za určitou
dobu; \emph{(2)}~poskytnout přibližně stejné množství dat z~již existujících prodejen;
\emph{(3)}~aplikovat MBA na oba soubory transakcí; \emph{(4)}~zaměřit se na pravidla
obsahující virtuální položky.

Zbývá se zbavit položek, které se objevují příliš zřídka. \textbf{Prořezávání podle
minimálního supportu} málo frekventované položky eliminuje, protože podniknuté akce se musí
týkat dostatečného počtu transakcí. Nabízejí se dvě možnosti: buď se vzácné položky
eliminují a~s~nimi i~odpovídající pravidla, nebo se použije taxonomie a~vytvoří se obecné
položky, které musí splnit stanovený prah. Pozor na to, že proměnný minimální support vede ke
\emph{kaskádovému efektu}.

\subsubsection{Algoritmus Apriori}

Prohledávat všechny kombinace položek nelze, protože jich je exponenciálně mnoho. Apriori
z~toho vybruslí jedinou vlastností, kterou frekventované kombinace mají a~která umožňuje
většinu prostoru vůbec nenavštívit.

\begin{defbox}[Frekventovaná množina položek a~Apriori princip]
\emph{Frekventovaná množina položek} (\emph{frequent itemset}) je kombinace (konjunkce)
kategorií, která na datech dosáhne předem stanovené frekvence, tedy supportu \emph{minsup}.

Klíčovou ideou Apriori je \textbf{vlastnost uzavřenosti dolů} (\emph{downward closure
property}): \emph{každá podmnožina frekventované množiny je také frekventovaná}. Ekvivalentně
řečeno, je-li nějaká množina nefrekventovaná, jsou nefrekventované i~všechny její nadmnožiny.

Při hledání kombinací délky $k$ se proto využijí už známé kombinace délky $k-1$, což je
\emph{prohledávání do šířky}. Aby byla kombinace délky $k$ vůbec vygenerována, požaduje se,
aby \emph{všechny} její podkombinace délky $k-1$ splňovaly požadavek na frekvenci.
\end{defbox}

\begin{defbox}[Algoritmus APRIORI (R.~Agrawal)]
\begin{algorithmic}[1]
\State $L_1\gets$ všechny kategorie, které dosahují alespoň požadované frekvence
\State $k\gets 2$
\While{$L_{k-1}\neq\emptyset$}
  \State $C_k\gets \textsc{Apriori-Gen}(L_{k-1})$ \Comment{generování kandidátů}
  \State $L_k\gets$ ty kombinace z~$C_k$, které dosáhly alespoň požadované frekvence
  \State $k\gets k+1$
\EndWhile
\end{algorithmic}

Vlastní generování kandidátů obstará \textbf{funkce \textsc{Apriori-Gen}$(L_{k-1})$}, která
nejprve kandidáty vytvoří spojováním a~pak ty beznadějné zase vyhodí:
\begin{algorithmic}[1]
\For{každou dvojici kombinací $\mathit{Comb}_p,\mathit{Comb}_q\in L_{k-1}$}
  \If{$\mathit{Comb}_p$ a~$\mathit{Comb}_q$ se shodují v~$k-2$ kategoriích}
     \State přidej $\mathit{Comb}_p\cup \mathit{Comb}_q$ do $C_k$ \Comment{krok spojení}
  \EndIf
\EndFor
\For{každou kombinaci $\mathit{Comb}\in C_k$} \Comment{krok prořezání}
  \If{některá její podkombinace délky $k-1$ není v~$L_{k-1}$}
     \State odstraň $\mathit{Comb}$ z~$C_k$
  \EndIf
\EndFor
\end{algorithmic}
\end{defbox}

\textbf{Generování pravidel.} Nalezením frekventovaných kombinací práce nekončí, protože
kombinace ještě není pravidlo. Z~kombinací s~vyhovující frekvencí se proto teprve vytvoří
asociační pravidla, přičemž platí, že \emph{frekvence nadkombinace} $\le$ \emph{frekvence
kombinace}. Každá kombinace $\mathit{Comb}$ se rozdělí na všechny možné dvojice podkombinací
$\mathit{Ant}$ a~$\mathit{Suc}$ tak, že
\[ \mathit{Suc}=\mathit{Comb}\setminus \mathit{Ant},\qquad
   \mathit{Ant}\cap \mathit{Suc}=\emptyset,\quad \mathit{Ant}\cup \mathit{Suc}=\mathit{Comb}. \]
Pro pravidlo $\mathit{Ant}\Rightarrow \mathit{Suc}$ pak \emph{support} odpovídá frekvenci
kombinace $\mathit{Comb}$ a~\emph{confidence} poměru frekvencí kombinací $\mathit{Comb}$
a~$\mathit{Ant}$. Obojí se tedy dá odečíst z~toho, co už algoritmus spočítal.

Přednosti a~slabiny celé metody se dají shrnout do dvou sloupců:

\begin{center}
\small
\begin{tabular}{@{}p{7.6cm}p{8.0cm}@{}}
\toprule
\textbf{Výhody MBA} & \textbf{Nevýhody MBA} \\
\midrule
jasné a~pochopitelné výsledky --- IF-THEN pravidla okamžitě použitelná & výpočetní náklady:
  prostor asociačních pravidel je exponenciální, $O(2^m)$ pro $m$ položek \\
dobývání bez cílových výstupních hodnot (bez apriorní znalosti) & nutnost taxonomií
  a~virtuálních položek \\
umí zpracovat data proměnné délky & obtížné určení vhodného počtu položek; položky by měly
  mít přibližně stejnou frekvenci \\
jednoduchý a~snadno pochopitelný výpočet & \emph{produkuje nadměrné množství pravidel} ---
  tisíce, desetitisíce, miliony \\
\bottomrule
\end{tabular}
\end{center}

\subsubsection{Vícenásobný minimální support: MS-Apriori}

\textbf{Problém vzácných položek.} Jediný \emph{minsup} mlčky předpokládá, že všechny položky
mají podobnou povahu a~frekvenci. V~praxi tomu tak není: některé položky se vyskytují velmi
často a~jiné jen vzácně, jak je vidět na dvojici kuchyňský robot proti chlebu a~mléku. Ať se
tedy jediná hodnota nastaví jakkoli, něco se pokazí:
\begin{itemize}
\item je-li \emph{minsup} příliš \textbf{vysoký}, pravidla se vzácnými položkami se vůbec
      nenajdou;
\item je-li příliš \textbf{nízký}, aby se našla pravidla se vzácnými i~s~častými položkami,
      hrozí \emph{kombinatorická explozivnost}, protože se časté položky pospojují všemi
      možnými způsoby.
\end{itemize}

\begin{defbox}[Model vícenásobných minimálních supportů (MIS)]
Minimální support pravidla se vyjádří pomocí \emph{minimálních supportů položek}
(\emph{minimum item support}, MIS) těch položek, které v~pravidle vystupují. Každá položka
smí mít vlastní $\mathrm{MIS}(i)$, a~zadáním různých hodnot tak uživatel vyjádří různé
požadavky na support pro různá pravidla.

\textbf{Minsup pravidla} $R: a_1,\dots,a_k\Rightarrow a_{k+1},\dots,a_r$ je
\emph{nejnižší} hodnota MIS mezi jeho položkami. Pravidlo tedy splňuje minimální support
tehdy, je-li jeho skutečný support $\ge \min\big(\mathrm{MIS}(a_1),\dots,\mathrm{MIS}(a_r)\big)$.

\textbf{Omezení rozdílu supportů} má zabránit tomu, aby se v~téže množině sešly velmi časté
a~velmi vzácné položky. Požaduje se proto
\[ \max_{i\in s}\{\mathrm{sup}(i)\}-\min_{i\in s}\{\mathrm{sup}(i)\}\le\varphi, \]
kde $0\le\varphi\le 1$ je uživatelem zadaný maximální rozdíl supportů, stejný pro všechny
množiny.
\end{defbox}

\begin{defbox}[Příklad --- minsup pravidla podle MIS]
Mějme položky \emph{chléb}, \emph{obuv}, \emph{oděvy} s~$\mathrm{MIS}(\text{chléb})=2\,\%$,
$\mathrm{MIS}(\text{obuv})=0{,}1\,\%$, $\mathrm{MIS}(\text{oděvy})=0{,}2\,\%$.

$R_1$: oděvy $\Rightarrow$ chléb $[\mathrm{sup}=0{,}15\,\%]$ minsup \textbf{nesplňuje},
protože $0{,}15\,\%<\mathrm{MIS}(\text{oděvy})=0{,}2\,\%$.

$R_2$: oděvy $\Rightarrow$ obuv $[\mathrm{sup}=0{,}15\,\%]$ minsup naopak \textbf{splňuje},
protože $0{,}15\,\%\ge\mathrm{MIS}(\text{obuv})=0{,}1\,\%$.
\end{defbox}

Za pružnost modelu se ale platí draze: \textbf{uzavřenost dolů v~novém modelu neplatí!} MIS
totiž dovoluje vyšší minimální supporty pro pravidla složená jen z~frekventovaných položek
a~nižší pro pravidla s~položkami méně frekventovanými. Ukažme si to na hodnotách
$\mathrm{MIS}(1)=10\,\%$, $\mathrm{MIS}(2)=20\,\%$, $\mathrm{MIS}(3)=5\,\%$,
$\mathrm{MIS}(4)=6\,\%$. Množina $\{1,2\}$ se supportem 9~\% je nefrekventovaná, přesto
$\{1,2,3\}$ a~$\{1,2,4\}$ frekventované být mohou. Množinu $\{1,2\}$ tedy \emph{nesmíme}
zahodit, i~když by ji klasické Apriori vyřadilo bez rozmýšlení.

\textbf{Řešení} je překvapivě jednoduché: položky se seřadí \emph{vzestupně podle hodnot MIS}
a~toto pořadí se pak používá ve všech dalších operacích algoritmu. Každá množina $w$ má díky
tomu formu $\{w[1],\dots,w[k]\}$, kde $\mathrm{MIS}(w[1])\le\cdots\le\mathrm{MIS}(w[k])$.

\begin{defbox}[Algoritmus MS-Apriori]
\begin{algorithmic}[1]
\State $M\gets \textsc{sort}(I,\mathit{MS})$ \Comment{setřídění položek podle $\mathrm{MIS}(i)$}
\State $L\gets \textsc{init-pass}(M,T)$ \Comment{první průchod daty; $L$ je množina semen}
\State $F_1\gets \{\{i\}\mid i\in L,\ i.\mathrm{count}/n\ge \mathrm{MIS}(i)\}$
\For{$k=2$; $F_{k-1}\neq\emptyset$; $k{+}{+}$}
  \If{$k=2$}
     \State $C_k\gets \textsc{level2-candidate-gen}(L,\varphi)$
  \Else
     \State $C_k\gets \textsc{MScandidate-gen}(F_{k-1},\varphi)$
  \EndIf
  \For{každou transakci $t\in T$}
     \For{každého kandidáta $c\in C_k$}
        \If{$c\subseteq t$} \State $c.\mathrm{count}{+}{+}$ \EndIf
        \If{$c\setminus\{c[1]\}\subseteq t$} \State $c.\mathrm{tailCount}{+}{+}$
           \Comment{pro generování pravidel} \EndIf
     \EndFor
  \EndFor
  \State $F_k\gets \{c\in C_k\mid c.\mathrm{count}/n\ge \mathrm{MIS}(c[1])\}$
\EndFor
\State \Return $F=\bigcup_k F_k$
\end{algorithmic}
\end{defbox}

Jednotlivé kroky si zaslouží komentář. Během \textbf{prvního průchodu daty}
(\textsc{init-pass}) algoritmus zjistí support každé položky. Pak v~setříděném pořadí najde
\emph{první} položku $i$, která splní $\mathrm{MIS}(i)$, a~vloží ji do množiny semen $L$. Pro
každou další položku $j$ za $i$ pak platí jednoduché kritérium: je-li
$j.\mathrm{count}/n\ge \mathrm{MIS}(i)$, patří $j$ do $L$ také.

\emph{Příklad.} Nechť $\mathrm{MIS}(1)=10\,\%$, $\mathrm{MIS}(2)=20\,\%$,
$\mathrm{MIS}(3)=5\,\%$, $\mathrm{MIS}(4)=6\,\%$ a~mějme 100 transakcí s~četnostmi
$\{3\}=6$, $\{4\}=3$, $\{1\}=9$, $\{2\}=25$. Pak $L=\{3,1,2\}$ a~$F_1=\{\{3\},\{2\}\}$.
Položka 4 není v~$L$, protože $4.\mathrm{count}/n<\mathrm{MIS}(3)=5\,\%$; množina $\{1\}$
naopak není v~$F_1$, protože $1.\mathrm{count}/n<\mathrm{MIS}(1)=10\,\%$.

\textbf{Generování kandidátů úrovně 2} (\textsc{level2-candidate-gen}) prochází položky
$l\in L$ v~daném pořadí. Splňuje-li $l$ podmínku $l.\mathrm{count}/n\ge\mathrm{MIS}(l)$,
přidá se pro každou položku $h$ za $l$ kandidát $\{l,h\}$, a~to tehdy, když
$h.\mathrm{count}/n\ge\mathrm{MIS}(l)$ \emph{a} zároveň
$|\mathrm{sup}(h)-\mathrm{sup}(l)|\le\varphi$. Podstatné je, že se pracuje s~$L$, a~nikoli
s~$F_1$: $F_1$ totiž neobsahuje položky, které splňují MIS nějaké dřívější položky, ale ne
své vlastní, což je v~uvedeném příkladu položka 1.

\textbf{Generování kandidátů vyšších úrovní} (\textsc{MScandidate-gen}) má dvě části.
\emph{Krok spojení} je stejný jako v~Apriori: spojí se $f_1=\{i_1,\dots,i_{k-2},i_{k-1}\}$
a~$f_2=\{i_1,\dots,i_{k-2},i'_{k-1}\}$, které se liší jen v~poslední položce, pokud
$i_{k-1}<i'_{k-1}$ a~$|\mathrm{sup}(i_{k-1})-\mathrm{sup}(i'_{k-1})|\le\varphi$.
V~\emph{kroku prořezání} se prochází každá $(k-1)$-podmnožina $s$ kandidáta $c$: není-li
$s\in F_{k-1}$, lze $c$ smazat. Existuje ovšem \textbf{jediná výjimka}. Neobsahuje-li $s$
položku $c[1]$, tedy položku s~nejnižší hodnotou MIS (a~taková $s$ je právě jedna), nesmíme
$c$ smazat. Nevíme totiž, zda $s$ nesplňuje
$\mathrm{MIS}(c[1])$; víme jen, že nesplňuje $\mathrm{MIS}(c[2])$, a~to by k~smazání stačilo,
jen pokud $\mathrm{MIS}(c[2])=\mathrm{MIS}(c[1])$.

\textbf{Proč \texttt{tailCount}.} U~Apriori platí, že je-li $f$ frekventovaná
a~$f_{sub}\subseteq f$, je frekventovaná i~$f_{sub}$, takže k~určení confidence pravidel
stačí supporty frekventovaných množin. U~MS-Apriori to nestačí. Problém vzniká tehdy, když by
v~závěru pravidla stála položka s~\emph{nejnižší} hodnotou MIS. Proto se zvlášť počítá
$c.\mathrm{tailCount}$, tedy support množiny $c$ bez její první položky.

\textbf{Jak zadat hodnoty MIS.} Nabízejí se dva postupy. Buď \emph{(a)}~se hodnota odvodí ze
skutečného supportu položky v~datech, například jako $\lambda\cdot\mathrm{sup}(i)$ pro
parametr $0\le\lambda\le 1$ stejný pro všechny položky, nebo \emph{(b)}~se položky seskupí do
bloků s~podobnými frekvencemi a~všem položkám z~bloku se dá stejná hodnota MIS. Nastavením
$\mathrm{MIS}=100\,\%$ (nebo více) se generování pravidel obsahujících pouze tyto položky
fakticky zakáže, což je praktický trik. Stojí za zapamatování, že model vícenásobných
minsupů \emph{zahrnuje} model jediného supportu jako speciální případ.

\subsubsection{Třídní asociační pravidla (CAR)}

Běžné dolování asociačních pravidel na žádný atribut necílí. Najde všechna pravidla, která
v~datech existují, a~libovolná položka může vystupovat jako podmínka i~jako závěr. V~některých
aplikacích ale uživatele zajímají právě \emph{cílové} atributy či třídy, například které
slovo je asociováno s~kterým tématem dokumentu. Pak se vyplatí hledání rovnou zaměřit.

\begin{defbox}[Třídní asociační pravidlo]
Nechť $T$ je množina $n$ transakcí, každá navíc označkovaná třídou $y$; dále nechť $I$ je
množina všech položek, $Y$ množina všech tříd a~$I\cap Y=\emptyset$.

\emph{Třídní asociační pravidlo} (\emph{class association rule}, CAR) je implikace
$X\to y$, kde $X\subseteq I$ a~$y\in Y$. Support a~confidence se definují stejně jako
u~běžných asociačních pravidel.
\end{defbox}

\begin{defbox}[Příklad --- CAR na textové kolekci]
\begin{center}
\small
\begin{tabular}{@{}ll@{}}
\toprule
Dokument & Třída \\
\midrule
1: Student, Učit, Škola & Vzdělávání \\
2: Student, Škola & Vzdělávání \\
3: Učit, Škola, Město, Hra & Vzdělávání \\
4: Baseball, Basketbal & Sport \\
5: Basketbal, Hráč, Divák & Sport \\
6: Baseball, Trenér, Hra, Tým & Sport \\
7: Basketbal, Tým, Město, Hra & Sport \\
\bottomrule
\end{tabular}
\end{center}
Zvolíme-li $\mathit{minsup}=20\,\%$ a~$\mathit{minconf}=60\,\%$, patří mezi nalezená CAR
například tato dvě:
\[ \text{Student, Škola}\to\text{Vzdělávání}\ [\mathrm{sup}=2/7,\ \mathrm{conf}=2/2],\qquad
   \text{Hra}\to\text{Sport}\ [\mathrm{sup}=2/7,\ \mathrm{conf}=2/3]. \]
\end{defbox}

Na rozdíl od běžných asociačních pravidel se CAR dají dolovat \textbf{přímo v~jednom kroku}.
Klíčovou operací je nalezení všech \emph{pravidlových položek} (\emph{rule-items}) tvaru
$(\mathit{condset},y)$, kde $\mathit{condset}\subseteq I$ a~$y\in Y$, které mají support nad
\emph{minsup}. Každá taková pravidlová položka totiž reprezentuje pravidlo
$\mathit{condset}\to y$. Apriori se na generování CAR dá snadno upravit.

Také zde se dá použít idea vícenásobných minimálních supportů, jen se uplatní na úrovni tříd.
Uživatel zadá \emph{různé minimální supporty různým třídám}, a~tím fakticky i~různý minsup
pravidlům každé třídy: pravidla třídy \emph{Ano} mohou mít minsup 5~\% a~pravidla třídy
\emph{Ne} 10~\%. Nastavením minimálního třídního supportu na 101~\% se generování pravidel
pro danou třídu zakáže.

\subsubsection{Dolování sekvenčních vzorů}

Asociační pravidla o~\emph{pořadí} transakcí nevědí nic, protože je vůbec neuvažují. V~mnoha
aplikacích je ale pořadí to podstatné. Už v~MBA je zajímavé vědět, zda lidé kupují položky
v~posloupnosti, tedy nejprve postel a~o~nějakou dobu později povlečení. Při dolování
z~webových logů (\emph{web usage mining}) se pak hledají navigační vzory podle posloupností
návštěv stránek.

\begin{defbox}[Sekvence, velikost, délka, podsekvence]
Nechť $I=\{i_1,\dots,i_m\}$ je množina položek. \emph{Sekvence} je uspořádaný seznam
množin položek. \emph{Množina položek} (\emph{prvek} sekvence) je neprázdná
$X\subseteq I$; sekvenci $s$ značíme $\langle a_1a_2\dots a_r\rangle$, kde $a_i$ je prvek.
Bez újmy na obecnosti předpokládáme položky v~prvku v~lexikografickém pořadí.

\emph{Velikost} sekvence je počet jejích prvků, kdežto \emph{délka} je počet položek;
sekvence délky $k$ se nazývá $k$-sekvence.

Sekvence $s_1=\langle a_1\dots a_r\rangle$ je \emph{podsekvencí} $s_2=\langle b_1\dots
b_v\rangle$ (a~$s_2$ je \emph{nadsekvencí} $s_1$, též $s_2$ \emph{obsahuje} $s_1$),
existují-li celá čísla $1\le j_1<j_2<\cdots<j_r\le v$ taková, že
$a_1\subseteq b_{j_1},\dots,a_r\subseteq b_{j_r}$.
\end{defbox}

\emph{Příklad.} Pro $I=\{1,\dots,9\}$ je $\langle\{3\}\{4,5\}\{8\}\rangle$ podsekvencí
$\langle\{6\}\{3,7\}\{9\}\{4,5,8\}\{3,8\}\rangle$, protože $\{3\}\subseteq\{3,7\}$,
$\{4,5\}\subseteq\{4,5,8\}$ a~$\{8\}\subseteq\{3,8\}$. Naopak
$\langle\{3\}\{8\}\rangle$ \emph{není} obsažena v~$\langle\{3,8\}\rangle$ ani naopak; na tom
je vidět, že na rozdělení položek do prvků záleží. Velikost sekvence
$\langle\{3\}\{4,5\}\{8\}\rangle$ je 3, její délka 4.

\textbf{Úloha} se pak formuluje takto: pro danou množinu $S$ vstupních datových sekvencí,
tedy pro sekvenční databázi, najít všechny sekvence, které mají uživatelem zadaný minimální
support. Takové sekvenci se říká \emph{frekventovaná sekvence} nebo \emph{sekvenční vzor}
a~jejím supportem je podíl datových sekvencí v~$S$, které ji obsahují.

Vzory hledá algoritmus \textbf{GSP} (\emph{Generalized Sequential Pattern}), a~to způsobem
obdobným Apriori: prohledává do šířky s~generováním a~testováním kandidátů, kde kandidáti
délky $k$ vznikají spojováním frekventovaných sekvencí délky $k-1$.

Časový rozměr se dá do MBA vpašovat i~jinak. \textbf{Analýza časových řad pomocí MBA} je
analýza \uv{příčina--následek}: zavedou se nové položky pro transakce, které nastaly
\emph{před} uvažovanou událostí (příčiny) a~\emph{po} ní (následky), a~duplicitní položky se
z~transakcí odstraní. \emph{Okno} je snímek všech dat za určité období, například všechny
transakce z~posledního měsíce, a~slouží k~zachycení trendů u~vzácných položek. Může být nutná
identifikace zákazníka.

Posledním rozšířením jsou \textbf{disociační pravidla} tvaru \uv{IF $A$ AND NOT $B$ THEN
$C$}. Získají se tak, že se zavedou nové položky komplementární k~původním: neobsahuje-li
transakce původní položku, obsahuje její negaci. Nevýhody jsou trojí. Počet položek se
zdvojnásobí, velikost transakcí naroste a~negované položky se vyskytují častěji než původní,
takže pravidla se všemi položkami negovanými se obtížně využívají.

\subsubsection{FP-Growth a~FP-strom}

Apriori pracuje ve stylu \emph{generuj a~testuj} a~za to platí hned dvakrát. Generování
kandidátů je nákladné prostorově i~časově a~nákladné je i~počítání supportu, protože kontrola
podmnožin je výpočetně drahá a~databáze se prochází \emph{opakovaně}, tedy s~vysokou zátěží
na I/O. \textbf{FP-Growth} proto hledá frekventované množiny \emph{bez generování kandidátů},
a~to ve dvou krocích: \emph{(1)}~postaví dvěma průchody daty kompaktní datovou strukturu
\emph{FP-strom} a~\emph{(2)}~extrahuje frekventované množiny přímo z~FP-stromu tím, že jej
projde.

\begin{defbox}[FP-strom]
Uzly FP-stromu odpovídají položkám a~každý má počítadlo. FP-Growth čte transakce po jedné
a~mapuje každou na \emph{cestu} ve stromu. Používá se přitom \emph{fixní pořadí} položek,
takže cesty se mohou překrývat, mají-li transakce stejný \emph{prefix}; v~takovém případě se
počítadla inkrementují. Mezi uzly, které obsahují tutéž položku, se udržují ukazatele tvořící
jednosměrné spojové seznamy. Čím více cest se překrývá, tím vyšší je komprese a~tím spíše se
FP-strom vejde do paměti.
\end{defbox}

\textbf{Konstrukce (2 průchody).} \emph{Průchod 1} zjistí support každé položky, nefrekventované
zahodí a~frekventované seřadí \emph{klesajícím} supportem, řekněme $a,b,c,d,e$; toto pořadí se
pak použije při stavbě stromu, aby se sdílely společné prefixy. \emph{Průchod 2} strom
postaví. Transakce $\{a,b\}$ vytvoří cestu \texttt{null}$\to a\to b$ s~počítadly 1. Transakce
$\{b,c,d\}$ s~ní nemá společný prefix, a~vytvoří proto \emph{disjunktní} cestu
\texttt{null}$\to b\to c\to d$; navíc se přidá ukazatel mezi oba uzly $b$. Transakce
$\{a,c,d,e\}$ už s~první transakcí společný prefix $a$ má, takže se cesty překryjí
a~počítadlo uzlu $a$ se inkrementuje.

\textbf{Velikost FP-stromu.} Obvykle bývá menší než nekomprimovaná data, protože transakce
typicky sdílejí položky, a~tedy i~prefixy. V~\emph{nejlepším případě} obsahují všechny
transakce stejnou množinu položek, což dá jedinou cestu. V~\emph{nejhorším případě} má každá
transakce unikátní množinu položek, strom je pak alespoň tak velký jako původní data a~nároky
na uložení jsou dokonce \emph{vyšší}, protože je nutné uložit ukazatele a~počítadla. Velikost
navíc závisí na pořadí položek a~klesající support je jen \emph{heuristika}, která
k~nejmenšímu stromu vést nemusí.

\begin{defbox}[Extrakce frekventovaných množin z~FP-stromu]
Algoritmus postupuje \emph{zdola nahoru}, od listů ke kořeni, metodou \emph{rozděl a~panuj}:
nejprve hledá frekventované množiny končící na $e$, pak na $de$, \dots, pak na $d$, $cd$
atd.

\begin{enumerate}
\item \textbf{Prefixový podstrom} (\emph{prefix path sub-tree}) končící danou položkou se
      získá pomocí spojových seznamů.
\item Ověř, zda je položka frekventovaná, a~to sečtením počítadel po spojovém seznamu.
      (V~příkladu pro $\mathit{minSup}=2$ vyjde $\mathrm{count}(e)=3$, takže $\{e\}$ je
      frekventovaná množina.)
\item Je-li položka frekventovaná, hledej rekurzivně frekventované množiny končící na
      $de$, $ce$, $be$, $ae$. K~tomu je nutné získat \emph{podmíněný FP-strom}.
\end{enumerate}

\textbf{Podmíněný FP-strom} (\emph{conditional FP-tree}) pro položku $e$ je takový strom,
který by vznikl, kdybychom uvažovali \emph{jen transakce obsahující} $e$ a~pak $e$ ze všech
transakcí odstranili. Získá se z~prefixového podstromu ve třech krocích:
\begin{enumerate}
\item aktualizuj počítadla po prefixových cestách tak, aby odrážela počet transakcí
      obsahujících $e$;
\item odstraň uzly obsahující $e$, protože informace o~nich už není potřeba;
\item odstraň z~prefixových cest nefrekventované položky. Má-li například $b$ support 1,
      znamená to, že $be$ má support 1, takže $be$ je nefrekventovaná a~$b$ lze odstranit.
\end{enumerate}
Podmíněný strom se pak použije k~nalezení množin končících na $de$, $ce$, $ae$; množina $be$
se neuvažuje, protože $b$ v~podmíněném stromu není. Pro každou z~nich se opět získají
prefixové cesty, extrahují se frekventované množiny a~vytvoří se podmíněný strom, a~tak
rekurzivně dál, například $e\to de\to ade$.
\end{defbox}

Také u~FP-Growth se přednosti a~slabiny dají postavit vedle sebe:

\begin{center}
\small
\begin{tabular}{@{}p{7.6cm}p{8.0cm}@{}}
\toprule
\textbf{Výhody FP-Growth} & \textbf{Nevýhody FP-Growth} \\
\midrule
jen 2 průchody daty & FP-strom se nemusí vejít do paměti \\
\uv{komprimuje} data & stavba FP-stromu je nákladná (ale poté se množiny čtou snadno) \\
žádné generování kandidátů, výrazně rychlejší než Apriori & prořezávat lze jen na úrovni
  jednotlivých položek; support lze spočítat teprve po vložení \emph{všech} dat do stromu \\
\bottomrule
\end{tabular}
\end{center}

\mimo{} Doplňkově se ještě rozlišují dva úsporné druhy frekventovaných množin.
\emph{Uzavřená} frekventovaná množina (\emph{closed}) je taková, jejíž žádná vlastní
nadmnožina nemá stejný support, kdežto \emph{maximální} (\emph{maximal}) je taková, jejíž
žádná vlastní nadmnožina není frekventovaná. Maximální množiny tvoří nejkompaktnější
reprezentaci, ale ztrácejí informaci o~supportech podmnožin. Uzavřené množiny jsou kompromis,
protože reprezentují všechny supporty bezeztrátově.

\subsection{Přístupy založené na principu učení s~učitelem}
\label{ssec:supervised}

\subsubsection{Formulace úlohy}

O~učení s~učitelem mluvíme tam, kde ke každému trénovacímu vzoru známe i~správnou odpověď.
Úloha je tím pádem dobře vymezená: máme se naučit vztah mezi popisem vzoru a~jeho správnou
odpovědí tak, aby fungoval i~na vzorech, které jsme nikdy neviděli.

\begin{defbox}[Klasifikace a~predikce]
Je dána trénovací množina $\{(x_1,y_1),\dots,(x_p,y_p)\}$, kde $x_i$ je vektor hodnot
vstupních atributů a~$y_i$ hodnota \emph{cílového} atributu. Úkolem je najít model
$\hat y=f(x)$, který dobře předpovídá cílový atribut i~pro dosud neviděné vstupy.

Podle povahy cílového atributu se úloha rozpadá na dvě. Je-li cílový atribut
\emph{diskrétní}, jde o~\textbf{klasifikaci}; je-li \emph{spojitý}, jde o~\textbf{regresi}
(predikci). Řešení má vždy dvě fáze: nejprve \emph{(1)}~\textbf{indukce} (učení, trénování)
modelu z~trénovacích dat a~teprve potom \emph{(2)}~\textbf{aplikace} modelu na nová data.
\end{defbox}

Zásadní je, že se model hodnotí na \emph{nových} datech. Schopnost naučit se trénovací data
nazpaměť je bezcenná, protože o~chování na datech dosud neviděných nic nevypovídá; takovému
selhání se říká \emph{přeučení} (\emph{overfitting}) a~vrací se jako téma u~každé metody
v~tomto oddílu. Jak se chyba modelu odhaduje a~jakými mírami se kvalita vyjadřuje, popisuje
kap.~\ref{sec:reprez}, odd.~\ref{ssec:eval}.

\subsubsection{Rozhodovací stromy}

Nejnázornějším klasifikátorem je rozhodovací strom: rozhodnutí o~třídě se v~něm rozpadne na
posloupnost jednoduchých dotazů na hodnoty jednotlivých atributů.

\begin{defbox}[Rozhodovací strom]
Nechť $D$ je databáze $\{d_1,\dots,d_p\}$, $\{A_1,\dots,A_n\}$ množina atributů
a~$C=\{C_1,\dots,C_m\}$ množina tříd. \emph{Rozhodovací strom} pro $D$ je strom, kde:
\begin{itemize}
\item každý vnitřní uzel je označen některým atributem $A_a$;
\item každá hrana je označena predikátem, který lze aplikovat na atribut rodičovského uzlu;
\item každý list nese označení některé třídy $C_j$.
\end{itemize}
Geometricky strom rozděluje prostor příznaků na \emph{obdélníkové oblasti} a~vzor se
klasifikuje podle toho, do které oblasti padne.
\end{defbox}

\begin{center}
\small
\begin{tabular}{@{}p{7.4cm}p{8.2cm}@{}}
\toprule
\textbf{Výhody} & \textbf{Nevýhody} \\
\midrule
snadná implementace a~interpretace & obtížné zpracování spojitých dat (nutná kategorizace
  $\to$ obdélníkové oblasti) \\
efektivnost & nelze vždy použít --- hranice tříd může být diagonální a~strom s~osově
  zarovnanými řezy ji umí jen schodovitě aproximovat (v~příkladu bankovního úvěru vede
  hranice napříč rovinou \emph{zůstatek na kontě} $\times$ \emph{příjem}) \\
extrakce jednoduchých pravidel & obtížné zpracování chybějících dat \\
použitelnost i~na velké datové soubory & přeučení $\Rightarrow$ nutné \emph{prořezávání} \\
  & neuvažují vzájemné korelace mezi atributy \\
\bottomrule
\end{tabular}
\end{center}

Ustálená terminologie pojmenovává tři věci. \emph{Rozhodovací atributy}, podle nichž se
datová množina dělí, označují rozhodovací uzly stromu, kdežto \emph{rozhodovací predikáty}
označují hrany. \emph{Zastavovací kritérium} pak říká, kdy se má větvení ukončit; typicky
zní \uv{všechny vzory z~redukované množiny patří do téže třídy}.

\begin{defbox}[Algoritmus TDIDT (\emph{Top Down Induction of Decision Trees})]
Strom se indukuje shora dolů metodou \emph{rozděl a~panuj}:
\begin{enumerate}
\item Vyber jeden atribut a~umísti jej jako \emph{kořen} částečného (pod)stromu.
\item Rozděl data přicházející do tohoto uzlu na podmnožiny podle hodnot vybraného atributu
      a~pro každou podmnožinu přidej nový uzel.
\item Existuje-li uzel takový, že vzory do něj přicházející nepatří všechny do téže třídy,
      zopakuj pro něj celý proces od kroku 1; jinak skonči.
\end{enumerate}
\end{defbox}

\begin{defbox}[Entropie a~volba dělicího atributu]
Cílem je vybrat takový atribut, který nejlépe oddělí vzory z~různých tříd. Mírou
neuspořádanosti, tedy \uv{překvapení} v~systému, je \emph{entropie}
\[ H=-\sum_{j=1}^{m}p_j\log_2 p_j, \]
kde $p_j$ je pravděpodobnost výskytu třídy $j$, odhadnutá relativní frekvencí na uvažované
množině vzorů, a~$m$ je počet tříd. Platí konvence $0\log_2 0=0$ a~pro převod z~dekadického
logaritmu vztah $\log_2 x=\log_{10}x/0{,}301$.

Atribut $A$ se hodnotí ve dvou krocích. Nejprve se pro každou jeho možnou hodnotu $v$
spočítá entropie $H(A(v))$ na množině příkladů pokrytých kategorií $A(v)$,
\[ H\big(A(v)\big)=-\sum_{j=1}^{m}p_j^{A(v)}\log_2 p_j^{A(v)}, \]
a~z~těchto dílčích hodnot se pak jako vážený součet sestaví \emph{vážená průměrná entropie}
atributu
\[ H(A)=\sum_{v}\frac{|A(v)|}{|D|}\,H\big(A(v)\big). \]
Vybírá se atribut s~\emph{minimální} hodnotou $H(A)$.
\end{defbox}

\begin{defbox}[Příklad --- žádost o~úvěr]
\begin{center}
\small
\begin{tabular}{@{}lllllc@{}}
\toprule
Klient & Příjem & Konto & Pohlaví & Nezaměstnaný & Úvěr \\
\midrule
C1 & vysoký & vysoké & žena & ne & ano \\
C2 & vysoký & vysoké & muž & ne & ano \\
C3 & nízký & nízké & muž & ne & ne \\
C4 & nízký & vysoké & žena & ano & ano \\
C5 & nízký & vysoké & muž & ano & ano \\
C6 & nízký & nízké & žena & ano & ne \\
C7 & vysoký & nízké & muž & ne & ano \\
C8 & vysoký & nízké & žena & ano & ano \\
C9 & nízký & střední & muž & ano & ne \\
C10 & vysoký & střední & žena & ne & ano \\
C11 & nízký & střední & žena & ano & ne \\
C12 & nízký & střední & muž & ne & ano \\
\bottomrule
\end{tabular}
\end{center}

\textbf{Krok 1 --- výběr atributu podle minimální entropie.} Začneme atributem
\emph{příjem} a~sestavíme pro něj a~pro \emph{úvěr} čtyřpolní tabulku: při vysokém příjmu
je 5 klientů s~úvěrem a~0 bez něj, při nízkém příjmu 3 s~úvěrem a~4 bez.
\begin{align*}
H(\text{příjem}(\text{vysoký}))&=-\tfrac55\log_2\tfrac55-\tfrac05\log_2\tfrac05=0,\\
H(\text{příjem}(\text{nízký}))&=-\tfrac37\log_2\tfrac37-\tfrac47\log_2\tfrac47=0{,}9852,\\
H(\text{příjem})&=\tfrac{5}{12}\cdot 0+\tfrac{7}{12}\cdot 0{,}9852=0{,}5747.
\end{align*}
Stejným postupem se spočítají zbylé tři atributy:
$H(\text{konto})=\tfrac13\cdot 0+\tfrac13\cdot 1+\tfrac13\cdot 1=0{,}6667$,
$H(\text{pohlaví})=\tfrac12\cdot 0{,}9183+\tfrac12\cdot 0{,}9183=0{,}9183$
a~$H(\text{nezaměstnaný})=\tfrac12\cdot 1+\tfrac12\cdot 0{,}65=0{,}825$.
Nejmenší entropii má \textbf{příjem}, podle něj se tedy strom větví poprvé.

\textbf{Krok 2.} Ve větvi příjem(vysoký) mají všechny vzory úvěr(ano), a~vzniká proto
rovnou list. Ve větvi příjem(nízký) zůstává 7~klientů, takže je nutné větvit dál. Vychází
$H(\text{konto})=\tfrac27 H(\dots)+\tfrac37 H(\dots)+\tfrac27
H(\dots)=0{,}3935$, $H(\text{pohlaví})=0{,}965$
a~$H(\text{nezaměstnaný})=0{,}9792$, takže se větví podle \textbf{konta}.

\textbf{Krok 3.} Větev konto(vysoké) obsahuje samá ano a~větev konto(nízké) samá ne, obě se
tedy uzavřou listem. Ve větvi konto(střední) zbývají 3~klienti a~z~hodnot
$H(\text{pohlaví})=\tfrac23\cdot 1+\tfrac13\cdot 0=0{,}6667$ a~$H(\text{nezaměstnaný})=0$
plyne třetí větvení podle \textbf{nezaměstnanosti}.

\textbf{Výsledný strom a~převod na pravidla.} Přečtením cest od kořene k~listům dostaneme
pět pravidel:
\begin{align*}
&\texttt{IF } \text{příjem(vysoký)} \texttt{ THEN } \text{úvěr(ano)}\\
&\texttt{IF } \text{příjem(nízký)}\wedge\text{konto(vysoké)} \texttt{ THEN } \text{úvěr(ano)}\\
&\texttt{IF } \text{příjem(nízký)}\wedge\text{konto(nízké)} \texttt{ THEN } \text{úvěr(ne)}\\
&\texttt{IF } \text{příjem(nízký)}\wedge\text{konto(střední)}\wedge\text{nezam.(ano)}
   \texttt{ THEN } \text{úvěr(ne)}\\
&\texttt{IF } \text{příjem(nízký)}\wedge\text{konto(střední)}\wedge\text{nezam.(ne)}
   \texttt{ THEN } \text{úvěr(ano)}
\end{align*}
\end{defbox}

Obecně platí, že každá cesta od kořene k~listu odpovídá právě jednomu pravidlu. Vnitřní uzly
na cestě spolu s~hodnotou uvedenou na odpovídající hraně tvoří levou stranu pravidla, tedy
jeho tělo, a~list s~cílovou třídou tvoří stranu pravou, tedy hlavu.

Kvalitu výsledného stromu ovlivňuje celá řada rozhodnutí, která se během indukce dělají:

\begin{itemize}
\item \emph{Volba rozhodovacích atributů} rozhoduje o~efektivitě stromu; užitečný přitom bývá
      \uv{informovaný} vstup experta na danou oblast.
\item \emph{Pořadí rozhodovacích atributů} umožňuje vyloučit redundantní porovnání.
\item \emph{Větvení} určuje, kolik dělení je vůbec nutné provést; náročnější je u~spojitých
      hodnot nebo u~atributů s~mnoha hodnotami.
\item \emph{Forma stromu} rozhoduje o~jeho hloubce. Vhodnější jsou vyvážené stromy s~co
      nejmenším počtem úrovní, byť za cenu složitějších porovnání; některé algoritmy tvoří
      pouze binární stromy, které bývají hluboké.
\item \emph{Zastavovací kritéria} určují, kdy indukce skončí. Základním kritériem je správná
      klasifikace trénovacích dat, \emph{předčasné zastavení} navíc zabrání budování velkých
      stromů a~přeučení, takže jde o~kompromis mezi přesností a~efektivitou. Za volbou
      jednodušších stromů stojí \emph{Occamova břitva} (Vilém Occam, 1320), tedy preference
      nejjednodušší hypotézy pro $D$.
\item \emph{Trénovací data a~jejich množství} působí na obě strany: příliš málo dat znamená,
      že strom nemusí být dost spolehlivý pro obecnější data, příliš mnoho dat zase přináší
      nebezpečí přeučení.
\item \emph{Prořezávání} (\emph{pruning}) zvyšuje přesnost i~efektivitu na testovacích datech.
      Odstraňují se jím redundantní porovnání, celé podstromy i~irelevantní atributy.
      Prořezávat lze buď už během indukce, což zabrání tvorbě nadměrně velkých stromů,
      nebo až na hotovém stromu.
\item \emph{Časová a~prostorová složitost} závisí na počtu trénovacích dat $p$, počtu atributů
      $n$ a~formě stromu. Indukce stromu stojí $O(n\,p\log p)$, klasifikace $q$ vzorů stromem
      hloubky $O(\log p)$ pak $O(q\log p)$.
\end{itemize}

Konkrétní podobu dává tomuto schématu \textbf{algoritmus ID3} (\emph{Iterative Dichotomiser},
Ross Quinlan, 1986). Učí funkce s~booleovskými hodnotami a~staví strom shora dolů
\emph{hladově}: v~každém uzlu zvolí ten atribut, který lokální trénovací data klasifikuje
nejlépe, a~za nejlepší se považuje atribut s~\emph{největším informačním ziskem}. Stavba
pokračuje, dokud nejsou všechna trénovací data správně klasifikována nebo dokud nejsou
využity všechny atributy.

\begin{defbox}[Algoritmus ID3 a~informační zisk]
\textsc{ID3}(\emph{EXAMPLES}, \emph{TARGET\_ATTRIBUTE}, \emph{ATTRIBUTES}):
\begin{algorithmic}[1]
\State vytvoř kořen $\mathit{ROOT}$
\If{všechny \emph{EXAMPLES} jsou pozitivní} \Return $\mathit{ROOT}$ s~jediným uzlem
  označeným \uv{$+$} \EndIf
\If{všechny \emph{EXAMPLES} jsou negativní} \Return $\mathit{ROOT}$ s~jediným uzlem
  označeným \uv{$-$} \EndIf
\If{$\mathit{ATTRIBUTES}=\emptyset$} \Return $\mathit{ROOT}$ s~uzlem označeným
  nejčastější hodnotou \emph{TARGET\_ATTRIBUTE} v~\emph{EXAMPLES} \EndIf
\State $A\gets$ atribut z~\emph{ATTRIBUTES}, který nejlépe klasifikuje \emph{EXAMPLES}
\State rozhodovací atribut pro $\mathit{ROOT}\gets A$
\For{každou možnou hodnotu $v_i$ atributu $A$}
   \State připoj k~$\mathit{ROOT}$ novou hranu odpovídající predikátu $A=v_i$
   \State $\mathit{EXAMPLES}_{v_i}\gets$ podmnožina \emph{EXAMPLES} s~hodnotou $v_i$ pro $A$
   \If{$\mathit{EXAMPLES}_{v_i}=\emptyset$}
      \State připoj k~hraně list označený nejčastější hodnotou \emph{TARGET\_ATTRIBUTE}
   \Else
      \State připoj podstrom \textsc{ID3}$(\mathit{EXAMPLES}_{v_i},
             \mathit{TARGET\_ATTRIBUTE}, \mathit{ATTRIBUTES}\setminus\{A\})$
   \EndIf
\EndFor
\State \Return $\mathit{ROOT}$
\end{algorithmic}

Kritériem volby atributu je \textbf{informační zisk} (\emph{information gain}), tedy očekávané
snížení entropie poté, co se data rozdělí podle uvažovaného atributu:
\[ \mathrm{GAIN}(D,A)=H(D)-\sum_{v\in \mathrm{VALUES}(A)}\frac{|D_v|}{|D|}H(D_v). \]
Hodnota $H(D)$ je přitom pro všechny atributy stejná, takže maximalizovat informační zisk
znamená totéž jako minimalizovat váženou entropii $H(A)$.
\end{defbox}

\textbf{Algoritmus C4.5} (Quinlan, 1993) je modifikací ID3 a~doplňuje jej hned v~několika
směrech:

\begin{itemize}
\item \textbf{Chybějící data} se při konstrukci stromu ignorují, tedy při vyhodnocení kritéria
      dělení se uvažují jen známé hodnoty daného atributu. Při klasifikaci se chybějící
      hodnoty odhadují podle hodnot známých u~jiných vzorů.
\item \textbf{Spojitá data} se kategorizují podle hodnot, kterých atribut nabývá na
      trénovacích vzorech.
\item \textbf{Prořezávání validací} vyžaduje rozdělit trénovací data na vlastní trénovací
      ($\tfrac23$) a~validační ($\tfrac13$) množinu. Při metodě \emph{reduced-error pruning}
      se podstromy nahrazují listy označenými nejčastější třídou příslušných trénovacích
      vzorů, přičemž nový strom \emph{nesmí} být na validační množině horší než původní;
      jinak se podstrom nenahradí. Při metodě \emph{subtree raising} se podstrom nahradí svým
      nejčastěji používaným podstromem, který se zvedne na vyšší úroveň, a~testuje se
      klasifikační přesnost.
\item \textbf{Prořezávání pravidel} (\emph{rule post-pruning}) jde jinou cestou. Z~trénovací
      množiny se nejprve odvodí strom, u~něhož je přeučení dovoleno, a~převede se na
      ekvivalentní množinu pravidel, jedno pravidlo na cestu od kořene k~listu. Pravidla se
      pak prořežou, tedy zobecní, vypouštěním všech možných předpokladů, pokud to nezhorší
      odhadovanou klasifikační přesnost, a~nakonec se seřadí podle očekávané přesnosti;
      v~tomto pořadí se pak používají ke klasifikaci. Prořezávání pravidel je snazší
      a~radikálnější než prořezávání celých stromů a~výsledná pravidla jsou průhledná
      a~snadno pochopitelná.
\item \textbf{Odhad přesnosti pravidel} se standardně dělá validační množinou, C4.5 však
      hodnotí přesnost \emph{na trénovacích datech} jako podíl vzorů správně klasifikovaných
      pravidlem ke všem vzorům, které pravidlo pokrývá, $\mathit{acc}=k/n$. Směrodatná
      odchylka přesnosti se odhaduje z~binomického rozdělení a~jako charakteristika pravidla
      se bere \emph{dolní mez} přesnosti pro zvolený interval spolehlivosti, na 95\% hladině
      \[ \mathit{acc}-1{,}96\sqrt{\frac{\mathit{acc}(1-\mathit{acc})}{n}}. \]
\item \textbf{Poměrný informační zisk} (\emph{information gain ratio}) nahrazuje dosavadní
      kritérium dělení:
      \[ \mathrm{GAIN\_RATIO}(D,A)=\frac{\mathrm{GAIN}(D,A)}
         {-\sum_{v}\frac{|D_v|}{|D|}\log_2\frac{|D_v|}{|D|}}. \]
      Dělí se podle atributu s~\emph{největším} poměrným informačním ziskem, přičemž
      jmenovatel penalizuje \emph{nadměrné větvení}, tedy atributy s~mnoha hodnotami.
\end{itemize}

\textbf{Algoritmus C5.0} je komerční verzí C4.5 určenou pro velké databáze. Indukce stromu
probíhá podobně, generování pravidel je ale zlepšené a~efektivnější a~vyšší přesnosti se
dosahuje \emph{boostingem}. Z~původních trénovacích dat se staví různé trénovací podmnožiny
a~každému trénovacímu vzoru se přiřadí váha odpovídající jeho vlivu při klasifikaci. Pro
každou kombinaci vah pak vznikne samostatný klasifikátor a~při klasifikaci má každý
klasifikátor svůj hlas, přičemž vítězí většinová třída.

\begin{defbox}[CART --- \emph{Classification And Regression Trees}]
CART generuje výhradně \emph{binární} stromy: při každém dělení dat se tvoří právě dvě hrany.
Nejlepší rozhodovací atribut se volí podle entropie, nebo podle \textbf{Gini indexu}. Pro
atribut o~$r$ hodnotách $v_1,\dots,v_r$ při $n$ vzorech a~$m$ třídách je
\[ G=\sum_{i=1}^{r}\frac{n_i}{n}G(v_i),\qquad G(v_i)=1-\sum_{j=1}^{m}p_j^2,
   \qquad \sum_{i=1}^{r}n_i=n, \]
kde $n_i$ je počet vzorů s~hodnotou $v_i$ a~$p_j$ pravděpodobnost třídy $j$ v~těchto $n_i$
vzorech. Čím \emph{nižší} je hodnota $G$, tím lepší je diskriminace.

V~uzlu $t$ se dělí podle \uv{nejlepšího} nalezeného bodu, který se hledá průchodem všech
možných hodnot $v$ rozhodovacího atributu $A$; při rovnosti se použije pravý podstrom.
Označme $P_L,P_R$ pravděpodobnosti zpracování levým, resp.\ pravým podstromem
a~$P(C_j\mid t_L),P(C_j\mid t_R)$ pravděpodobnosti, že vzor patří do třídy $C_j$ a~je v~levém,
resp.\ pravém podstromu uzlu $t$. Kritérium dělení pak zní
\[ \Phi=2P_LP_R\sum_{j=1}^{m}\big|P(C_j\mid t_L)-P(C_j\mid t_R)\big|, \]
a~nabývá maxima tehdy, jsou-li všechny vzory dané třídy seskupeny v~témže podstromu.

Pro CART je charakteristických několik vlastností. Pořadí atributů ve stromu odpovídá jejich
vlivu při klasifikaci. Chybějící data se ignorují a~nezapočítávají se do vyhodnocení
rozhodovacích kritérií. Zastavovacím kritériem je zjištění, že žádné další dělení už
klasifikační přesnost nezlepší. Vysoká přesnost na trénovací množině přitom nemusí odrážet
přesnost na datech testovacích.
\end{defbox}

\begin{defbox}[CHAID --- \emph{Chi-square Automatic Interaction Detection}]
Kritériem větvení je zde $\chi^2$. Algoritmus navíc automaticky \emph{seskupuje hodnoty}
kategoriálních atributů: hodnoty se postupně slučují od původního počtu, dokud nezůstanou jen
dvě grupy, a~teprve pak se vybere atribut spolu s~kategorizací, které jsou pro dělení dat
v~daném kroku nejlepší. Při větvení tak nevzniká tolik větví, kolik má atribut hodnot.

\textbf{Seskupování hodnot atributu} probíhá takto:
\begin{enumerate}
\item Opakuj, dokud nezůstanou jen dvě grupy hodnot:
      \emph{(a)}~vyber pár kategorií atributu, které jsou si z~hlediska $\chi^2$ nejvíce
      podobné a~lze je sloučit; \emph{(b)}~novou kategorii uvažuj jako možné seskupení
      v~daném kroku.
\item Pro každou možnou grupu hodnot testuj $\chi^2$ testem její relevanci k~predikované
      proměnné.
\item Vyber seskupení, které dává \uv{nejlepší} rozdělení s~nejvýznamnějším rozštěpením.
\item Urči, zda toto seskupení statisticky významně přispívá k~rozlišení vzorů z~různých tříd.
\end{enumerate}
\end{defbox}

\subsubsection{Bayesovská klasifikace}

Zatímco strom rozhoduje posloupností testů, bayesovský přístup počítá pravděpodobnost
jednotlivých tříd a~vybírá tu nejpravděpodobnější.

\begin{defbox}[Bayesova formule a~hypotéza MAP]
Pravděpodobnost, že hypotéza $H$ platí, pozorujeme-li $E$, je
\[ P(H\mid E)=\frac{P(E\mid H)\,P(H)}{P(E)}, \]
kde $P(H)$ je \emph{apriorní} pravděpodobnost hypotézy, $P(H\mid E)$ \emph{aposteriorní}
(podmíněná) pravděpodobnost a~$P(E)$ pravděpodobnost jevu $E$.

Za \emph{nejpravděpodobnější hypotézu} $H_{MAP}$ (\emph{maximum a~posteriori}) se považuje ta
s~největší aposteriorní pravděpodobností. Jmenovatel $P(E)$ je pro všechny hypotézy stejný,
takže jej lze zanedbat:
\[ H_{MAP}=\arg\max_{H\in\mathcal H} P(H\mid E)=\arg\max_{H\in\mathcal H} P(E\mid H)\,P(H). \]
\end{defbox}

\begin{defbox}[Naivní bayesovský klasifikátor]
\textbf{Předpoklad:} jevy $E_1,\dots,E_k$ jsou \emph{podmíněně nezávislé} při platnosti
hypotézy $H$. V~reálných úlohách bývá tento předpoklad splněn jen zřídka, a~právě odtud
pochází označení \uv{naivní}. Za tohoto předpokladu je
\[ P(H\mid E_1,\dots,E_k)=\frac{P(E_1,\dots,E_k\mid H)P(H)}{P(E_1,\dots,E_k)}
   =\frac{P(H)}{P(E_1,\dots,E_k)}\prod_{i=1}^{k}P(E_i\mid H), \]
a~klasifikuje se hypotézou s~největší aposteriorní pravděpodobností
\[ H_{MAP}=\arg\max_{H_t} P(H_t)\prod_{i=1}^{k}P(E_i\mid H_t). \]
Potřebné pravděpodobnosti se odhadnou přímo z~dat jako
$P(H_t)=P(\text{třída }t)=\frac{\text{počet vzorů z~}t}{\text{počet všech
vzorů}}$ a~$P(E_k\mid H_t)=\frac{\text{počet vzorů z~}t\text{ splňujících }E_k}
{\text{počet vzorů z~}t}$.

\textbf{Výhodou} klasifikátoru je, že si poradí i~s~\emph{nekompletně popsanými} vzory.
\textbf{Nevýhodou} je nulová aposteriorní pravděpodobnost $P(E_k\mid H_t)$ v~případě, že
v~trénovací množině není žádný odpovídající vzor; podobně jsou podhodnocené i~ty členy,
u~nichž se $E_k$ a~$H_t$ vyskytnou společně jen zřídka.
\end{defbox}

\begin{defbox}[Příklad --- bayesovská klasifikace úvěru]
\emph{Jednoduchý případ} vystačí s~jediným pozorováním. Nechť
$P(\text{půjčit})=0{,}667$, $P(\text{zamítnout})=0{,}333$,
$P(\text{vysoký příjem}\mid\text{půjčit})=0{,}91$ a~
$P(\text{vysoký příjem}\mid\text{zamítnout})=0{,}12$. Porovnáním
$0{,}91\cdot 0{,}667=0{,}607$ proti $0{,}12\cdot 0{,}333=0{,}040$ vychází
$H_{MAP}=\text{půjčit}$.

\emph{Naivní klasifikátor} nad daty z~příkladu o~úvěru (12 klientů) potřebuje odhadnout tyto
pravděpodobnosti:
$P(\text{úvěr(ano)})=8/12=0{,}667$, $P(\text{úvěr(ne)})=4/12=0{,}333$;
$P(\text{konto(střední)}\mid\text{ano})=2/8=0{,}25$,
$P(\text{konto(střední)}\mid\text{ne})=2/4=0{,}5$;
$P(\text{nezam.(ne)}\mid\text{ano})=5/8=0{,}625$,
$P(\text{nezam.(ne)}\mid\text{ne})=1/4=0{,}25$.

Pro klienta se středním kontem, který není nezaměstnaný, pak vyjde
\[ 0{,}667\cdot 0{,}25\cdot 0{,}625=0{,}1042 \quad\text{vs.}\quad
   0{,}333\cdot 0{,}5\cdot 0{,}25=0{,}0416, \]
klient se tedy zařadí do třídy \textbf{úvěr(ano)}.
\end{defbox}

\subsubsection{Support Vector Machines}

Metodu SVM vynalezl V.~Vapnik se spolupracovníky už v~70.~letech v~Rusku, na Západě se ale
stala známou teprve v~roce 1992. Základem jsou \emph{lineární} klasifikátory, které hledají
hyperrovinu oddělující dvě třídy, pozitivní a~negativní. Tam, kde lineární separace nestačí,
nastupují \emph{kernelové} funkce.

SVM stojí na důkladném teoretickém základu a~zároveň v~aplikacích dosahuje vyšší přesnosti
než většina jiných metod, a~to zejména na vysokorozměrných datech. Pro klasifikaci textu
patří k~vůbec nejlepším klasifikátorům.

\begin{defbox}[Lineární SVM, separabilní případ]
Trénovací množina $D=\{(x_1,y_1),\dots,(x_r,y_r)\}$, $x_i\in X\subseteq\mathbb R^n$,
$y_i\in\{1,-1\}$. SVM hledá lineární funkci $f(x)=\langle w\cdot x\rangle+b$ tak, že
\[ y_i=\begin{cases}1 & \text{pro }\langle w\cdot x_i\rangle+b\ge 0,\\
-1 & \text{pro }\langle w\cdot x_i\rangle+b<0.\end{cases} \]
Hyperrovina $\langle w\cdot x\rangle+b=0$ se nazývá \emph{rozhodovací hranice}
(\emph{decision boundary}). Takových hyperrovin existuje mnoho a~SVM z~nich vybírá tu
s~\emph{největším okrajem} (\emph{margin}), protože podle teorie strojového učení právě ona
minimalizuje mez chyby.

Vzdálenost bodu $x_i$ od hyperroviny je $\dfrac{|\langle w\cdot x_i\rangle+b|}{\|w\|}$, kde
$\|w\|=\sqrt{\langle w\cdot w\rangle}$. Pro okrajové hyperroviny $H_+$ a~$H_-$ platí
\[ d_+=d_-=\frac{1}{\|w\|},\qquad \text{margin}=d_++d_-=\frac{2}{\|w\|}. \]

Maximalizace okraje se tak převede na \textbf{optimalizační úlohu}
\[ \text{minimalizuj }\ \frac{\langle w\cdot w\rangle}{2}
   \quad\text{za podmínek}\quad y_i\big(\langle w\cdot x_i\rangle+b\big)\ge 1,\
   i=1,\dots,r, \]
kde je omezení zapsáno souhrnně za obě třídy najednou: $\langle w\cdot x_i\rangle+b\ge 1$
pro $y_i=1$ a~$\langle w\cdot x_i\rangle+b\le -1$ pro $y_i=-1$.
\end{defbox}

\begin{defbox}[Duální formulace a~podpůrné vektory]
Úloha se řeší standardní \emph{Lagrangeovou} metodou:
\[ L_P=\frac{\langle w\cdot w\rangle}{2}-\sum_{i=1}^{r}\alpha_i
   \big[y_i(\langle w\cdot x_i\rangle+b)-1\big],\qquad \alpha_i\ge 0. \]
Optimální řešení musí splňovat \emph{Kuhn--Tuckerovy podmínky}, které jsou pro \emph{konvexní}
účelovou funkci s~\emph{lineárními} omezeními nejen nutné, ale i~postačující. Z~podmínky
komplementarity plyne, že $\alpha_i>0$ mohou mít jen body ležící \emph{na} okrajových
hyperrovinách, neboť právě pro ně je $y_i(\langle w\cdot x_i\rangle+b)-1=0$. Těmto bodům se
říká \textbf{podpůrné vektory} (\emph{support vectors}); pro všechny ostatní body je
$\alpha_i=0$.

Vynulováním parciálních derivací $L_P$ podle $w$ a~$b$ a~zpětnou substitucí dostaneme
\emph{Wolfeho duální} úlohu
\[ L_D=\sum_{i=1}^{r}\alpha_i-\frac12\sum_{i,j}y_iy_j\alpha_i\alpha_j
   \langle x_i\cdot x_j\rangle\ \longrightarrow\ \max, \]
jejíž maximum se nabývá pro tytéž hodnoty $w,b,\alpha_i$ jako minimum $L_P$. Numericky se
duální úloha řeší snáz než úloha primární.

Tvar \textbf{rozhodovací hranice} a~způsob \textbf{testování} vzoru $z$ jsou pak
\[ \langle w\cdot x\rangle+b=\sum_{i\in \mathit{sv}}y_i\alpha_i\langle x_i\cdot x\rangle+b=0,
\qquad
\mathrm{sign}\Big(\sum_{i\in \mathit{sv}}\alpha_iy_i\langle x_i\cdot z\rangle+b\Big). \]
Vrátí-li výraz 1, je $z$ klasifikován jako pozitivní, jinak jako negativní.
\end{defbox}

\begin{defbox}[Kernelový trik]
Reálná data mohou vyžadovat nelineární separaci. Řešením je \emph{transformovat} je do jiného,
obvykle mnohem vícerozměrného prostoru, v~němž už třídy oddělí \emph{lineární} hranice:
$\phi:X\to F$, $x\mapsto\phi(x)$. Transformovaný prostor $F$ se nazývá \emph{prostor
příznaků} (\emph{feature space}), původní prostor \emph{vstupní}.

\textbf{Problém explicitní transformace} spočívá v~tom, že počet dimenzí prostoru příznaků
může být pro užitečné transformace enormní i~při rozumném počtu vstupních atributů; explicitní
výpočet se tak stává kvůli \emph{prokletí dimenzionality} nezvládnutelným.

\textbf{Klíčové pozorování} je, že v~duální formulaci vyžadují jak konstrukce optimální
hyperroviny, tak vyhodnocení rozhodovací funkce \emph{pouze skalární součiny}
$\langle\phi(x)\cdot\phi(z)\rangle$, nikdy transformovaný vektor $\phi(x)$ samotný. Umíme-li
tedy počítat skalární součin přímo ze vstupních vektorů, nemusíme $\phi$ znát vůbec. Takové
funkci se říká \emph{kernelová funkce}
\[ K(x,z)=\langle\phi(x)\cdot\phi(z)\rangle. \]
Nahradit ve všech výrazech skalární součin kernelem znamená použít \textbf{kernelový trik}.

\emph{Příklad (polynomiální kernel $d=2$, 2-D vstup):}
\begin{align*}
\langle x\cdot z\rangle^2&=(x_1z_1+x_2z_2)^2=x_1^2z_1^2+2x_1z_1x_2z_2+x_2^2z_2^2\\
&=\big\langle (x_1^2,x_2^2,\sqrt2 x_1x_2)\cdot(z_1^2,z_2^2,\sqrt2 z_1z_2)\big\rangle
 =\langle\phi(x)\cdot\phi(z)\rangle.
\end{align*}
Kernel $\langle x\cdot z\rangle^d$ je tedy skalární součin v~transformovaném prostoru, aniž
bychom $\phi$ kdy vyčíslili. Na otázku, zda je daná funkce kernelem, tj.\ zda je skalárním
součinem v~nějakém prostoru příznaků, odpovídá \emph{Mercerova věta}. Mezi běžně používané
kernely patří lineární (identita), polynomiální, gaussovský (RBF) a~sigmoidální.
\end{defbox}

Metoda má i~několik omezení, která je dobré znát. SVM pracuje jen v~reálném prostoru, takže
kategoriální atributy je nutné převést na numerické, a~provádí jen \emph{dvoutřídní}
klasifikaci. Pro více tříd se proto sahá po strategii \emph{jeden proti ostatním}
(\emph{one-against-rest}) nebo po kódech opravujících chyby (\emph{error-correcting output
coding}). Hyperrovina, kterou SVM vyprodukuje, je navíc pro člověka obtížně pochopitelná
a~kernely to ještě zhoršují; SVM se proto používá tam, kde se srozumitelnost pro člověka
nevyžaduje. \mimo{} Pro nelineárně separabilní data se zavádějí \emph{proměnné vůle}
(\emph{slack variables}) $\xi_i\ge 0$ a~minimalizuje se
$\frac{\langle w\cdot w\rangle}{2}+C\sum_i\xi_i$, kde parametr $C$ řídí kompromis mezi šířkou
okraje a~počtem chyb.

\subsubsection{Logistická regrese}
\label{ssec:logreg}

Logistická regrese má dvojí tvář. Je to nejdůležitější případ nelineární regrese
(kap.~\ref{sec:priprava}) a~zároveň plnohodnotný klasifikátor: závislá proměnná $y$ je
kategoriální, typicky dvouhodnotová, a~model nepředpovídá přímo její hodnotu, nýbrž
\emph{pravděpodobnost}, že $y$ nabude určité hodnoty, a~to v~závislosti na kombinaci hodnot
nezávislých veličin.

\begin{defbox}[Logistická regrese]
Pro binární $y\in\{0,1\}$ se modeluje logaritmus \uv{podmíněné šance}
(\emph{conditional odds}):
\[ \ln\frac{P(y=1\mid x_1,\dots,x_m)}{1-P(y=1\mid x_1,\dots,x_m)}
   =\beta_0+\beta_1x_1+\cdots+\beta_mx_m=z, \]
odkud
\[ P(y=1\mid x)=\frac{e^{z}}{1+e^{z}}=\frac{1}{1+e^{-z}}=\sigma(z)
   \qquad\text{(sigmoida).} \]

\textbf{Učení} probíhá maximalizací věrohodnosti (\emph{maximum likelihood}):
\[ L=\prod_{i=1}^{p}P\big(y_i\mid x_{i,1},\dots,x_{i,m}\big)
   =\prod_{i=1}^{p}\Big(\frac{1}{1+e^{-z_i}}\Big)^{y_i}
     \Big(\frac{1}{1+e^{z_i}}\Big)^{1-y_i}. \]
Místo maximalizace $L$ se v~praxi minimalizuje \emph{křížová entropie} $E=-\log L$:
\[ E=-\sum_i\Big[y_i\log\sigma(z_i)+(1-y_i)\log\big(1-\sigma(z_i)\big)\Big]. \]
Derivace sigmoidy má tvar $\sigma'=\sigma(1-\sigma)$, takže po zjednodušení vyjdou velmi
jednoduchá adaptační pravidla
\[ \Delta\beta_j=-\eta\frac{\partial E}{\partial\beta_j}
   =\eta\sum_i\big(y_i-\sigma(z_i)\big)x_{i,j},\qquad
   \Delta\beta_0=\eta\sum_i\big(y_i-\sigma(z_i)\big), \]
kde $\eta$ je rychlost učení. Chyba tedy vstupuje do úpravy vah přímo jako rozdíl
\emph{pozorované} a~\emph{predikované} hodnoty.
\end{defbox}

Logistická regrese bývá běžnou volbou základního klasifikátoru všude tam, kde se požaduje
pravděpodobnostní výstup a~interpretovatelné koeficienty. Používá se například při predikci
vazeb v~sociálních sítích (kap.~\ref{sec:sna}).

\subsubsection{Diskriminační analýza}

\begin{defbox}[Diskriminační analýza]
Úlohou je \emph{klasifikace vzorů do známých tříd}, a~hledá se proto závislost nominální
veličiny, která určuje příslušnost ke třídě, na ostatních numerických veličinách.
Předpokládáme, že ke každé třídě $C_j$, $j=1,\dots,m$ existuje \emph{diskriminační funkce}
$g_j$; vzor $x=(x_1,\dots,x_m)$ patří do třídy $C_j$, právě když
\[ g_j(x)=\max_{k} g_k(x). \]

\textbf{Lineární diskriminační analýza} volí
$g_j(x)=\beta_{j0}+\beta_{j1}x_1+\cdots+\beta_{jm}x_m$. Jsou-li třídy jen dvě, stačí místo
$g_1,g_2$ hledat jedinou funkci $g=g_1-g_2$ a~klasifikovat podle jejího \emph{znaménka}:
\uv{$+$} znamená třídu 1, \uv{$-$} třídu 2.

\textbf{Optimální klasifikace ve smyslu minimální chyby} nastane tehdy, jsou-li diskriminační
funkce $\propto$ aposteriorní (podmíněné) pravděpodobnosti zařazení pozorování $x$ do
třídy $C_j$, tedy $g_j(x)=P(C_j\mid x)=\dfrac{p(x\mid C_j)P(C_j)}{p(x)}$; pro dvě třídy
$g(x)=p(x\mid C_1)P(C_1)-p(x\mid C_2)P(C_2)$.
\end{defbox}

Konkrétní tvar diskriminační funkce závisí na rozdělení dat. Pro \emph{normální rozdělení}
s~různými středními hodnotami $\mu_1,\mu_2$ a~obecnými kovariančními maticemi $S_1,S_2$
vychází funkce \emph{kvadratická}. Jsou-li kovarianční matice \emph{shodné}
($S_1=S_2=S$), redukuje se na \emph{lineární}
\[ g(x)=(\mu_1-\mu_2)^{\!\top}S^{-1}x-\tfrac12(\mu_1-\mu_2)^{\!\top}S^{-1}(\mu_1+\mu_2)
   -\ln\frac{P(C_2)}{P(C_1)}, \]
a~jsou-li navíc obě třídy stejně pravděpodobné a~$S=I$, zbude jednoduché srovnávání
vzdáleností od středních hodnot. Pro normální rozdělení se tedy celé hledání diskriminačních
funkcí redukuje na \textbf{odhad středních hodnot} výběrovými průměry a~\textbf{kovariančních
matic} korigovanými výběrovými rozptyly.

\subsubsection{Neuronové sítě a~extrémní učicí stroje (ELM)}

\begin{defbox}[Síť ELM (\emph{Extreme Learning Machine}, G.-B.~Huang)]
Jde o~dopřednou síť s~\emph{jednou} skrytou vrstvou: $n$ vstupních neuronů, $L$ skrytých
a~1 výstupní. Výstup skrytých neuronů je
$h(x)=\big[G(a_1,b_1,x),\dots,G(a_L,b_L,x)\big]$, kde $a_i$ je váhový vektor a~$b_i$ bias
$i$-tého neuronu; například $G(a_i,b_i,x)=g(\langle a_i\cdot x\rangle+b_i)$ (sigmoida) nebo
$G(a_i,b_i,x)=g(b_i\|x-a_i\|)$ (RBF). Výstup celé sítě je
$f_L(x)=\sum_{i=1}^{L}\beta_iG(a_i,b_i,x)$.

\textbf{Klíčová myšlenka:} lze-li spojitou funkci $f$ aproximovat dopřednou sítí se skrytou
vrstvou neuronů počítajících $G$, pak \emph{není nutné parametry skrytých neuronů adaptovat}.
Mohou být vygenerovány \emph{náhodně}, dokonce bez znalosti trénovacích dat. Pro každou
spojitou $f$ a~náhodně generovanou posloupnost $\{(a_i,b_i)\}$ totiž platí
s~pravděpodobností 1
\[ \lim_{L\to\infty}\Big\|f-\textstyle\sum_{i=1}^{L}\beta_iG(a_i,b_i,\cdot)\Big\|=0 \]
při nastavení $\beta_i$ tak, aby minimalizovaly $\|f-f_L\|$.

\textbf{Maticový model:} soustava $\sum_{i=1}^{L}\beta_iG(a_i,b_i,x_j)=t_j$,
$j=1,\dots,N$ je ekvivalentní zápisu $H\beta=T$, kde $H$ je \emph{výstupní matice skryté
vrstvy}; její $j$-tý řádek tvoří výstupy skrytých neuronů pro vzor $x_j$ a~$i$-tý sloupec
výstupy $i$-tého neuronu pro všechny vzory.

\textbf{Algoritmus učení} se tím smrskne na tři neiterativní kroky:
\begin{algorithmic}[1]
\State zvol \emph{náhodně} parametry skrytých neuronů $(a_i,b_i)$, $i=1,\dots,L$
\State spočítej výstupní matici skryté vrstvy $H$
\State spočítej výstupní váhy $\beta=H^{+}T$
\end{algorithmic}
kde $H^{+}$ je Mooreova--Penroseova pseudoinverze $H$. Stabilnější řešení s~lepší generalizací
dá regularizace, tedy přičtení $I/C$ na diagonálu $H^{\!\top}H$:
$f(x)=h(x)\big(I/C+H^{\!\top}H\big)^{-1}H^{\!\top}T$.
\end{defbox}

Výhody ELM plynou přímo z~té tříkrokové procedury. Učicí procedura je jednoduchá
a~neiterativní, zatímco tradiční algoritmy jsou podstatně složitější. Trénování je díky tomu
velmi rychlé a~přímočaré: odpadá řešení lokálních minim, hledání vhodné inicializace parametrů
i~starost o~přeučení. Parametry skrytých neuronů $a_i,b_i$ jsou nezávislé na trénovacích datech
i~na sobě navzájem, takže je lze vygenerovat \emph{už před} předložením trénovacích dat;
konvenční techniky naproti tomu data \uv{vidět} musí. Poslední výhodou je volnost ve výběru
přenosové funkce, protože ELM je použitelný s~\emph{libovolnou} omezenou po částech spojitou
nekonstantní přenosovou funkcí, kdežto gradientní algoritmy vyžadují funkce hladké.

\subsubsection{Ensemble metody}

Ensemble metody nespoléhají na jediný dokonalý model, ale kombinují predikce mnoha
klasifikátorů. Liší se přitom v~tom, kterou složku chyby snižují.

\begin{defbox}[Bagging (\emph{bootstrapped aggregating})]
Bagging se snaží snížit \emph{rozptyl} klasifikátoru. Je-li rozptyl jednoho klasifikátoru
$\sigma^2$, pak rozptyl průměru $k$ nezávislých a~stejně rozdělených (i.i.d.) klasifikátorů
klesne na $\sigma^2/k$.

Takové i.i.d.\ klasifikátory se aproximují \textbf{bootstrappingem}: body se vybírají
\emph{s~opakováním} uniformně z~původních dat. Vzorek má přibližně stejnou velikost jako
původní data, může obsahovat duplikáty a~obsahuje podíl
\[ 1-\Big(1-\frac1n\Big)^{\!n}\approx 1-\frac1e=63{,}2\,\% \]
různých bodů z~původních dat, neboť pravděpodobnost, že bod ve vzorku \emph{není}, je
$(1-1/n)^n$. Vylosuje se $k$ nezávislých bootstrapových vzorků velikosti $n$, na každém se
natrénuje klasifikátor a~predikce pro testovací vzor se určí \emph{většinovým hlasem}.

\textbf{Vhodnost:} ideálním základem pro bagging jsou rozhodovací stromy, protože mají
\emph{nízké zkreslení} a~\emph{vysoký rozptyl}, jsou-li dostatečně hluboké.
\textbf{Omezení:} bagging \emph{nesnižuje zkreslení} modelu. Je-li zkreslení hlavním zdrojem
chyby, může přesnost dokonce mírně zhoršit, a~to kvůli korelacím mezi komponentami, které
porušují předpoklad i.i.d.
\end{defbox}

\begin{defbox}[Náhodné lesy (\emph{random forests})]
Mají-li $k$ klasifikátorů s~rozptylem $\sigma^2$ vzájemnou \emph{pozitivní korelaci} $\rho$,
je rozptyl zprůměrované predikce
\[ \rho\sigma^2+\frac{(1-\rho)\sigma^2}{k}. \]
Člen $\rho\sigma^2$ je \emph{nezávislý} na počtu komponent $k$, a~právě on omezuje přínos
baggingu. Bootstrapované stromy přitom pozitivně korelované bývají.

\emph{Náhodný les} je ensemble rozhodovacích stromů, do jejichž \emph{konstrukce} je
explicitně vložena náhodnost, aby byly komponenty méně korelované. Vedou k~tomu dvě cesty.
\emph{(1)}~\textbf{Bagging} znamená, že se stromy trénují na mírně odlišných datech, totiž na
bootstrapových vzorcích. \emph{(2)}~\textbf{Omezený počet příznaků} znamená, že se stromu
v~každém uzlu předloží jen \emph{náhodná podmnožina} $m$ příznaků a~informační zisk se počítá
jen na ní. Druhá cesta navíc zrychluje trénování, protože je k~prohledání méně příznaků,
a~stromy není nutné prořezávat; snížený rozptyl přitom nezhoršuje zkreslení. Výstupem lesa je
většinový hlas stromů.

\textbf{Algoritmus} je pro každý z~$N$ stromů stejný: vytvoř nový bootstrapový vzorek
trénovací množiny a~nad ním trénuj rozhodovací strom tak, že v~každém uzlu vybereš náhodně
$m$ příznaků, spočítáš informační zisk jen na nich a~vybereš optimální; opakuj, dokud není
strom hotový.

\textbf{Parametry} jsou dva. Počet příznaků $m$ se běžně volí jako
$\sqrt{\text{počet příznaků}}$, na tento parametr ostatně nejsou náhodné lesy příliš citlivé.
S~počtem stromů lze pokračovat tak dlouho, dokud chyba klesá.
\end{defbox}

\begin{defbox}[Boosting a~AdaBoost (Freund \& Schapire, 1996)]
\textbf{Princip:} každému trénovacímu vzoru je přiřazena \emph{váha}, různé klasifikátory se
trénují s~těmito vahami a~váhy se iterativně upravují podle výkonu klasifikátoru. Budoucí
modely tedy \emph{závisí na předchozích}: každý klasifikátor se staví tímtéž algoritmem, ale
na jinak váženém trénovacím souboru. \textbf{Idea} je soustředit se v~dalších iteracích na
\emph{nesprávně klasifikované} vzory tím, že se zvýší jejich relativní váha. Iterativním
opakováním a~vytvořením vážené kombinace klasifikátorů tak lze získat klasifikátor s~nižším
celkovým \emph{zkreslením}, na rozdíl od baggingu, který snižuje rozptyl.

\textsc{AdaBoost}(data $D$, bázový klasifikátor $A$, max.\ počet kol $T$):
\begin{algorithmic}[1]
\State $t\gets 0$; pro každý vzor $i$ inicializuj $W_1(i)=1/n$
\Repeat
  \State $t\gets t+1$
  \State urči váženou chybu $\epsilon_t$ na $D$, aplikuje-li se $A$ na data s~vahami
         $W_t(\cdot)$
  \State $\alpha_t\gets \tfrac12\ln\big((1-\epsilon_t)/\epsilon_t\big)$
  \For{každý \emph{nesprávně} klasifikovaný $X_i\in D$}
     \State $W_{t+1}(i)\gets W_t(i)e^{\alpha_t}$
  \EndFor
  \For{každý \emph{správně} klasifikovaný $X_i\in D$}
     \State $W_{t+1}(i)\gets W_t(i)e^{-\alpha_t}$
  \EndFor
  \State normalizuj $W_{t+1}(i)\gets W_{t+1}(i)/\sum_j W_{t+1}(j)$
\Until{$t\ge T$ \textbf{nebo} $\epsilon_t=0$ \textbf{nebo} $\epsilon_t\ge 0{,}5$}
\State klasifikuj testovací vzory ensemblem s~vahami komponent $\alpha_t$
\end{algorithmic}
Zde $\epsilon_t$ je podíl trénovacích vzorů, které model v~$t$-té iteraci predikoval
nesprávně, a~to po vážení vahami $W_t$.

\textbf{Vlastnosti:} boosting produkuje \emph{posloupnost} klasifikátorů se stejným bázovým
učicím algoritmem, kde každý závisí na předchozím a~soustředí se na jeho chyby. Při testování
se výsledky celé posloupnosti kombinují. \textbf{Citlivost na šum} je jeho slabinou: je-li
odlehlých hodnot velmi mnoho, důraz na \uv{obtížné} příklady výkon spíš zhorší.
\end{defbox}

Náhodné lesy a~boosting se tedy hodí do různých situací:

\begin{center}
\small
\begin{tabular}{@{}p{7.6cm}p{8.0cm}@{}}
\toprule
\textbf{Náhodné lesy} & \textbf{Boosting} \\
\midrule
rychlé --- mohou běžet \emph{paralelně} a~v~každém uzlu prohledávají malou množinu příznaků
  & vyčerpávající --- v~každé fázi prohledává \emph{celou} množinu příznaků a~každá fáze
  závisí na předchozí \\
stromy jsou levné na trénování, lze jich ve stejném čase vypěstovat více i~na velkých
  a~složitých datech & musí běžet \emph{sekvenčně}, což může být nákladné \\
překvapivě dobré výsledky vzhledem k~malé části prostoru problému, kterou každý strom vidí
  & při \emph{stejném počtu} stromů boosting náhodné lesy překonává (lesy v~každé fázi
  prohledávají jen podmnožinu) \\
\bottomrule
\end{tabular}
\end{center}

\subsubsection{Srovnání klasifikačních metod}

Následující tabulka staví probrané metody vedle sebe podle srozumitelnosti výsledného modelu,
dosahované přesnosti a~rychlosti:

\begin{center}
\small
\begin{tabular}{@{}p{2.9cm}p{2.3cm}p{2.2cm}p{2.3cm}p{4.6cm}@{}}
\toprule
\textbf{Metoda} & \textbf{Srozumitel\-nost} & \textbf{Přesnost} & \textbf{Rychlost}
  & \textbf{Poznámka} \\
\midrule
Rozhodovací strom & vysoká & střední & vysoká & obdélníkové oblasti, nutné prořezávání,
  ignoruje korelace atributů \\
Pravidla z~C4.5 & nejvyšší & střední & vysoká & prořezávání pravidel je radikálnější než
  stromů \\
Naivní Bayes & střední & překvapivě dobrá & velmi vysoká & zvládne nekompletní vzory;
  problém s~nulovými pravděpodobnostmi \\
SVM & velmi nízká & vysoká & nízká (trénink) & nejlepší pro vysokorozměrná data (text);
  jen 2 třídy, jen reálný prostor \\
Logistická regrese & vysoká & střední & vysoká & pravděpodobnostní výstup,
  interpretovatelné koeficienty \\
Diskriminační analýza & střední & dobrá při normalitě & vysoká & stačí odhadnout
  $\mu_j$ a~$S_j$ \\
Neuronová síť / ELM & velmi nízká & vysoká & ELM velmi vysoká & ELM negeneruje lokální
  minima, netrénuje skrytou vrstvu \\
Náhodný les & nízká & velmi vysoká & vysoká (paralelně) & snižuje rozptyl, nezhoršuje
  zkreslení \\
Boosting / AdaBoost & nízká & velmi vysoká & nízká (sekvenčně) & snižuje zkreslení; citlivý
  na šum a~odlehlé hodnoty \\
\bottomrule
\end{tabular}
\end{center}

\subsection{Klastrová analýza}
\label{ssec:clustering}

\subsubsection{Formulace úlohy a~míry vzdálenosti}

Předchozí oddíly se zabývaly klasifikací, kde byla u~každého trénovacího vzoru známa jeho
třída. Shlukování žádné takové označení nemá a~strukturu si musí v~datech najít samo.

\emph{Klastrová analýza} (shlukování, \emph{cluster analysis}) rozděluje pozorované vzory do
skupin, tedy klastrů či shluků, které sdružují vzájemně podobné vzory. Aby to bylo možné,
musí platit jeden \textbf{předpoklad:} umíme mezi dvěma vzory změřit \emph{vzdálenost}.
Každý vzor je přitom charakterizován $m$ numerickými veličinami, takže vzdálenost se počítá
z~jejich hodnot.

\begin{defbox}[Míry vzdálenosti]
Pro vzory $x_1=(x_{11},\dots,x_{1m})$ a~$x_2=(x_{21},\dots,x_{2m})$ se používají zejména tyto
čtyři míry:
\begin{enumerate}
\item \textbf{Hammingova} vzdálenost (\emph{Manhattan}, $L_1$) sečte absolutní rozdíly
      v~jednotlivých souřadnicích:
      $d_H(x_1,x_2)=\sum_{i=1}^{m}|x_{1i}-x_{2i}|$
\item \textbf{Euklidovská} vzdálenost ($L_2$) je běžná přímá vzdálenost v~prostoru:
      $d_E(x_1,x_2)=\sqrt{\sum_{i=1}^{m}(x_{1i}-x_{2i})^2}$
\item \textbf{Čebyševova} vzdálenost ($L_\infty$) se dívá jen na největší rozdíl mezi
      souřadnicemi:
      $d_C(x_1,x_2)=\max_i |x_{1i}-x_{2i}|$
\item \textbf{Minkowského metrika} ($L_q$) všechny tři předchozí zastřešuje, protože jsou
      jejími speciálními případy:
      $d_q(x_1,x_2)=\Big(\sum_{i=1}^{m}|x_{1i}-x_{2i}|^{q}\Big)^{1/q}$
\end{enumerate}
Konkrétně platí $d_H=d_1$, $d_E=d_2$ a~$d_C=\lim_{q\to\infty}d_q$. Vezmeme-li vzor
$(0,\dots,0)$ a~vzor jednotkové vzdálenosti, dají všechny tři metriky shodně hodnotu 1.
\end{defbox}

Kterou metriku zvolit, závisí na zúčastněných veličinách a~na jejich rozsahu. Veličina
měřená ve~velkých číslech by jinak přehlušila ostatní, a~proto je nutné veličiny
\textbf{normalizovat}: dělením průměrem, směrodatnou odchylkou, rozsahem $\max-\min$
a~podobně (viz kap.~\ref{sec:priprava}). Uvedené metriky navíc mlčky předpokládají, že
všechny veličiny mají \emph{stejný rozptyl}. Mají-li rozptyl \emph{různý}, sáhne se po jiné
míře:

\begin{defbox}[Mahalanobisova vzdálenost]
\[ d_M^2(x_1,x_2)=(x_1-x_2)^{\!\top}S^{-1}(x_1-x_2), \]
kde $S$ je kovarianční matice. Tato míra zohledňuje jak různé rozptyly veličin, tak korelace
mezi nimi; pro $S=I$ se redukuje na kvadrát euklidovské vzdálenosti.
\end{defbox}

Vzdálenost mezi jednotlivými vzory ale nestačí. Shlukovací algoritmy potřebují umět změřit
i~vzdálenost dvou celých skupin, a~na to existuje několik různých konvencí.

\begin{defbox}[Vzdálenost mezi dvěma klastry $C_k$ a~$C_l$]
\textbf{Metoda nejbližšího souseda} (\emph{single linkage}) bere minimum ze vzdáleností
jejich prvků:
$D(C_k,C_l)=\min_{x_i,x_j} d(x_i,x_j)$, $x_i\in C_k$, $x_j\in C_l$.

\textbf{Metoda nejvzdálenějšího souseda} (\emph{complete linkage}) pracuje zrcadlově
a~vezme maximum:
$D(C_k,C_l)=\max_{x_i,x_j} d(x_i,x_j)$.

\textbf{Metoda průměrné vzdálenosti} (\emph{average linkage}) zprůměruje vzdálenosti mezi
všemi dvojicemi vzorů, přičemž $n_k$ a~$n_l$ jsou počty vzorů v~obou klastrech:
\[ D(C_k,C_l)=\frac{1}{n_kn_l}\sum_{i=1}^{n_k}\sum_{j=1}^{n_l} d(x_i,x_j). \]

\textbf{Centroidní metoda} měří vzdálenost středů klastrů, tedy $D(C_k,C_l)=d(c_k,c_l)$, kde
$c_k$ a~$c_l$ jsou středy (centroidy) obou klastrů.
\end{defbox}

\emph{Centroid} je střed klastru a~slouží jako jeho \uv{prototyp}, tedy jediný bod, který celý
klastr zastupuje. Tentýž klastr přitom může být reprezentován jedním centroidem, ale i~více
centroidy současně; záleží na tvaru klastru a~na zvolené metrice.

\subsubsection{Metoda $k$-means}

Nejrozšířenější shlukovací algoritmus vychází z~prosté úvahy: klastr je popsán svým
centroidem a~vzor patří tam, kde má centroid nejblíž. Jenže jedno závisí na druhém, protože
centroidy se počítají z~rozdělení vzorů a~rozdělení vzorů se řídí centroidy. Řešením je obě
části střídavě přepočítávat tak dlouho, dokud se rozdělení neustálí.

\begin{defbox}[Algoritmus $k$-means (Lloyd, 1982)]
\begin{algorithmic}[1]
\State rozděl vzory \emph{náhodně} do $k$ klastrů $C_1,\dots,C_k$ (disjunktních, ve~sjednocení
       všechny vzory)
\State \label{alg:km2} urči centroidy všech klastrů v~aktuálním rozdělení,
       např.\ $c_k=\frac{1}{|C_k|}\sum_{x\in C_k}x$
\For{každý vzor $x$}
   \State urči vzdálenosti $d(x,c_k)$ pro $1\le k\le K$
   \State nechť $d(x,c_l)=\min_k d(x,c_k)$
   \If{$x$ nepatří do klastru $l$} \State přesuň $x$ do klastru $l$ \EndIf
\EndFor
\If{došlo k~přesunu} \State goto krok~\ref{alg:km2} \Else \State \textbf{konec} \EndIf
\end{algorithmic}
Jakmile algoritmus skončí, jsou klastry reprezentovány svými centroidy.

\textbf{Varianty.} Místo náhodného počátečního rozdělení lze za centroidy rovnou prohlásit
\emph{prvních $k$ vzorů}, čímž v~první iteraci odpadá krok~\ref{alg:km2}. Druhá obměna
upravuje centroidy \emph{po každém} přesunu, tedy uvnitř cyklu, a~ne až po jeho skončení.
\end{defbox}

\textbf{Rozšíření Lloydova algoritmu.} Střídání minimalizačních kroků neposkytuje žádnou
záruku aproximace. Lloydův algoritmus ale cenu počátečního shlukování \emph{nikdy nezvýší},
takže výsledek stojí a~padá s~tím, odkud se vyjde. Vyplatí se proto dodat mu
\emph{prokazatelně dobrou inicializaci}:
\begin{itemize}
\item \textbf{$k$-means++} (Arthur, Vassilvitskii, 2007) vybírá $k$ centroidů z~$P$ vzorů
      postupně, a~to s~pravděpodobností úměrnou kvadrátu vzdálenosti k~nejbližšímu už
      existujícímu centroidu:
      $\dfrac{d^2(x_i,c)}{\sum_{j=1}^{P}d^2(x_j,c)}$. Vzdálené body tedy mají větší šanci
      stát se centroidem.
\item \textbf{LocalSearch++} (Lattanzi, Sohler, 2019) nejprve spustí $k$-means++ pro volbu $k$
      počátečních centroidů. Poté vzorkuje nové vzory s~analogickou pravděpodobností
      a~nahrazuje jimi existující centroidy, ovšem jen tehdy, pokud se tím sníží celková
      kvadratická vzdálenost vzorů k~jejich nejbližším centroidům.
\end{itemize}

Průměr není jediná možná volba reprezentanta. \textbf{Metoda $k$-medians} používá jako
účelovou funkci Manhattanovu vzdálenost, tedy $d(x_i,c)=\|x_i-c\|_1$. Lze ukázat, že
optimálním reprezentantem $c$ je pak \emph{medián} datových bodů v~klastru, počítaný
\emph{podél každé dimenze} zvlášť. Právě proto, že se medián volí v~každé dimenzi nezávisle,
výsledný $m$-rozměrný reprezentant obvykle \emph{nepatří} do původní datové množiny. Zisk je
ve~větší odolnosti: $k$-medians volí reprezentanty \emph{robustněji} než $k$-means, protože
medián není na odlehlé hodnoty tak citlivý jako průměr.

\subsubsection{Metoda $k$-medoids, CLARA a~CLARANS}

$k$-means i~$k$-medians reprezentanta \emph{dopočítávají}, takže výsledný bod v~datech ležet
nemusí. Následující rodina metod jde opačnou cestou a~za reprezentanty prohlašuje přímo
některé z~pozorovaných bodů.

\begin{defbox}[Algoritmus $k$-medoids]
Pro datovou množinu $D$ o~$n$ $d$-rozměrných bodech $x_1,\dots,x_n$ se hledá $k$
reprezentantů $y_1,\dots,y_k$ \emph{z~původní databáze} $D$, jejichž počet $k$ zadává
uživatel. Reprezentanti mají minimalizovat účelovou funkci danou součtem vzdáleností bodů
k~jejich nejbližším reprezentantům:
\[ O=\sum_{i=1}^{n}\min_{j} d(x_i,y_j). \]
Jedna iterace stojí $O(k\,n\,d)$ času a~obvykle stačí malý konstantní počet iterací.

\textbf{Proč volit reprezentanty z~$D$?} Důvody jsou dva. \emph{(1)}~Klastry, které najde
$k$-means, mohou být zkresleny odlehlými hodnotami: reprezentanti se pak ocitnou v~prázdných
oblastech, jež nejsou pro většinu bodů klastru reprezentativní, a~různé klastry mohou
částečně splynout. Částečně se to dá řešit pečlivým zacházením s~odlehlými hodnotami.
\emph{(2)}~U~složitých dat, například u~množiny časových řad různé délky, může být obtížné
optimálního centrálního reprezentanta vůbec spočítat. $k$-medoids lze proto definovat
prakticky pro \emph{libovolný typ dat}; stačí poskytnout vhodnou funkci podobnosti či
vzdálenosti.

\textbf{Postup} je \emph{hill-climbing}. Množina reprezentantů $S$ se inicializuje body
z~$D$ a~pak se iterativně zlepšuje tím, že se bod z~$S$ vymění za datový bod z~$D$. Které
výměny zkoušet, řeší dvě strategie:
\begin{enumerate}
\item Vyzkoušet všech $|S|\cdot|D|$ možností a~vybrat tu nejlepší. Je to \emph{extrémně
      drahé}.
\item Pracovat s~\emph{náhodně vybranou} množinou $r$ párů $(x,y)$, kde $x\in D$ a~$y\in S$,
      a~provést nejlepší z~nich. Tato varianta vyžaduje čas proporcionální $r\cdot|D|$.
      Algoritmus konverguje, když se účelová funkce už nezlepší, nebo když průměrné zlepšení
      klesne pod uživatelem zadaný práh.
\end{enumerate}
Ve srovnání s~$k$-means je $k$-medoids \emph{pomalejší}, zato ho lze použít na různé typy dat.
\end{defbox}

\begin{defbox}[Škálovatelné varianty: CLARA a~CLARANS]
V~mnoha aplikacích jsou data tak velká, že se nevejdou do hlavní paměti a~musí zůstat na
disku, což svazuje ruce při návrhu algoritmu. CLARA a~CLARANS jsou škálovatelné implementace
$k$-medoids a~dědí po bázovém algoritmu výhody i~nevýhody: CLARANS je sice snadno
generalizovatelný na různé typy dat, ale nese si i~jeho poměrně vysokou výpočetní složitost.

\textbf{CLARA} (\emph{Clustering LARge Applications}) vychází z~varianty $k$-medoids zvané
\emph{PAM} (\emph{Partitioning Around Medoids}). PAM testuje na výměnu \emph{všech} $k(n-k)$
párů medoid--nemedoid, vymění pár s~nejlepším zlepšením účelové funkce a~takto pokračuje až do
konvergence k~lokálnímu optimu. Každá iterace si vyžádá $O(kn^2)$ výpočtů vzdálenosti, tedy
$O(kn^2d)$ času, což je \emph{značně drahé}. CLARA proto celý algoritmus aplikuje jen na menší
\emph{vzorek} velikosti $f\cdot n$ (kde $f\ll 1$) a~zbylé nevzorkované body pouze přiřadí
k~nalezeným medoidům. Postup se zopakuje nad několika nezávisle volenými vzorky stejné
velikosti a~vybere se nejlepší nalezené shlukování. Složitost iterace tím klesne na
$O(kf^2n^2d+k(n-k))$, což je pro malá $f$ o~řády rychlejší. \textbf{Hlavní problém} zůstává:
předvybrané vzorky nemusí dobrou volbu medoidů vůbec obsahovat.

\textbf{CLARANS} (\emph{Clustering Large Applications based on RANdomized Search}) se tomuto
problému vyhýbá tím, že pracuje s~\emph{plnými} daty. Iterativně zkouší výměny \emph{náhodných}
medoidů za \emph{náhodné} nemedoidy a~po každém pokusu kontroluje, zda se kvalita zlepšila.
Když ano, výměna zůstane definitivně; když ne, přičte se neúspěšný pokus. Lokální optimum je
nalezeno, jakmile počet neúspěšných pokusů dosáhne uživatelem zadané meze \emph{MaxAttempt}.
Celý proces se opakuje \emph{MaxLocal}-krát a~vybere se nejlepší z~nalezených lokálních optim.
Díky tomu CLARANS prozkoumá \emph{větší rozmanitost} prostoru řešení než CLARA.
\end{defbox}

\subsubsection{Hierarchické shlukování}

Všechny dosavadní metody vyžadovaly, aby uživatel počet klastrů $k$ zadal dopředu. To je
nepříjemné, protože právě to obvykle nevíme. Hierarchické metody volbu odsouvají až na
konec: vybudují celou posloupnost postupně vnořených rozdělení a~teprve z~ní se vybere ta
úroveň, která dává smysl.

\begin{defbox}[Aglomerativní hierarchické shlukování]
Tento přístup jde \uv{zdola nahoru} (\emph{bottom-up}).

\textbf{Inicializace} má dva kroky: \emph{(1)}~urči vzájemné vzdálenosti mezi všemi vzory
a~\emph{(2)}~přiřaď každý vzor do samostatného klastru.

\textbf{Hlavní cyklus} běží, dokud je klastrů více než jeden. V~každém průchodu
\emph{(1)}~najdi dva vzájemně nejbližší klastry a~slouč je a~\emph{(2)}~pro nově vzniklý
klastr spočítej jeho vzdálenost od ostatních klastrů.

Výsledek se zobrazuje jako \textbf{dendrogram}, který zleva doprava ukazuje postupné slučování
klastrů; výška spojení odpovídá vzdálenosti, při níž ke sloučení došlo. \emph{Optimální počet
klastrů není znám předem} a~volí se až dodatečně, řezem dendrogramu ve~vhodné výšce.
\end{defbox}

Opačným směrem postupuje \textbf{divizivní} hierarchické shlukování (\uv{shora dolů}).
Začíná jediným klastrem, který obsahuje všechny vzory, a~ten postupně dělí. Nejznámějším
příkladem v~kontextu sítí je Girvanové--Newmanův algoritmus (kap.~\ref{sec:sna}).

\begin{defbox}[CURE (\emph{Clustering Using REpresentatives})]
Aglomerativní hierarchický algoritmus, který je \emph{velmi rychlý} a~umí odhalit klastry
\emph{libovolného tvaru}, což třeba CLARANS nedokáže. Postupuje v~pěti krocích:
\begin{enumerate}
\item Vyber vzorek $s$ bodů z~databáze $D$ o~$n$ bodech.
\item Rozděl těchto $s$ bodů na $p$ částí velikosti $s/p$.
\item Každou část shlukuj \emph{nezávisle} hierarchickým slučováním do $k'$ klastrů. Celkový
      počet $k'\cdot p$ klastrů přes všechny části je stále větší než uživatelem požadované
      $k$.
\item Nad těmito $k'\cdot p$ klastry proveď hierarchické shlukování až na požadovaných $k$.
\item Každý z~$(n-s)$ nevzorkovaných bodů přiřaď ke klastru s~nejbližším reprezentantem.
\end{enumerate}
\end{defbox}

\subsubsection{Vektorová kvantizace: algoritmus LVQ}

Předchozí metody si vystačily s~neoznačenými daty. LVQ je proti nim výjimkou: pracuje
se~vzory, u~nichž třídu známe, a~tuto informaci využívá k~posouvání reprezentantů.

\begin{defbox}[LVQ --- \emph{Learning Vector Quantization}]
Slouží pro \textbf{shlukování s~učitelem} (třídy jsou označeny $C$):
\begin{algorithmic}[1]
\State inicializuj všechny váhové vektory $w_j(0)$, rychlosti učení $\mu(0)$, $k\gets 0$
\While{není splněna zastavovací podmínka}
  \For{každý trénovací vzor $x$}
     \State urči index $j=q$ váhového vektoru s~\emph{minimální} euklidovskou vzdáleností,
            tj.\ $\min_j\|x-w_j\|_2^2$
     \If{třída $x$ souhlasí s~třídou $w_q$}
        \State $w_q(k+1)\gets w_q(k)+\mu(k)\big(x-w_q(k)\big)$ \Comment{přitáhni}
     \Else
        \State $w_q(k+1)\gets w_q(k)-\mu(k)\big(x-w_q(k)\big)$ \Comment{odtlač}
     \EndIf
  \EndFor
  \State $k\gets k+1$; sniž rychlost učení, např.\ $\mu(k)=\mu(k-1)/(k+1)$ pro $k>0$
\EndWhile
\end{algorithmic}
Váhové vektory se obvykle inicializují \emph{prvními $m$ vektory} trénovací množiny, které
musí zahrnovat vzory ze všech tříd. Neurony se tím automaticky \uv{kalibrují} s~třídami.
\end{defbox}

\subsubsection{Mřížkové a~hustotní metody}

Metody postavené na reprezentantech mají společné omezení: klastr je u~nich vždy okolí
jednoho bodu, takže vycházejí kulovité. Protáhlé či jinak nepravidelné tvary vyžadují změnu
úhlu pohledu. Neptáme se, jak daleko je bod od středu, ale kde v~prostoru je bodů hodně
a~kde málo.

\begin{defbox}[Mřížkové metody (např.\ MAFIA)]
\begin{itemize}
\item Data se nejdřív diskretizují do $p$ (ekvidistantních) intervalů, což pro $d$-rozměrná
      data dá $p^d$ hyperkrychlí.
\item Tyto hyperkrychle pak slouží jako stavební bloky pro shlukování.
\item Pomocí prahu hustoty $\tau$ se identifikují \emph{husté} hyperkrychle, tedy ty
      s~alespoň $\tau$ body.
\item Klastr libovolného tvaru se skládá z~více hustých oblastí, které jsou \emph{propojeny}.
      Dvě mřížkové oblasti přitom považujeme za propojené, sdílejí-li stěnu.
\end{itemize}
\textbf{Cílem} je určit klastry jako husté a~vzájemně propojené oblasti. Formálně: hustým
mřížkovým buňkám odpovídají uzly grafu a~hrany reprezentují sousedskou propojenost. Data
pokrytá buňkami z~jedné komponenty souvislosti tohoto grafu pak tvoří jeden výsledný klastr.
\end{defbox}

\begin{defbox}[DBSCAN --- \emph{Density-Based Spatial Clustering of Applications with Noise}]
Myšlenka je podobná, jen stavebními bloky nejsou mřížkové buňky, ale \emph{jednotlivé datové
body}, přesněji kruhové oblasti okolo nich. Hustota je dána počtem bodů ležících v~jejich
poloměru $\varepsilon$ (včetně bodu samotného) a~podle ní se body dělí na tři druhy:
\begin{enumerate}
\item \textbf{Jádrový bod} (\emph{core point}) je takový, jehož kulová oblast o~poloměru
      $\varepsilon$ obsahuje alespoň $\tau$ datových bodů.
\item \textbf{Hraniční bod} (\emph{border point}) má ve~své oblasti \emph{méně} než $\tau$
      bodů, ale zároveň v~ní leží alespoň jeden jádrový bod ve~vzdálenosti $\varepsilon$.
\item \textbf{Šumový bod} (\emph{noise point}) není ani jádrový, ani hraniční.
\end{enumerate}

\textbf{Algoritmus} má tři kroky:
\begin{enumerate}
\item Urči jádrové, hraniční a~šumové body.
\item Zkonstruuj graf propojenosti, v~němž každý uzel odpovídá jádrovému bodu. Hrana mezi
      dvěma jádrovými body existuje \emph{právě když} leží od sebe ve~vzdálenosti nejvýše
      $\varepsilon$.
\item Najdi všechny komponenty souvislosti tohoto grafu; ty odpovídají klastrům postaveným na
      jádrových bodech. Hraniční body se přiřadí ke klastru, s~nímž mají nejvyšší úroveň
      propojenosti, a~šumové body se ohlásí jako \emph{odlehlé hodnoty}.
\end{enumerate}
\end{defbox}

Probrané metody se dají srovnat podle několika mála otázek: jaký tvar klastrů zvládnou najít,
jestli je nutné jejich počet zadat předem a~jak se chovají vůči odlehlým hodnotám.

\begin{center}
\small
\begin{tabular}{@{}p{2.6cm}p{2.3cm}p{2.4cm}p{2.2cm}p{5.0cm}@{}}
\toprule
\textbf{Metoda} & \textbf{Tvar klastrů} & \textbf{Počet klastrů} & \textbf{Odlehlé hodnoty}
  & \textbf{Poznámka} \\
\midrule
$k$-means & kulovité & zadán & citlivá & rychlá; závislá na inicializaci $\to$
  $k$-means++ \\
$k$-medians & kulovité & zadán & robustnější & medián po dimenzích; reprezentant není
  z~dat \\
$k$-medoids & kulovité & zadán & robustní & reprezentanti \emph{z~dat}; libovolný typ dat;
  pomalejší \\
CLARA / CLARANS & kulovité & zadán & robustní & škálovatelné $k$-medoids (vzorkování /
  náhodné prohledávání) \\
Hierarchické & podle metriky spojování & určí se z~dendrogramu & citlivé & bez nutnosti
  zadat $k$ předem; $O(n^2)$--$O(n^3)$ \\
CURE & libovolný & zadán & robustní & vzorkování $+$ hierarchické slučování po částech \\
LVQ & kulovité & dán počtem neuronů & --- & shlukování \emph{s~učitelem} \\
Mřížkové (MAFIA) & libovolný & vyplyne & šum vypadne & citlivé na granularitu mřížky \\
DBSCAN & libovolný & vyplyne & \emph{detekuje je} & parametry $\varepsilon$, $\tau$;
  problém s~různou hustotou klastrů \\
\bottomrule
\end{tabular}
\end{center}

\subsection{Shrnutí ke~zkoušce}

\textbf{Asociační pravidla.}
\begin{itemize}
\item \textbf{Support} je podíl transakcí obsahujících $i$ i~$j$, kdežto \textbf{confidence}
      je podíl transakcí s~$i$ i~$j$ mezi transakcemi obsahujícími $i$; \textbf{lift} se
      spočítá jako $p(i\wedge j)/\big(p(i)p(j)\big)$. Je-li lift $<1$, je pravidlo horší než
      náhoda a~lepší může být negace jeho závěru. Typický úkol u~tabule je spočítat všechny
      tři veličiny z~malé tabulky transakcí.
\item \textbf{Apriori} stojí na \emph{uzavřenosti dolů}, tedy na tom, že každá podmnožina
      frekventované množiny je sama frekventovaná, a~prohledává \emph{do šířky}. Procedura
      \textsc{Apriori-Gen} se skládá z~kroku spojení, který kombinuje množiny shodné
      v~$k-2$ kategoriích, a~z~kroku prořezání, jenž vyřadí kandidáta, chybí-li některá jeho
      podkombinace délky $k-1$. Prostor pravidel má velikost $O(2^m)$ a~MBA produkuje
      \emph{nadměrné množství pravidel}.
\item \textbf{Taxonomie a~virtuální položky} řeší otázku, na jaké úrovni obecnosti položky
      brát: vzácné položky se posunou výše v~taxonomii, časté níže. Virtuální položky nesou
      informaci o~transakci samotné (den, čas, zákazník); u~obojího hrozí redundance
      pravidel.
\item \textbf{MS-Apriori} útočí na problém vzácných položek tím, že každá položka dostane
      vlastní $\mathrm{MIS}$; minsup celého pravidla je pak \emph{minimum} z~MIS jeho položek,
      k~čemuž se přidává omezení rozdílu supportů $\varphi$. Cenou za to je, že
      \emph{uzavřenost dolů přestává platit}, a~algoritmus se musí přizpůsobit: položky se
      setřídí vzestupně podle MIS, místo $F_1$ se pracuje s~množinou semen $L$ a~pro generování
      pravidel slouží \texttt{tailCount}. U~zkoušky se hodí umět uvést protipříklad
      k~uzavřenosti dolů.
\item \textbf{CAR} pracuje s~transakcemi označkovanými třídou a~hledá pravidla tvaru
      $X\to y$. Dolují se \emph{přímo v~jednom kroku} přes pravidlové položky
      $(\mathit{condset},y)$ a~různým třídám lze zadat různé minsupy, přičemž hodnota
      101~\% danou třídu fakticky zakáže.
\item \textbf{Sekvenční vzory}: sekvence je uspořádaný seznam množin položek, její
      \emph{velikost} je počet prvků a~\emph{délka} počet položek; podsekvence se definuje
      přes vnořené inkluze. Algoritmus \textbf{GSP} pracuje obdobně jako Apriori. Pozor na
      to, že $\langle\{3\}\{8\}\rangle$ a~$\langle\{3,8\}\rangle$ se navzájem
      \emph{neobsahují}.
\item \textbf{FP-Growth} si vystačí se dvěma průchody daty a~negeneruje žádné kandidáty.
      FP-strom sdílí \emph{prefixy}, k~čemuž slouží heuristické uspořádání položek podle
      klesajícího supportu, a~doplňují ho spojové seznamy položek. Extrakce jde zdola nahoru
      metodou \emph{rozděl a~panuj} přes \emph{prefixové podstromy} a~\emph{podmíněné
      FP-stromy}: aktualizovat počítadla, odstranit uzly položky a~odstranit nefrekventované
      položky. V~nejhorším případě, kdy jsou všechny transakce unikátní, vyjde strom větší než
      samotná data.
\end{itemize}

\textbf{Učení s~učitelem.}
\begin{itemize}
\item \textbf{Rozhodovací stromy} se budují metodou TDIDT, tedy rozděl a~panuj, a~atribut se
      volí podle \emph{minimální vážené entropie}
      $H(A)=\sum_v\frac{|A(v)|}{|D|}H(A(v))$, což je ekvivalentní maximálnímu
      \emph{informačnímu zisku}. Očekává se schopnost projít příklad s~úvěrem (příjem 0,5747;
      konto 0,6667; pohlaví 0,9183; nezaměstnaný 0,825) a~převést hotový strom na pravidla.
      Indukce stojí $O(np\log p)$, klasifikace jednoho vzoru $O(q\log p)$.
\item \textbf{ID3} (Quinlan 1986) je hladový a~řídí se informačním ziskem. \textbf{C4.5}
      (1993) přidává práci s~chybějícími a~spojitými daty, prořezávání validací
      (\emph{reduced-error pruning}, \emph{subtree raising}), \emph{rule post-pruning}, odhad
      přesnosti dolní mezí $\mathit{acc}-1{,}96\sqrt{\mathit{acc}(1-\mathit{acc})/n}$
      a~\emph{poměrný} informační zisk, který penalizuje nadměrné větvení. \textbf{C5.0} je
      komerční verze s~boostingem. \textbf{CART} staví \emph{binární} stromy, používá Gini
      index $G(v_i)=1-\sum_j p_j^2$ (nižší je lepší) a~kritérium
      $\Phi=2P_LP_R\sum_j|P(C_j|t_L)-P(C_j|t_R)|$. \textbf{CHAID} se opírá o~$\chi^2$
      a~automaticky seskupuje hodnoty kategoriálních atributů.
\item \textbf{Bayesovská klasifikace} hledá $H_{MAP}=\arg\max_H P(E\mid H)P(H)$, přičemž
      jmenovatel lze zanedbat. \emph{Naivní} předpoklad podmíněné nezávislosti vede na
      $\arg\max_{H_t}P(H_t)\prod_i P(E_i\mid H_t)$. Metoda si poradí s~nekompletními vzory,
      naráží ale na problém nulových pravděpodobností. Příklad je potřeba umět spočítat
      (0,1042 vs.\ 0,0416).
\item \textbf{SVM} maximalizuje okraj $2/\|w\|$, což vede na primární úlohu
      $\min\frac{\langle w\cdot w\rangle}{2}$ za podmínek
      $y_i(\langle w\cdot x_i\rangle+b)\ge 1$. Kuhn--Tuckerovy podmínky jsou pro konvexní
      úlohu s~lineárními omezeními nejen nutné, ale \emph{i~postačující}; podmínka
      $\alpha_i>0$ platí jen pro \emph{podpůrné vektory} a~úloha se řeší přes Wolfeho duál.
      \textbf{Kernelový trik} využívá toho, že duál potřebuje jen skalární součiny, takže
      stačí zvolit $K(x,z)=\langle\phi(x)\cdot\phi(z)\rangle$ a~samotné $\phi$ nikdy
      nevyčíslíme (Mercerova věta). Omezení: jen 2 třídy, jen reálný prostor
      a~neinterpretovatelný výsledek.
\item \textbf{Logistická regrese} modeluje $\ln\frac{p}{1-p}=\beta^{\!\top}x$, což dává
      $p=\sigma(z)=1/(1+e^{-z})$. Učí se maximální věrohodností, ekvivalentně minimalizací
      křížové entropie, a~protože $\sigma'=\sigma(1-\sigma)$, vychází adaptační pravidlo
      $\Delta\beta_j=\eta\sum_i(y_i-\sigma(z_i))x_{ij}$.
\item \textbf{Diskriminační analýza} zavádí pro každou třídu funkci $g_j(x)$ a~klasifikuje se
      podle jejího maxima; pro dvě třídy stačí znaménko rozdílu $g=g_1-g_2$. Optimální volbou
      je $g_j\propto P(C_j\mid x)$. Pro normální rozdělení vyjde funkce kvadratická,
      a~jsou-li kovarianční matice shodné, dokonce \emph{lineární} --- pak stačí odhadnout
      $\mu_j$ a~$S$.
\item \textbf{ELM} volí parametry skrytých neuronů náhodně, dokonce i~bez znalosti dat,
      a~výstupní váhy dopočítá \emph{jedinou} pseudoinverzí $\beta=H^{+}T$; regularizace se
      přidá jako $I/C$ na diagonálu. Je proto rychlé, nemá lokální minima a~funguje
      s~libovolnou omezenou po částech spojitou přenosovou funkcí.
\item \textbf{Ensembly}: \emph{bagging} snižuje \emph{rozptyl}, pro i.i.d.\ komponenty
      z~$\sigma^2$ na $\sigma^2/k$, přičemž bootstrapový vzorek obsahuje přibližně
      $63{,}2\,\%$ různých bodů, tedy $1-1/e$. Ideální komponentou jsou stromy s~nízkým
      zkreslením a~vysokým rozptylem; zkreslení bagging nesnižuje. \emph{Náhodné lesy} navíc
      řeší korelaci komponent, protože rozptyl je $\rho\sigma^2+(1-\rho)\sigma^2/k$ a~první
      člen na $k$ vůbec nezávisí; náhodnost se vnáší jednak bootstrapem, jednak náhodnou
      podmnožinou $m\approx\sqrt{\#\text{příznaků}}$ příznaků v~každém uzlu.
      \emph{Boosting} a~konkrétně \textbf{AdaBoost} naopak snižuje \emph{zkreslení}: váhy
      chybně klasifikovaných vzorů rostou, hlas komponenty je
      $\alpha_t=\frac12\ln\frac{1-\epsilon_t}{\epsilon_t}$, běh je \emph{sekvenční}
      a~metoda je citlivá na šum. U~zkoušky se často ptá na srovnání lesů a~boostingu:
      lesy jsou paralelní a~levné, boosting vyčerpávající a~při stejném počtu stromů
      přesnější.
\end{itemize}

\textbf{Klastrová analýza.}
\begin{itemize}
\item \textbf{Metriky}, které je třeba znát: Hammingova neboli Manhattan ($L_1$),
      euklidovská ($L_2$), Čebyševova ($L_\infty$) a~Minkowského ($L_q$), jejímiž
      speciálními případy ostatní jsou. Pro veličiny s~různým rozptylem slouží
      \textbf{Mahalanobisova} vzdálenost $d_M^2=(x_1-x_2)^{\!\top}S^{-1}(x_1-x_2)$. Před
      shlukováním je vždy potřeba data \emph{normalizovat}.
\item \textbf{Vzdálenost dvou klastrů} se měří čtyřmi způsoby: metodou nejbližšího souseda
      (minimum), nejvzdálenějšího souseda (maximum), průměrné vzdálenosti a~centroidní
      metodou.
\item \textbf{$k$-means} (Lloyd 1982) vyjde z~náhodného rozdělení, spočítá centroidy, přesune
      vzory k~nejbližšímu centroidu a~opakuje, dokud k~přesunům dochází. Cenu shlukování nikdy
      nezvýší, ale nedává záruku aproximace, a~proto se doplňuje o~lepší inicializaci:
      \textbf{$k$-means++} volí centroid s~pravděpodobností úměrnou kvadrátu vzdálenosti
      k~nejbližšímu už zvolenému centroidu, \textbf{LocalSearch++} na to navazuje výměnami.
\item \textbf{$k$-medians} používá Manhattanovu vzdálenost a~medián po dimenzích, takže je
      robustnější, zatímco \textbf{$k$-medoids} volí reprezentanty \emph{z~dat}, zlepšuje je
      metodou hill-climbing s~výměnami, stojí $O(knd)$ na iteraci a~zvládne libovolný typ dat.
      \textbf{PAM} testuje všech $k(n-k)$ výměn za cenu $O(kn^2d)$; \textbf{CLARA} pracuje se
      vzorkem, čímž se dostane na $O(kf^2n^2d+k(n-k))$, ovšem s~rizikem špatně zvoleného
      vzorku, a~\textbf{CLARANS} prohledává náhodně nad plnými daty s~parametry
      \emph{MaxAttempt} a~\emph{MaxLocal}.
\item \textbf{Hierarchické shlukování} zdola nahoru začíná s~každým vzorem ve~vlastním
      klastru a~postupně slučuje dva nejbližší; výstupem je \textbf{dendrogram} a~počet
      klastrů se volí až jeho řezem. \textbf{CURE} vezme vzorek, rozdělí ho na $p$ částí,
      každou nezávisle shlukuje hierarchicky na $k'$ klastrů, výsledné klastry hierarchicky
      sloučí na $k$ a~nakonec přiřadí zbytek bodů; je rychlé a~zvládne klastry libovolného
      tvaru.
\item \textbf{LVQ} je shlukování \emph{s~učitelem}: vítězný váhový vektor se ke vzoru
      přitáhne, souhlasí-li třídy, a~odtlačí, když nesouhlasí. Rychlost učení $\mu$ přitom
      klesá jako $\mu(k-1)/(k+1)$.
\item \textbf{Mřížkové metody} (MAFIA) rozdělí prostor na $p^d$ hyperkrychlí, prahem hustoty
      $\tau$ vyberou husté buňky a~klastry sestaví jako propojené husté oblasti, formálně jako
      komponenty souvislosti grafu buněk. \textbf{DBSCAN} dělí body na \emph{jádrové}
      (alespoň $\tau$ bodů v~okolí o~poloměru $\varepsilon$), \emph{hraniční} (méně než $\tau$, ale
      jádrový bod do vzdálenosti $\varepsilon$) a~\emph{šumové}; klastry jsou komponenty
      souvislosti grafu jádrových bodů a~šumové body se hlásí jako odlehlé. Zvládne tedy
      klastry \emph{libovolného tvaru} a~odlehlé hodnoty detekuje sám.
\end{itemize}

%=====================================================================
\newpage

\section{Charakteristická a~diskriminační pravidla a~jejich zajímavost}
\label{sec:pravidla}
%=====================================================================

\mimo{} Celá kapitola stojí nad rámec slidů obou kurzů. Okruh ji ale jmenuje výslovně a~jde
o~standardní látku, která je zde zpracovaná podle Han--Kamber. Jak generalizace souvisí
s~OLAP, ukazuje kap.~\ref{sec:paradigmata}.

\subsection{Charakterizace vs.\ diskriminace}

Deskriptivní dobývání znalostí odpovídá na dvě otázky, které spolu těsně souvisejí, a~přesto
se nesmějí zaměňovat. První se ptá, jak vybraná skupina objektů vypadá; druhá, čím se tato
skupina liší od zbytku dat.

\begin{defbox}[Charakterizace vs.\ diskriminace]
\textbf{Charakterizace dat} (\emph{data characterization}) shrnuje obecné vlastnosti
\emph{cílové třídy} (\emph{target class}). Odpovídá tedy na otázku \uv{Jak vypadá typický
klient s~hypotékou?} a~jejím výsledkem jsou \emph{charakteristická pravidla}
(\emph{characteristic rules}).

\textbf{Diskriminace dat} (\emph{data discrimination}) porovnává obecné vlastnosti cílové
třídy s~vlastnostmi jedné či více \emph{kontrastních tříd} (\emph{contrasting classes}).
Odpovídá tedy na otázku \uv{Čím se klient s~hypotékou liší od klienta bez hypotéky?} a~jejím
výsledkem jsou \emph{diskriminační pravidla} (\emph{discriminant rules}).

Rozdíl mezi oběma úlohami je zásadní. Vlastnost může být pro cílovou třídu \emph{typická},
tedy může ji mít vysoký podíl objektů uvnitř třídy, a~přesto nemá žádnou \emph{rozlišovací}
sílu, protože je stejně typická i~pro kontrastní třídu.
\end{defbox}

Obě úlohy řeší jedna a~tatáž metoda, \emph{indukce orientovaná na atributy}
(\emph{attribute-oriented induction}, AOI), kterou navrhli Han, Cai a~Cercone (1991).

Alternativou k~ní je datově-kostkový (\emph{data cube}) přístup známý z~OLAP, ve kterém
generalizaci odpovídá operace \emph{roll-up} a~specializaci operace \emph{drill-down}. Každý
z~obou přístupů má svou přednost: AOI je flexibilnější, protože si poradí i~s~atributy, pro
které není kostka předem materializovaná, kdežto OLAP je rychlejší, protože sahá po
předpočítaných agregacích.

\subsection{Koncepční hierarchie}

Aby se dala hodnota atributu nahradit obecnější, musí být napřed řečeno, co je vůči čemu
obecnější. Tuto informaci nese koncepční hierarchie.

\begin{defbox}[Koncepční hierarchie]
\emph{Koncepční hierarchie} (\emph{concept hierarchy}) atributu $A$ je částečné uspořádání
jeho hodnot od nejkonkrétnějších (listy) po nejobecnější (kořen \texttt{ANY}). Podle toho,
odkud se uspořádání bere, se rozlišují čtyři typy:
\begin{itemize}
\item \emph{schématovou} (\emph{schema hierarchy}) dává přímo struktura dat: ulice $\prec$
      město $\prec$ kraj $\prec$ stát;
\item \emph{seskupovací} (\emph{set-grouping}) vzniká tím, že se hodnoty explicitně seskupí,
      třeba \{20--29\} $\prec$ \emph{mladý};
\item \emph{operační} (\emph{operation-derived}) se odvodí výpočtem z~hodnoty samotné:
      e-mailová adresa $\to$ doména;
\item \emph{pravidlová} (\emph{rule-based}) je definovaná pravidly, která smějí sahat přes
      více atributů najednou.
\end{itemize}
Hierarchii buď zadá expert, nebo se generuje automaticky: u~numerických atributů
diskretizací, u~kategoriálních podle počtu různých hodnot, protože atribut s~více různými
hodnotami leží obvykle na nižší úrovni hierarchie.
\end{defbox}

\subsection{Indukce orientovaná na atributy (AOI)}

Myšlenka AOI je jednoduchá. Dokud je tabulka příliš podrobná, nahrazují se hodnoty
obecnějšími a~řádky, které se tím staly shodnými, se slučují do jednoho.

\begin{defbox}[Algoritmus AOI]
\textbf{Vstup:} databáze, dotaz vymezující cílovou třídu, seznam atributů, prahy generalizace
$T_i$ pro jednotlivé atributy a~celkový práh $T$ pro velikost výsledné relace.

\begin{algorithmic}[1]
\State \textbf{sběr dat:} vyber dotazem množinu objektů cílové třídy (\emph{initial working
       relation})
\State ke každému objektu přidej atribut \texttt{count} $=1$
\For{každý atribut $A_i$}
   \If{$A_i$ má mnoho různých hodnot a~\emph{neexistuje} pro něj koncepční hierarchie}
      \State \textbf{odstranění atributu} (\emph{attribute removal})
   \ElsIf{$A_i$ má mnoho různých hodnot, ale hierarchie existuje}
      \State \textbf{generalizace atributu} (\emph{attribute generalization}): nahraď hodnoty
      \State jejich předky o~úroveň výš, dokud počet různých hodnot $\le T_i$
   \EndIf
\EndFor
\State \textbf{agregace:} slouč identické zobecněné objekty a~sečti jejich \texttt{count}
\While{počet řádků výsledné relace $>T$}
   \State zobecni dále atribut, jehož generalizace nejvíce sníží počet řádků
\EndWhile
\State \textbf{prezentace:} zobraz \emph{primární relaci} (\emph{prime relation}) jako
       tabulku, křížovou
\State tabulku, graf, nebo ji převeď na kvantitativní pravidla
\end{algorithmic}
\end{defbox}

Veškerou práci přitom odvedou \textbf{dvě základní generalizační operace:}
\begin{itemize}
\item \textbf{Odstranění atributu} přichází ke~slovu tehdy, má-li atribut příliš mnoho
      různých hodnot a~zároveň buď (a)~pro něj neexistuje generalizační operátor, nebo
      (b)~jeho vyšší úrovně jsou už vyjádřeny jinými atributy. Takový atribut se z~relace
      jednoduše odstraní; typicky jde o~\emph{jméno}, \emph{telefon} nebo \emph{rodné číslo}.
\item \textbf{Generalizace atributu} se použije tam, kde generalizační operátor existuje,
      tedy kde je po ruce koncepční hierarchie. Hodnoty se nahradí hodnotami na vyšší úrovni,
      jak ukazuje řetězec \emph{datum narození} $\to$ \emph{věk} $\to$ \emph{věková
      kategorie} nebo \emph{město} $\to$ \emph{kraj} $\to$ \emph{region}.
\end{itemize}

Kam až generalizace zajde, \textbf{řídí dva prahy}. \emph{Prah generalizace atributu} $T_i$
udává maximální povolený počet různých hodnot jednoho atributu a~volí se typicky mezi 2 a~8;
\emph{prah velikosti relace} $T$ omezuje počet řádků výsledné primární relace, obvykle na 10
až 30. Obě meze je nutné trefit. Příliš vysoká generalizace vede k~triviálním výsledkům typu
\uv{všichni klienti jsou lidé}, příliš nízká skončí u~tabulky, ve které se nikdo nevyzná.
Vhodná úroveň se proto hledá interaktivně, střídáním operací drill-down a~roll-up.

Cena za to je nízká. Realizuje-li se generalizace přes hašovací tabulku zobecněných záznamů,
stačí jediný průchod daty a~agregace, takže složitost činí $O(n\,a)$ pro $n$ objektů a~$a$
atributů. AOI je tedy \emph{lineární} v~počtu objektů, což je pro velké databáze zásadní.

\subsection{Míry \emph{t}-weight a~\emph{d}-weight}

Primární relace sama o~sobě říká jen to, jaké profily se v~datech vyskytují. Aby z~ní šlo
udělat pravidlo, potřebuje každý zobecněný záznam ještě číslo, které vyjádří jeho váhu.
Taková čísla jsou dvě a~každé míří na jiný typ pravidla.

\begin{defbox}[t-weight a~d-weight]
Nechť primární relace obsahuje zobecněné záznamy (\emph{generalized tuples}) $q_1,\dots,q_N$.

\textbf{Typičnost} (\emph{typicality}, t-weight) záznamu $q_a$ v~cílové třídě se definuje jako
\[ t\text{-}\mathrm{weight}(q_a)=\frac{\mathrm{count}(q_a)}{\sum_{i=1}^{N}\mathrm{count}(q_i)}
   \in[0,1], \]
tj.\ jako podíl objektů cílové třídy, které záznam $q_a$ pokrývá. Přes všechny záznamy
primární relace dávají t-weight dohromady 1, tedy 100~\%. Tato míra slouží pro
\emph{charakteristická} pravidla.

\textbf{Diskriminační váha} (d-weight) záznamu $q_a$ se měří vůči cílové třídě $C_{target}$
a~kontrastním třídám $C_1,\dots,C_m$:
\[ d\text{-}\mathrm{weight}(q_a)=\frac{\mathrm{count}(q_a\in C_{target})}
   {\sum_{j=1}^{m}\mathrm{count}(q_a\in C_j)+\mathrm{count}(q_a\in C_{target})}\in[0,1], \]
tj.\ jako podíl objektů pokrytých záznamem $q_a$, které patří do cílové třídy. Tato míra
slouží pro \emph{diskriminační} pravidla. Vysoká d-weight ($\to 1$) znamená, že záznam cílovou
třídu silně odlišuje, kdežto hodnota blízká apriornímu podílu cílové třídy znamená nulovou
rozlišovací schopnost.
\end{defbox}

\begin{defbox}[Kvantitativní charakteristické, diskriminační a~popisné pravidlo]
\textbf{Kvantitativní charakteristické pravidlo} vyjadřuje nutnou podmínku, tedy implikaci
\emph{z} třídy:
\[ \forall X\ \ \mathrm{target\_class}(X)\Rightarrow \mathrm{cond}_1(X)\,[t{:}\,w_1]\ \vee\
   \cdots\ \vee\ \mathrm{cond}_N(X)\,[t{:}\,w_N] \]
Čte se: patří-li $X$ do cílové třídy, pak splňuje $\mathrm{cond}_1$ s~typičností $w_1$, nebo
$\mathrm{cond}_2$ s~typičností $w_2$, \dots

\textbf{Kvantitativní diskriminační pravidlo} vyjadřuje naopak postačující podmínku, tedy
implikaci \emph{do} třídy:
\[ \forall X\ \ \mathrm{target\_class}(X)\Leftarrow \mathrm{cond}(X)\,[d{:}\,w] \]
Čte se: splňuje-li $X$ podmínku $\mathrm{cond}$, patří s~pravděpodobností $w$ do cílové třídy.

\textbf{Kvantitativní popisné pravidlo} (\emph{quantitative description rule}) obě předchozí
podoby kombinuje, takže podmínka je zároveň nutná i~postačující:
\[ \forall X\ \ \mathrm{target\_class}(X)\Leftrightarrow \mathrm{cond}_1(X)\,
   [t{:}\,w_1,\ d{:}\,w_1']\ \vee\ \cdots\ \vee\ \mathrm{cond}_N(X)\,[t{:}\,w_N,\ d{:}\,w_N'] \]
\end{defbox}

\begin{defbox}[Příklad --- charakterizace a~diskriminace klientů banky]
\textbf{Cílová třída:} klienti s~hypotékou (200 osob). \textbf{Kontrastní třída:} klienti bez
hypotéky (800 osob).

Nejprve proběhne AOI: odstraní se \emph{jméno}, \emph{rodné číslo} a~\emph{telefon},
generalizuje se \emph{datum narození} $\to$ \emph{věková kategorie}, \emph{plat} $\to$
\emph{příjmová kategorie} a~\emph{obec} $\to$ \emph{region}. Vznikne tato primární relace:

\begin{center}
\small
\begin{tabular}{@{}lllrrrr@{}}
\toprule
Věk & Příjem & Region & \multicolumn{2}{c}{počet} & $t$-weight & $d$-weight \\
 & & & cíl & kontrast & (cíl) & \\
\midrule
30--40 & vysoký & Praha & 90 & 10 & \textbf{45~\%} & \textbf{90~\%} \\
30--40 & střední & Praha & 50 & 150 & 25~\% & 25~\% \\
40--50 & vysoký & mimo Prahu & 40 & 160 & 20~\% & 20~\% \\
20--30 & nízký & mimo Prahu & 20 & 480 & 10~\% & 4~\% \\
\midrule
\multicolumn{3}{@{}l}{Celkem} & 200 & 800 & 100~\% & \\
\bottomrule
\end{tabular}
\end{center}

\textbf{Výpočet.} U~prvního záznamu vyjde $t\text{-}\mathrm{weight}=90/200=45\,\%$
a~$d\text{-}\mathrm{weight}=90/(90+10)=90\,\%$. Pro druhý záznam je $t=50/200=25\,\%$
a~$d=50/(50+150)=25\,\%$, pro čtvrtý $t=20/200=10\,\%$ a~$d=20/(20+480)=4\,\%$.

\textbf{Kvantitativní charakteristické pravidlo:}
\begin{align*}
\forall X\ \ \text{hypotéka}(X)\Rightarrow\ & \big(\text{věk}(X){=}\text{30--40}\wedge
  \text{příjem}(X){=}\text{vysoký}\wedge\text{region}(X){=}\text{Praha}\big)\ [t{:}\,45\,\%]\\
\vee\ & \big(\text{věk}(X){=}\text{30--40}\wedge\text{příjem}(X){=}\text{střední}\wedge
  \text{region}(X){=}\text{Praha}\big)\ [t{:}\,25\,\%]\ \vee\ \cdots
\end{align*}

\textbf{Kvantitativní diskriminační pravidlo:}
\[ \forall X\ \ \text{hypotéka}(X)\Leftarrow \big(\text{věk}(X){=}\text{30--40}\wedge
   \text{příjem}(X){=}\text{vysoký}\wedge\text{region}(X){=}\text{Praha}\big)\ [d{:}\,90\,\%] \]

\textbf{Interpretace rozdílu.} Právě poslední řádek tabulky ukazuje, proč jsou obě míry
potřeba. Profil \uv{20--30, nízký příjem, mimo Prahu} pokrývá 10~\% klientů s~hypotékou, což
není zanedbatelná typičnost, jeho d-weight je ale pouhá 4~\%, tedy výrazně \emph{méně} než
apriorní podíl cílové třídy $200/1000=20\,\%$. Tento profil hypotéku spíše \emph{vylučuje}.
U~druhého řádku leží d-weight jen nepatrně \emph{nad} apriorním podílem (25~\% proti 20~\%),
takže rozlišovací síla je zanedbatelná. A~u~třetího řádku je d-weight \emph{přesně rovna}
apriornímu podílu (20~\%), takže rozlišovací síla je \textbf{nulová}: profil o~příslušnosti
k~cílové třídě neříká vůbec nic, přestože pokrývá pětinu klientů s~hypotékou.
\end{defbox}

\subsection{Měření zajímavosti pravidel}

Když už jsou pravidla na světě, zbývá rozhodnout, která z~nich stojí za pozornost. Zajímavost
odvozených pravidel se posuzuje obdobně jako u~asociačních pravidel (odd.~\ref{ssec:ar}),
s~důrazem na následující míry:

\begin{center}
\small
\begin{tabular}{@{}lp{11.0cm}@{}}
\toprule
\textbf{Míra} & \textbf{Význam} \\
\midrule
$t$-weight & typičnost profilu uvnitř cílové třídy; pravidlo s~nízkou $t$-weight popisuje
  okrajový jev \\
$d$-weight & rozlišovací síla profilu; \emph{zajímavá} jsou pravidla s~$d$-weight výrazně
  vyšší (nebo nižší) než apriorní podíl třídy \\
Lift / zdvih & $\frac{d\text{-weight}}{P(C_{target})}$ --- kolikrát profil zvýší šanci na
  cílovou třídu \\
Pokrytí (\emph{coverage}) & podíl všech objektů pokrytých pravidlem; nízké pokrytí
  $\Rightarrow$ statisticky nespolehlivé \\
$\chi^2$ & statistická významnost vztahu profilu a~třídy \\
Informační zisk & snížení entropie třídy po znalosti profilu \\
Jednoduchost & délka pravidla (počet podmínek), počet pravidel; Occamova břitva \\
Neočekávanost & rozpor s~modelem světa uživatele (subjektivní míra) \\
Akceschopnost & lze na jeho základě jednat (subjektivní míra) \\
\bottomrule
\end{tabular}
\end{center}

\textbf{Prakticky} je pravidlo zajímavé tehdy, splňuje-li současně čtyři požadavky: musí být
(a)~\emph{srozumitelné}, (b)~\emph{platné} na nových datech s~vysokou spolehlivostí,
(c)~\emph{potenciálně užitečné} a~(d)~\emph{nové}. První tři se dají měřit objektivně,
poslední vyžaduje znalost uživatele.

Z~toho plyne \textbf{klasická past}. Pravidlo s~vysokou $t$-weight, ale s~$d$-weight rovnou
apriornímu podílu třídy není žádným objevem, i~když na první pohled vypadá impozantně. Věta
\uv{80~\% klientů s~hypotékou má bankovní účet} zní dobře jen do chvíle, než si člověk
uvědomí, že účet má i~80~\% klientů bez hypotéky.

\subsection{Shrnutí ke~zkoušce}

\begin{itemize}
\item \textbf{Charakterizace} popisuje cílovou třídu, kdežto \textbf{diskriminace} ji
      porovnává s~kontrastními třídami. Celý rozdíl se dá shrnout jedním vztahem: typičnost
      $\neq$ rozlišovací síla.
\item \textbf{AOI} postupuje tak, že dotazem sebere data a~ke~každému objektu přidá
      \texttt{count}. Pak každý atribut buď \emph{odstraní} (má-li mnoho hodnot a~neexistuje
      pro něj hierarchie), nebo \emph{generalizuje} podle koncepční hierarchie. Následně
      agreguje stejné záznamy a~generalizuje dál, dokud relace není menší než práh $T$;
      nakonec výsledek prezentuje. Složitost $O(na)$ je lineární v~počtu objektů.
\item \textbf{Koncepční hierarchie} se dělí na schématovou, seskupovací, operační
      a~pravidlovou.
\item \textbf{t-weight} $=\frac{\mathrm{count}(q_a)}{\sum_i\mathrm{count}(q_i)}$ měří
      typičnost uvnitř cílové třídy a~její součet přes všechny záznamy dává 100~\%. Používá se
      pro \emph{charakteristická} pravidla, tedy pro nutnou podmínku ($\Rightarrow$).
\item \textbf{d-weight} $=\frac{\mathrm{count}(q_a\in C_{target})}{\sum_j\mathrm{count}
      (q_a\in C_j)}$ měří rozlišovací sílu a~používá se pro \emph{diskriminační} pravidla,
      tedy pro postačující podmínku ($\Leftarrow$). Popisné pravidlo obě podoby kombinuje
      ($\Leftrightarrow$).
\item Obě váhy je potřeba umět spočítat z~tabulky a~vysvětlit, proč vysoká t-weight neznamená
      vysokou d-weight; klíčem je srovnání s~\emph{apriorním} podílem třídy.
\item \textbf{Zajímavost} má stránku objektivní (t/d-weight, lift, pokrytí, $\chi^2$,
      informační zisk, jednoduchost) a~subjektivní (neočekávanost, akceschopnost).
\end{itemize}

%=====================================================================
\newpage

\section{Reprezentace, vyhodnocování a~vizualizace znalostí}
\label{sec:reprez}
%=====================================================================

Věta okruhu spojuje tři samostatné otázky a~podle nich je členěná i~tato kapitola. První se ptá,
v~jaké formě nalezenou znalost vůbec vyjádříme. Druhá chce vědět, jak ověříme, že je ta znalost
dobrá. Třetí se týká toho, jak ji zobrazíme. Prostřední část, tedy vyhodnocování klasifikátorů,
je jediná, kterou slidy pokrývají v~celém rozsahu, a~je zároveň ta, na kterou se u~zkoušky ptá
nejčastěji.

\subsection{Formy reprezentace znalostí}

Znalost vydolovaná z~dat může mít velmi různou podobu a~volba formy rozhoduje o~tom, kdo jí
nakonec porozumí. Následující přehled shrnuje obvyklé formy reprezentace, metody, z~nichž
vznikají, a~cenu, kterou se za jejich srozumitelnost platí.

\mimoq{tabulka jako přehled; jednotlivé metody v~ní jsou ve slidech doloženy}
\begin{center}
\small
\begin{tabular}{@{}p{3.4cm}p{4.4cm}p{7.0cm}@{}}
\toprule
\textbf{Forma} & \textbf{Typický zdroj} & \textbf{Vlastnosti} \\
\midrule
IF-THEN pravidla & asociační pravidla, C4.5 & nejsrozumitelnější; lokální (každé pravidlo se
  čte samostatně); riziko konfliktů a~nadprodukce \\
Rozhodovací strom & ID3/C4.5/CART & globální model, snadno čitelný do hloubky 4--5; nad ni
  nepřehledný \\
Rozhodovací tabulka & diskretizovaná data & úplný výčet kombinací; přehledná jen pro málo
  atributů \\
Matematická funkce & regrese, SVM, NN & kompaktní a~přesná; obtížně interpretovatelná (kromě
  lineárního modelu s~normalizovanými vahami) \\
Prototypy / centroidy & $k$-means, $k$-medoids, LVQ & \uv{typický zástupce} skupiny; velmi
  intuitivní pro byznys \\
Pravděpodobnostní model & naivní Bayes, bayesovská síť & explicitní nejistota; graf sítě je
  čitelný \\
Grafy a~sítě & analýza sociálních sítí & vztahy mezi entitami; vizuálně silné \\
Instance (paměť případů) & $k$-NN, CBR & \uv{model} je samotná data; vysvětlení odkazem na
  podobné případy \\
\bottomrule
\end{tabular}
\end{center}

\begin{defbox}[Interpretovatelnost a~vysvětlitelnost]
\mimoq{vysvětlitelnost a~post-hoc metody; interpretovatelnost jako kritérium hodnocení modelu
je ve slidech --- viz odd.~\ref{ssec:eval}}
\emph{Interpretovatelnost} (\emph{interpretability}) je schopnost člověka porozumět tomu, jak
model dospěl k~výstupu. \emph{Vysvětlitelnost} (\emph{explainability}) je schopnost podat
srozumitelné odůvodnění konkrétního rozhodnutí.

Modely se z~tohoto hlediska dělí na dvě skupiny. \emph{Inherentně interpretovatelné modely}
umožňují nahlédnout přímo dovnitř: patří sem lineární modely, rozhodovací stromy, sady pravidel
a~prototypy. U~ostatních zbývá \emph{post-hoc vysvětlení černých skříněk}, tedy dodatečná
analýza už hotového modelu. Do této skupiny patří důležitost atributů (permutační, SHAP),
lokální surogátní modely (LIME), grafy částečné závislosti (PDP), extrakce pravidel
z~neuronové sítě, analýza citlivosti a~kontrafaktuální vysvětlení typu \uv{kdyby byl váš příjem
o~5000 vyšší, úvěr by byl schválen}.

Interpretovatelnost přitom není jen pohodlí analytika. V~regulovaných doménách, jako jsou
úvěry, pojištění nebo medicína, ji vyžaduje právo: GDPR ve svém čl.~22 zakládá právo na
vysvětlení automatizovaného rozhodnutí a~totéž platí podle AI Actu pro vysoce rizikové systémy.
\end{defbox}

\subsection{Vyhodnocování modelů}
\label{ssec:eval}

Model, který nikdo neposoudil, není znalost, ale hypotéza. Následující oddíly proto postupně
odpovídají na tři otázky: podle čeho se klasifikační metody vůbec porovnávají, jak se odhaduje
jejich chyba a~jakými mírami se popisuje kvalita klasifikace, když samotná přesnost nestačí.

\subsubsection{Kritéria hodnocení klasifikačních metod}

Kvalita modelu není jediné číslo. Klasifikační metody se posuzují podle několika kritérií
zároveň, a~ta si zpravidla navzájem konkurují:

\begin{itemize}
\item \textbf{Prediktivní přesnost} (\emph{accuracy}) je podíl správně klasifikovaných vzorů:
      \[ \mathit{Accuracy}=\frac{\text{počet správných klasifikací}}
         {\text{počet všech testovaných případů}}. \]
\item \textbf{Efektivita} se měří časem, a~to jak časem potřebným pro konstrukci modelu, tak
      časem pro jeho použití.
\item \textbf{Robustnost} vypovídá o~tom, jak si metoda poradí se šumem a~s~chybějícími
      hodnotami.
\item \textbf{Škálovatelnost} znamená, že metoda zůstane efektivní i~pro databáze uložené na
      disku.
\item \textbf{Interpretovatelnost} se ptá, jakou jasnost a~jaký vhled model svému uživateli
      poskytuje.
\item \textbf{Kompaktnost modelu} se posuzuje podle velikosti stromu, počtu pravidel a~podobně.
\end{itemize}

\subsubsection{Metody odhadu chyby}

Přesnost naměřená na datech, ze kterých se model naučil, je nadhodnocená. Následující postupy
se liší tím, jak velkou část dat obětují testování, a~volí se hlavně podle velikosti dostupné
datové sady.

\begin{defbox}[Holdout, křížová validace, leave-one-out, validační množina]
Při metodě \textbf{holdout} (\emph{holdout set}) se dostupná množina $D$ rozdělí na dvě
disjunktní podmnožiny: \emph{trénovací} $D_{train}$ slouží k~naučení modelu a~\emph{testovací}
$D_{test}$ k~jeho otestování. Platí přitom \textbf{zásadní pravidlo}, že trénovací množina se
nesmí použít při testování a~testovací při učení, protože jen neviděná testovací množina dává
\emph{nezkreslený} odhad přesnosti. Holdout je vhodný především pro \emph{velká} $D$.

\textbf{$k$-násobná křížová validace} (\emph{$k$-fold cross-validation}) rozdělí data uniformně
na $k$ stejně velkých disjunktních podmnožin. Každá podmnožina se jednou použije jako testovací
a~zbývajících $k-1$ tvoří trénovací data. Procedura se tedy provede $k$-krát a~získá se $k$
přesností $\mathit{Accuracy}_i$, jejichž průměr je výsledným odhadem:
\[ \mathit{Accuracy}_{CV}=\frac1k\sum_{i=1}^{k}\mathit{Accuracy}_i. \]
Běžně se volí 10-násobná a~5-násobná křížová validace a~celá metoda se používá spíše pro
\emph{menší} data.

\textbf{Leave-one-out} je speciální případ křížové validace, ve kterém každá testovací
podmnožina obsahuje jediný vzor a~na všem ostatním se trénuje. Pro $m$ vzorů tak dostáváme $m$-násobnou
křížovou validaci. Uplatní se u~\emph{velmi malých} dat.

\textbf{Validační množina} (\emph{validation set}) vzniká tak, že se data rozdělí nikoli na dvě,
ale na \emph{tři} části: trénovací, validační a~testovací. Validační množina slouží k~odhadu
\emph{parametrů} učicích algoritmů, protože se použijí ty hodnoty, které na ní dávají nejlepší
přesnost. K~odhadu parametrů lze místo ní použít i~křížovou validaci.
\end{defbox}

\begin{defbox}[Interval spolehlivosti pro křížovou validaci]
Odhad z~křížové validace je průměr $k$ čísel, a~proto se k~němu dá připojit i~jeho nejistota.
Přibližný $N\%$ interval spolehlivosti má tvar
\[ \mathit{Accuracy}_{CV}\pm t_{N,k-1}\cdot s_{\mathit{Accuracy}},\qquad
   s_{\mathit{Accuracy}}=\sqrt{\frac{1}{k(k-1)}\sum_{i=1}^{k}
   \big(\mathit{Accuracy}_i-\mathit{Accuracy}_{CV}\big)^2}. \]
Hodnoty $t_{N,k-1}$ pro dvoustranné intervaly jsou:
\begin{center}
\small
\begin{tabular}{@{}lrrr@{}}
\toprule
& \multicolumn{3}{c}{hladina spolehlivosti $N$} \\
$k$ & 90~\% & 95~\% & 99~\% \\
\midrule
2 & 2,92 & 4,30 & 9,92 \\
10 & 1,81 & 2,23 & 3,17 \\
30 & 1,70 & 2,04 & 2,75 \\
$\infty$ & 1,64 & 1,96 & 2,58 \\
\bottomrule
\end{tabular}
\end{center}
V~praxi se často pracuje rovnou s~$t=1{,}96$, což je hodnota pro $k\to\infty$ a~$N=95\,\%$.
\end{defbox}

\subsubsection{Matice záměn a~odvozené míry}

Přesnost je jen \emph{jedna} míra, doplněná vztahem $\mathit{error}=1-\mathit{accuracy}$,
a~v~některých aplikacích se pro hodnocení nehodí. Při dolování z~textu nás mohou zajímat pouze
dokumenty určitého tématu, tedy malá část velké kolekce. Podobně při klasifikaci silně
\emph{nevyvážených} dat, jako je detekce průniků do sítě nebo finančních podvodů, jde jen
o~minoritní třídu. Tvoří-li průniky 1~\% případů, dosáhneme 99\% přesnosti tím, že
\emph{neuděláme nic}, a~přitom neodhalíme jediný průnik. Zajímavá třída se v~této terminologii
nazývá \emph{pozitivní} a~ostatní \emph{negativní}.

\begin{defbox}[Matice záměn, precision, recall, $F_1$]
\begin{center}
\small
\begin{tabular}{@{}llcc@{}}
\toprule
& & \multicolumn{2}{c}{\textbf{klasifikováno jako}} \\
& & pozitivní & negativní \\
\midrule
\textbf{skutečnost} & pozitivní & TP & FN \\
                    & negativní & FP & TN \\
\bottomrule
\end{tabular}
\end{center}
Jednotlivá pole znamenají toto: $TP$ jsou správné klasifikace pozitivních vzorů (\emph{true
positive}), $FN$ nesprávné klasifikace pozitivních vzorů (\emph{false negative}), $FP$
nesprávné klasifikace negativních vzorů (\emph{false positive}) a~$TN$ správné klasifikace
negativních vzorů.

Z~matice se odvozují dvě základní míry:
\[ p=\mathit{precision}=\frac{TP}{TP+FP},\qquad r=\mathit{recall}=\frac{TP}{TP+FN}. \]
\emph{Precision} $p$ je počet správně klasifikovaných pozitivních vzorů dělený počtem všech
vzorů klasifikovaných jako pozitivní. \emph{Recall} $r$ je počet správně klasifikovaných
pozitivních vzorů dělený počtem všech skutečně pozitivních vzorů v~testovací množině. Podstatné
je, že \textbf{obě míry hodnotí klasifikaci pouze na pozitivní třídě.}

Obě míry spojuje do jednoho čísla \textbf{$F_1$-míra}, která je jejich \emph{harmonickým
průměrem}:
\[ F_1=\frac{2pr}{p+r}=\frac{2}{\frac1p+\frac1r}. \]
Harmonický průměr dvou čísel leží blíž tomu \emph{menšímu} z~nich, a~proto vyjde $F_1$ velká
jen tehdy, jsou-li velké \emph{obě} míry.
\end{defbox}

Jak snadno může jedna míra klamat, ukazuje \emph{ilustrace pasti}. Klasifikátor, který správně
označí jediný pozitivní vzor a~žádný negativní neoznačí chybně, má $p=100\,\%$, ale $r=1\,\%$.
Vysoká precision sama o~sobě tedy o~užitečnosti modelu neříká nic.

\subsubsection{Skórování, řazení a~lift analýza}

\emph{Skórování} (\emph{scoring}) je klasifikaci příbuzné, ale zajímá nás u~něj jediná,
pozitivní třída, například třída kupujících v~marketingové databázi. Místo definitivního
zařazení dostane každý testovací vzor \emph{odhad pravděpodobnosti} (\emph{probability
estimate}, PE), že do pozitivní třídy patří. Klasifikační systémy se ke~skórování použít dají,
musí ale pravděpodobnostní odhad poskytnout: u~rozhodovacích stromů se k~tomu hodí hodnota
confidence v~listu.

Oskórované vzory se pak \emph{seřadí} podle PE a~rozdělí do $n$ přihrádek (\emph{bins}),
například do 10. Podle počtu pozitivních vzorů v~každé přihrádce se nakreslí \textbf{lift
křivka} a~celému postupu se říká \emph{lift analýza}.

\begin{defbox}[Příklad --- lift analýza pro cílenou kampaň]
Chceme rozeslat propagační materiály potenciálním zákazníkům, konkrétně nabídku hodinek. Každá
zásilka stojí \$0,50 a~prodej jedněch hodinek přinese \$5 zisku. Otázka zní, kolik zásilek
poslat a~komu.

Testovací množina má 10\,000 vzorů, z~nichž 500 je pozitivních. Po oskórování a~seřazení do
10 přihrádek po 1000 vzorech vypadá rozdělení takto:

\begin{center}
\small
\begin{tabular}{@{}lrrrrrrrrrr@{}}
\toprule
Přihrádka & 1 & 2 & 3 & 4 & 5 & 6 & 7 & 8 & 9 & 10 \\
\midrule
Pozitivních & 210 & 120 & 60 & 40 & 22 & 18 & 12 & 7 & 6 & 5 \\
\% z~celku & 42 & 24 & 12 & 8 & 4,4 & 3,6 & 2,4 & 1,4 & 1,2 & 1,0 \\
Kumulativně & 42 & 66 & 78 & 86 & 90,4 & 94 & 96,4 & 97,8 & 99 & 100 \\
\bottomrule
\end{tabular}
\end{center}

První přihrádka obsahuje 42~\% všech pozitivních případů a~první tři dohromady 78~\%. Náhodné
rozeslání by přitom v~každé přihrádce zasáhlo 10~\% pozitivních, takže právě rozdíl mezi lift
křivkou a~touto diagonálou vyjadřuje přínos modelu. Odpověď na otázku \uv{kolik zásilek} pak
vychází z~ekonomiky: zásilkou do první přihrádky získáme
$210\cdot\$5-1000\cdot\$0{,}50=\$550$, kdežto do desáté
$5\cdot\$5-1000\cdot\$0{,}50=-\$475$.
\end{defbox}

\subsubsection{ROC křivka a~AUC}

\begin{defbox}[ROC křivka (\emph{Receiver Operating Characteristic})]
ROC křivka vynáší \emph{míru správně pozitivních} (TPR) proti \emph{míře falešně pozitivních}
(FPR):
\[ \mathit{TPR}=\frac{TP}{TP+FN},\qquad \mathit{FPR}=\frac{FP}{FP+TN},\qquad
   \mathit{TNR}=\frac{TN}{FP+TN}. \]
$\mathit{TPR}$ je podíl skutečně pozitivních případů, které jsou správně klasifikovány, takže
je to v~podstatě \emph{recall} pozitivní třídy; nazývá se též \emph{senzitivita}.
$\mathit{FPR}$ je naproti tomu podíl skutečně negativních případů klasifikovaných jako
pozitivní. Pro $\mathit{TNR}$, což je recall negativní třídy, se používá název
\emph{specificita} a~platí $\mathit{FPR}=1-\text{specificita}$.

Každá ROC křivka začíná v~bodě $(0,0)$, kdy je vše klasifikováno jako negativní, a~končí
v~$(1,1)$, kdy je vše pozitivní. \textbf{Hlavní diagonála} odpovídá náhodnému hádání, tedy
predikci pozitivní třídy s~fixní pravděpodobností, při níž $\mathit{TPR}=\mathit{FPR}$.
\end{defbox}

Křivka se sestrojuje z~klasifikace ve formě žebříčku testovacích případů. \textbf{Konstrukce
ROC křivky} má čtyři kroky:
\begin{enumerate}
\item Získej žebříček setříděním testovacích případů \emph{klesajícím} výstupem klasifikátoru,
      například aposteriorní pravděpodobností.
\item Rozděl žebříček \emph{na každém} testovacím případě na dvě části a~považuj každý případ
      v~horní části za pozitivní a~každý v~dolní za negativní.
\item Pro každé takové rozdělení spočítej dvojici $(\mathit{FPR},\mathit{TPR})$. Je-li horní
      část prázdná, dostaneme bod $(0,0)$; je-li prázdná dolní část, bod $(1,1)$.
\item Spoj sousední body.
\end{enumerate}

\begin{defbox}[AUC --- plocha pod ROC křivkou]
Dva klasifikátory se mohou v~ROC prostoru \emph{křížit}. V~příkladu ze slidů je pro
$\mathit{FPR}<0{,}43$ lepší $C_1$, zatímco pro $\mathit{FPR}>0{,}43$ lepší $C_2$, takže na
otázku \uv{který je lepší} neexistuje jednoznačná odpověď. Celkově proto lze použít
\emph{plochu pod ROC křivkou} (\emph{area under the curve}, AUC): platí-li
$\mathrm{AUC}(C_i)>\mathrm{AUC}(C_j)$, říkáme, že $C_i$ je lepší než $C_j$. Perfektní
klasifikátor má $\mathrm{AUC}=1$ a~náhodné hádání $\mathrm{AUC}=0{,}5$.
\end{defbox}

\subsection{Vizualizační techniky}

Zbývá poslední ze tří otázek: jak nalezenou znalost i~samotná data ukázat. Každá technika
odpovídá na jiný typ dotazu a~každá má hranici, za kterou přestává být čitelná.

\mimoq{tabulka jako přehled; krabicový graf, dendrogram a~ROC/lift křivka samy ve slidech jsou}
\begin{center}
\small
\begin{tabular}{@{}p{3.9cm}p{5.0cm}p{6.0cm}@{}}
\toprule
\textbf{Technika} & \textbf{K~čemu} & \textbf{Omezení / poznámka} \\
\midrule
Histogram & rozdělení jedné numerické veličiny & citlivý na volbu šířky košů \\
Krabicový graf (\emph{box plot}) & medián, kvartily, odlehlé hodnoty; srovnání skupin &
  neukáže bimodalitu (řeší \emph{violin plot}) \\
Bodový graf (\emph{scatter plot}) & vztah dvou numerických veličin & při mnoha bodech
  překryv $\to$ průhlednost, hexbin \\
Matice bodových grafů (SPLOM) & všechny dvojice atributů najednou & použitelné do zhruba
  10 atributů \\
Paralelní souřadnice & vysokorozměrná data --- objekt je lomená čára přes svislé osy
  atributů & silně závisí na pořadí os; nutné škálování; křížení čar $=$ negativní korelace \\
Teplotní mapa (\emph{heatmap}) & korelační matice, matice záměn & volba barevné škály
  (divergentní pro korelace) \\
Dendrogram & výsledek hierarchického shlukování & pořadí listů není jednoznačné \\
Projekce (PCA, MDS, t-SNE, UMAP) & 2-D pohled na vysokorozměrná data & t-SNE/UMAP zkreslují
  globální vzdálenosti --- \emph{velikosti} shluků a~mezery mezi nimi nelze interpretovat \\
ROC / lift křivka & kvalita klasifikátoru & viz odd.~\ref{ssec:eval} \\
Mozaikový graf & kontingenční tabulka kategoriálních atributů & vizuální alternativa
  k~$\chi^2$ \\
Vizualizace sítě & sociální sítě, grafy vztahů & \uv{chlupatá koule} u~velkých grafů $\to$
  agregace komunit \\
\bottomrule
\end{tabular}
\end{center}

\textbf{Vizualizace hraje v~KDD tři různé role.} První z~nich je \emph{explorativní analýza}
ve fázi porozumění datům, kdy graf odhalí rozdělení hodnot, odlehlé i~chybějící hodnoty
a~nečekané vztahy. Klasickým dokladem je Anscombeho kvartet: čtyři datové sady mají identické
průměry, rozptyly i~korelace, ale zcela odlišnou strukturu, kterou odhalí jedině obrázek. Druhou
rolí je \emph{vizualizace modelu} samotného, tedy stromu, dendrogramu, sítě nebo hranice
rozhodování. Třetí je \emph{prezentace výsledků} manažerovi, kde se sází na agregaci
a~jednoduchost podle zásady jeden graf $=$ jedno sdělení.

\textbf{Zásady dobré vizualizace} se dají shrnout do několika bodů. Poměr datové a~inkoustové
složky má být co nejvyšší, jak požaduje Tufte. Graf nesmí zkreslovat, takže sloupcové grafy
patří na nulovou osu a~osy mají být lineární, nebo zřetelně označené jako logaritmické. Barva
se volí podle typu dat: kategorie žádají kvalitativní paletu, pořadí sekvenční a~odchylka od
středu divergentní, a~to vždy s~ohledem na barvoslepost. Popisky a~jednotky patří do grafu
bezpodmínečně. Naopak trojrozměrné efekty a~koláčové grafy s~mnoha výsečemi je lepší vynechat.

\subsection{Vyhodnocení znalostí z~pohledu uživatele}

Pojem \uv{vyhodnocování získaných znalostí} se používá ve třech různých rovinách a~ty je nutné
od sebe odlišovat:
\begin{itemize}
\item \emph{Statistické vyhodnocení modelu} pracuje s~přesností, maticí záměn, ROC a~AUC
      a~s~odhadem chyby křížovou validací (odd.~\ref{ssec:eval}).
\item \emph{Vyhodnocení zajímavosti jednotlivých vzorů} se opírá o~support, confidence, lift,
      $\chi^2$, t-weight a~d-weight (odd.~\ref{ssec:ar} a~kap.~\ref{sec:pravidla}).
\item \emph{Vyhodnocení z~pohledu uživatele} se ptá, zda je nalezená znalost srozumitelná,
      platná, užitečná a~nová.
\end{itemize}
První dvě roviny jsou objektivní a~měřitelné, kdežto třetí je nutně subjektivní: tentýž vzor
může být pro jednoho uživatele objevem a~pro druhého notorietou. Část výsledků bývá nezajímavá
nebo z~hlediska uživatele zřejmá, část se dá použít přímo a~část je nutné prezentovat
srozumitelněji. Užitečné bývá shrnout výsledky do \emph{analytické zprávy}, přičemž výstupem
celého procesu může být i~provedení rozumné akce.

\subsection{Shrnutí ke~zkoušce}

\begin{itemize}
\item \textbf{Formy reprezentace znalostí} sahají od IF-THEN pravidel (nejsrozumitelnější,
      lokální) přes rozhodovací strom (globální, čitelný do hloubky 4--5), rozhodovací tabulku,
      matematickou funkci (kompaktní, ale neinterpretovatelná), prototypy a~centroidy
      a~pravděpodobnostní model až po grafy a~sítě a~instance ($k$-NN, CBR). U~každé formy je
      třeba umět říct, z~jaké metody vzniká a~čím se platí za její srozumitelnost.
\item \textbf{Kritéria hodnocení metod} jsou prediktivní přesnost, efektivita (čas konstrukce
      i~čas použití), robustnost, škálovatelnost, interpretovatelnost a~kompaktnost modelu.
\item \textbf{Odhad chyby} má čtyři podoby. \emph{Holdout} se hodí na velká data a~trénovací
      s~testovací množinou se u~něj nesmí míchat. \emph{$k$-násobná křížová validace} se
      používá na menší data, běžné je $k$ rovné 5 nebo 10, přesnost vyjde jako průměr $k$
      přesností a~doplní se interval spolehlivosti s~$t_{N,k-1}$. \emph{Leave-one-out} je
      $m$-násobná CV pro velmi malá data. \emph{Validační množina} je třetí část dat, na které
      se ladí \emph{parametry} algoritmu.
\item \textbf{Matice záměn} obsahuje $TP$, $FN$, $FP$ a~$TN$, z~nichž se počítá
      $p=\frac{TP}{TP+FP}$, $r=\frac{TP}{TP+FN}$ a~$F_1=\frac{2pr}{p+r}$, tedy harmonický
      průměr, který leží blíž menšímu z~obou čísel. Je nutné umět vysvětlit, \emph{proč}
      accuracy nestačí u~nevyvážených dat: při 1~\% průniků dá strategie \uv{nedělat nic}
      $\Rightarrow$ 99~\% přesnosti.
\item \textbf{ROC} vynáší $\mathit{TPR}=\frac{TP}{TP+FN}$, což je senzitivita čili recall
      pozitivní třídy, proti $\mathit{FPR}=\frac{FP}{FP+TN}$; $\mathit{TNR}$ je specificita
      a~rovná se $1-\mathit{FPR}$. Křivka vede z~$(0,0)$ do $(1,1)$ a~diagonála odpovídá
      náhodnému hádání. Je třeba umět popsat konstrukci ze žebříčku, tedy dělení na každém
      případu. \textbf{AUC} má hodnotu 1 pro perfektní a~0,5 pro náhodný klasifikátor
      a~potřebujeme ji proto, že se ROC křivky dvou klasifikátorů mohou \emph{křížit}.
\item \textbf{Skórování a~lift}: místo třídy se každému vzoru přiřadí odhad pravděpodobnosti
      (u~stromu confidence v~listu), vzory se seřadí a~rozdělí do přihrádek a~lift křivka se
      poměřuje s~diagonálou náhodného výběru. U~zkoušky se k~tomu očekává ekonomická úvaha
      o~tom, kolik zásilek se vyplatí poslat.
\item \textbf{Tři roviny \uv{vyhodnocování}} jsou statistická kvalita modelu, zajímavost vzorů
      a~vyhodnocení uživatelem, přičemž první dvě jsou objektivní a~třetí subjektivní.
      \textbf{Čtyři kritéria zajímavosti znalosti} zní: srozumitelnost, platnost, užitečnost
      a~novost.
\item \textbf{Vizualizace} se vybírá podle úlohy: histogram, box plot, scatter, SPLOM,
      paralelní souřadnice, heatmapa, dendrogram, projekce, ROC a~lift křivka. V~KDD má
      tři role: exploraci (dokládá ji Anscombeho kvartet), vizualizaci modelu a~prezentaci.
      Zásady jsou Tufteho poměr datové a~inkoustové složky, nezkreslovat osami a~volit barvu
      podle typu dat. Znát je nutné i~omezení: t-SNE a~UMAP zkreslují globální vzdálenosti.
\end{itemize}

%=====================================================================
\newpage

\section{Modely pro analýzu sociálních sítí, centralita a~detekce komunit}
\label{sec:sna}
%=====================================================================

Věta okruhu jmenuje tři věci: \emph{modely} pro analýzu sociálních sítí, \emph{míry
centrality} a~\emph{detekci komunit}. Tomu odpovídá i~členění kapitoly. Nejprve přijde na
řadu reprezentace sítí a~míry, kterými se popisuje jednotlivý aktér (centralita) a~síť jako
celek (koheze). Na to navazují modely: jaké vlastnosti mají reálné sítě, jak se v~nich šíří
vliv, co znamená homofilie, segregace a~strukturní balance a~jak se analyzují vazby (HITS,
PageRank). Poslední vrstvu tvoří detekce komunit a~predikce vazeb.

\subsection{Vymezení oboru a~motivace}

\begin{defbox}[Analýza sociálních sítí]
\emph{Analýza sociálních sítí} (\emph{social network analysis}, SNA) je interdisciplinární
obor, který vyrostl ze tří oblastí:
\begin{itemize}
\item \emph{sociální vědy} dodávají otázky a~problémy, tedy snahu vysvětlit sociální chování
      analýzou společností;
\item \emph{analýza sítí} umožňuje formulovat a~řešit problémy, které mají síťovou strukturu;
\item \emph{teorie grafů} poskytuje matematický aparát pro reprezentaci a~analýzu grafů.
\end{itemize}
Zásadní je posun perspektivy: SNA se soustředí na \textbf{vztahy} mezi sociálními entitami
(jednotlivci, skupinami, institucemi) \emph{spíše než na entity samotné} a~jejich atributy.
Studium společnosti ze síťové perspektivy chápe jednotlivce jako \emph{zasazené} do sítě
vztahů a~vysvětlení sociálního chování hledá ve struktuře těchto sítí, nikoli
v~jednotlivcích samotných.
\end{defbox}

Praktickou motivací je \textbf{analýza vazeb} (\emph{link analysis}), jejímž cílem je najít
vztahy mezi daty a~tyto vazby vizualizovat. Uplatní se v~telekomunikacích,
v~kriminalistice a~právu, kde jsou skupiny zločinců vzájemně propojené a~analýza vztahů může
pomoci je odhalit, a~v~marketingu, kde jde o~vztahy mezi zákazníky.

\textbf{Šest stupňů odloučení} (Milgram, 1967) je experiment o~fenoménu malého světa.
Náhodně vybraní lidé z~Nebrasky měli poslat dopis burzovnímu makléři v~Bostonu, a~to jen
přes prostředníky: poslat jej směli výhradně někomu, koho znali osobně, tedy na křestní
jméno. Mezi dopisy, které cíl našly, byl průměrný počet vazeb \emph{šest}. Do cíle jich
ovšem mnoho nedorazilo. Experiment předznamenal pozdější objevy, zejména huby a~tendenci
postupovat podle geografie či profese. Podobně vypovídají \emph{Erdősova čísla}: Paul Erdős
napsal přes 1500 článků s~507 spoluautory, průměrné Erdősovo číslo mezi matematiky je 4,65
a~maximum 13.

Druhou motivací jsou \textbf{bezškálové sítě} (\emph{scale-free}, SF), v~nichž mají některé
uzly extrémně mnoho spojení (\emph{huby}), zatímco většina uzlů jen několik. Důsledkem je
dvojaká povaha takové sítě: \emph{odolnost proti náhodné poruše} a~současně
\emph{zranitelnost koordinovaným útokem}. Selhat může až 80~\% náhodně vybraných uzlů, aniž
se síť fragmentuje, ale eliminace 5--15~\% \emph{hubů} způsobí rozpad celého systému.

Aplikace této úvahy sahají daleko. V~ochraně proti šíření počítačových virů internetem
z~ní plyne, že vymýcení internetových virů je v~SF architektuře v~podstatě nemožné.
V~medicíně vede k~vakcinačním kampaním cíleným na huby a~k~novým lékům mířícím na klíčové
molekuly. V~byznysu vysvětluje kaskádové finanční krachy a~podpírá marketing, tedy šíření
\uv{nákazy} a~influencer marketing.

Bezškálovou strukturu mají \emph{sociální} sítě (vědecká spolupráce prostřednictvím
spoluautorství článků, Hollywood a~herci ve stejném filmu), \emph{biologické} sítě
(metabolismus buňky, proteinové regulační sítě) i~\emph{sociotechnické} sítě (internet
tvořený routery a~spoji, WWW tvořené stránkami a~URL).

\subsection{Reprezentace sítí}

Než se dá cokoli z~toho měřit, musí být síť zapsaná, a~to způsobem, s~nímž umí pracovat
algoritmus.

\begin{defbox}[Základní reprezentace]
Síť je graf $G=(V,E)$, kde \emph{uzly} (vrcholy) jsou aktéři a~\emph{hrany} (vazby) vztahy
mezi nimi. Zapsat ji lze třemi ekvivalentními způsoby:
\begin{itemize}
\item \textbf{obrázkem grafu};
\item \textbf{seznamem hran} (\emph{edge list}), tedy dvojicemi uzlů, u~ohodnocených
      grafů doplněnými o~sloupec vah;
\item \textbf{maticí sousednosti} (\emph{adjacency matrix}) $A=(a_{ij})$, kde $a_{ij}=1$,
      je-li uzel $i$ spojen s~uzlem $j$, a~0 jinak; u~vážených grafů obsahuje matice
      váhy.
\end{itemize}
U~\textbf{orientovaných} grafů (kdo koho kontaktuje) matice sousednosti \emph{nemusí} být
symetrická, u~\textbf{neorientovaných} (kdo koho zná) symetrická je. Tentýž seznam hran má
proto v~obou případech jinou interpretaci.
\end{defbox}

\textbf{Váhy hran} nesou informaci o~tom, jak silný vztah hrana zachycuje, a~jejich význam
závisí na tom, co se v~dané úloze měří. Vyjadřovat mohou frekvenci interakce v~období
pozorování nebo počet vyměněných předmětů. Jinde jde o~individuální percepci silné vazby,
o~náklady komunikace nebo výměny (například o~vzdálenost), případně o~kombinaci předchozích.
Samy hrany pak reprezentují interakce, toky informací či zboží, podobnosti a~afiliace, nebo
sociální vztahy.

\textbf{Ego sítě a~\uv{celé} sítě.} \emph{Ego síť} je osobní síť jednoho aktéra
(\emph{ego}) a~jeho sousedů (\emph{alterů}). Samotné ego přitom pro další analýzu
\emph{není} potřebné, protože všichni alteři jsou s~egem spojeni už z~definice. \uv{Celá}
síť oproti tomu obsahuje i~izolované uzly (\emph{isolates}). Platí ovšem, že \textbf{žádná
studovaná síť není v~praxi \uv{celá}}: obvykle jde jen o~částečný obraz skutečné sítě, čemuž
se říká \emph{problém specifikace hranic} (\emph{boundary specification problem}).

\subsection{Síla vazeb}

Sociální média propojují uživatele extrémně snadno, takže uživatel může mít online tisíce
\uv{přátel}. Tyto vazby ale nejsou stejně silné a~hrají různou roli při formování komunit
a~difuzi informací. \emph{Silné vazby} jsou blízcí přátelé, \emph{slabé vazby} známí. Právě
na ty druhé upozornil Granovetter (1973) tezí o~\textbf{síle slabých vazeb}: občasná setkání
se vzdálenými známými mohou poskytnout důležitou informaci o~nových příležitostech,
například při hledání práce.

\begin{defbox}[Homofilie, tranzitivita, přemostění]
\textbf{Homofilie} (\emph{homophily}) je tendence vztahovat se k~lidem s~podobnými
charakteristikami (status, přesvědčení, \dots). Vede k~tvorbě homogenních skupin (klastrů),
v~nichž je snazší vytvářet vztahy. \emph{Extrémní} homogenizace ale může působit proti
inovaci a~generování idejí, takže \emph{heterofilie} je v~některých kontextech žádoucí.
Homofilní vazby mohou být silné i~slabé.

\textbf{Tranzitivita} (\emph{transitivity}) říká, že \uv{existuje-li vazba mezi $A$ a~$B$
a~mezi $B$ a~$C$, budou v~tranzitivní síti spojeni i~$A$ a~$C$.} Silné vazby jsou tranzitivní
častěji než slabé, takže tranzitivita existenci silných vazeb \emph{indikuje}, není to však
ani nutná, ani postačující podmínka. Tranzitivita a~homofilie společně vedou ke vzniku
\emph{klik}, tedy plně propojených klastrů.

\textbf{Přemostění} (\emph{bridging}) stojí na pojmu \emph{mostu} (\emph{bridge}): most je
uzel či hrana spojující různé skupiny. Obvykle jde o~slabé vazby, ale ne každá slabá vazba je
most. Mosty usnadňují mezigrupovou komunikaci, zvyšují sociální kohezi a~podporují inovaci.
\end{defbox}

\textbf{Určování síly vazby z~topologie sítě.} Mosty spojující dvě různé komunity jsou slabé
vazby a~hrana je \emph{most} tehdy, jestliže její odstranění rozpojí její koncové uzly.
V~reálných sítích ale mosty nejsou časté, a~proto se definice \emph{uvolní} na
\textbf{zkratkový most} (\emph{shortcut bridge}): kontroluje se, zda se po odstranění hrany
\emph{zvětší} vzdálenost jejích koncových uzlů, konkrétně na více než 2. Čím větší tato
vzdálenost je, tím \emph{slabší} vazba. Například $d(2,5)=4$ po odstranění hrany $(2,5)$,
zatímco $d(5,6)=2$ po odstranění $(5,6)$ $\Rightarrow$ $(5,6)$ je silnější vazba než $(2,5)$.

\begin{defbox}[Překryv sousedství (\emph{neighborhood overlap})]
Pro dva uzly $i,j$ spojené vazbou platí, že čím větší je jejich překryv, tím \emph{silnější}
je vazba mezi nimi:
\[ O(i,j)=\frac{\text{počet společných přátel } i \text{ a } j}
   {\text{počet přátel, kteří jsou sousedy alespoň }i\text{ nebo }j}
   =\frac{|N_i\cap N_j|}{|N_i\cup N_j|-2}, \]
kde $-2$ ve jmenovateli vylučuje $i$ a~$j$ samotné.

\emph{Příklady:} $O(2,5)=0$; $O(5,6)=\frac{|\{4\}|}{|\{2,4,5,6,7,10\}|-2}=\frac14$;
pro graf $G=(\{1,2\},\{(1,2)\})$ je $O(1,2)=0$.
\end{defbox}

\textbf{Určování síly vazby z~profilů a~interakcí.} Na Twitteru, kde lze sledovat jiné bez
jejich potvrzení, není skutečná síť přátelství určena sítí follower--followee, nýbrž
\emph{frekvencí}, s~jakou spolu dva uživatelé mluví. A~právě tato skutečná síť má silnější
vliv na používání Twitteru. Na Facebooku lze sílu vazeb predikovat z~různých informací:
z~příspěvků iniciovaných přítelem, ze zpráv vyměněných na zdi nebo z~počtu společných
přátel. Numerickou sílu vazby se lze i~učit, například metodou maximální věrohodnosti:
podobnost profilů určuje sílu vazby a~síla vazby určuje interakci uživatelů, takže se
maximalizuje věrohodnost na základě pozorovaných profilů a~interakcí.

\subsection{Klíčoví hráči: míry centrality}

Následující míry se všechny ptají, kdo je v~síti klíčový. Každá z~nich ale měří jiný druh
důležitosti, a~proto se jejich pořadí uzlů zpravidla neshodují.

\begin{defbox}[Stupňová centralita (\emph{degree centrality})]
Vstupní stupeň uzlu je počet vazeb vedoucích do uzlu a~výstupní stupeň počet vazeb vedoucích
z~uzlu. V~neorientovaných grafech jsou oba stupně shodné.

Míra zachycuje \emph{propojenost}, \emph{vliv} nebo \emph{popularitu} uzlu a~pomáhá zjistit,
které uzly jsou centrální z~hlediska šíření informace a~ovlivňování ostatních ve svém
\uv{sousedství}.
\end{defbox}

Stupňová centralita vystačí s~bezprostředním okolím uzlu. Následující dvě míry, mezilehlá
a~blízkostní, už ale potřebují vědět, jak daleko od sebe uzly leží, a~tedy pojem cesty.

\begin{defbox}[Cesty a~nejkratší cesty]
\emph{Cesta} mezi dvěma uzly je libovolná posloupnost neopakujících se uzlů, která je spojuje.
\emph{Nejkratší cesta} spojuje dva uzly nejmenším počtem hran a~tento počet se nazývá
\emph{vzdálenost} uzlů. Kratší cesty jsou preferované pro rychlou komunikaci či výměnu,
i~když ne vždy; příkladem opačné situace je šíření nemoci.

\textbf{Algoritmy nejkratších cest.} Volba algoritmu se řídí dvěma otázkami: hledáme cesty
z~jednoho zdroje, nebo mezi všemi páry uzlů, a~připouštíme záporné váhy? Pro \emph{jeden
zdroj} a~obecné grafy, tedy orientované a~se zápornými vahami, slouží Bellmanův--Fordův
algoritmus se složitostí $O(|V||E|)$. Nemá-li graf záporné váhy, stačí Dijkstrův algoritmus,
$O(|V|^2)$. Hledá-li se cesta jen mezi dvěma uzly, lze algoritmus zastavit hned po jejím
nalezení.

Pro \emph{všechny páry} uzlů se při povolených záporných váhách používá Floydův--Warshallův
algoritmus, $O(|V|^3)$. U~řidších grafů se zápornými vahami je výhodnější Johnsonův
algoritmus, $O(|E||V|+|V|^2\log_2|V|)$, a~bez záporných vah opět Dijkstra, $O(|V|^3)$.
\end{defbox}

\begin{defbox}[Mezilehlá centralita (\emph{betweenness centrality})]
Mezilehlá centralita je počet nejkratších cest procházejících uzlem, dělený všemi nejkratšími
cestami v~síti. Pro uzel $v$ se pro každou dvojici $s,t$ spočítá počet nejkratších cest mezi
$s$ a~$t$ procházejících $v$, ten se vydělí všemi nejkratšími cestami mezi $s$ a~$t$
a~výsledné hodnoty se sečtou přes všechny dvojice. Někdy se výsledek ještě normalizuje tak,
aby nejvyšší hodnota byla 1.

Míra indikuje uzly, které leží na \emph{komunikačních cestách} mezi jinými uzly, a~pomáhá
určit místa, kde by se síť mohla \emph{rozpadnout}.
\end{defbox}

\begin{defbox}[Blízkostní centralita (\emph{closeness centrality})]
Blízkostní centralita je míra \emph{dosahu}, tedy rychlosti, s~jakou se informace z~daného
uzlu dostane k~ostatním. Spočítá se \emph{průměrná} délka všech nejkratších cest z~uzlu ke
všem ostatním uzlům sítě, jinak řečeno kolik \uv{skoků} je v~průměru potřeba, a~vezme se
\emph{převrácená hodnota}. Právě kvůli tomuto převrácení znamenají vyšší hodnoty vyšší
blízkost a~jsou \uv{lepší}. Míra se hodí tam, kde jde především o~rychlost šíření.
\end{defbox}

\begin{defbox}[Vlastní centralita (\emph{eigenvector centrality})]
Vlastní centralita uzlu je proporcionální součtu vlastních centralit všech uzlů přímo
spojených s~uvažovaným uzlem:
uzel s~vysokou vlastní centralitou je spojen s~jinými uzly s~vysokou vlastní centralitou.
Metrika vychází z~myšlenky, že moc a~status aktéra jsou \emph{rekurzivně} definovány mocí
a~statusem jeho sousedů, a~měří tedy, jak dobře je aktér spojen s~jinými \emph{dobře
propojenými} aktéry.

Formálně: pro $G=(V,E)$ s~maticí sousednosti $A=(a_{i,j})$ platí
\[ x_i=\frac1\lambda\sum_{j\in M(i)}x_j=\frac1\lambda\sum_{j\in V}a_{i,j}x_j
   \quad\Longleftrightarrow\quad Ax=\lambda x, \]
kde $M(i)$ je množina sousedů $i$. Vlastních čísel $\lambda$, pro něž nenulové řešení
existuje, je mnoho. \emph{Požadavek nezápornosti} všech složek vlastního vektoru však
znamená, že žádanou míru centrality dává \emph{jen největší} vlastní číslo. Vlastní vektor
je navíc určen až na společný násobek, takže je nutná normalizace, například součtem přes
všechny uzly $=1$ nebo $=|V|$.

Míra blízce souvisí s~Google PageRank, kde odkazy ze silně odkazovaných stránek váží více
(odd.~\ref{ssec:pagerank}).
\end{defbox}

Rozdíl mezi čtyřmi mírami je nejlépe vidět na otázce, kterou každá z~nich klade, a~na tom, co
znamená v~konkrétní síti:

\begin{center}
\small
\begin{tabular}{@{}p{2.7cm}p{12.6cm}@{}}
\toprule
\textbf{Míra} & \textbf{Interpretace v~sociálních sítích} \\
\midrule
Stupňová & Kolik lidí může tato osoba dosáhnout \emph{přímo}? --- \emph{Síť hudebních
  spoluprací:} s~kolika lidmi spolupracoval? \\
Mezilehlá & Jak pravděpodobně leží tato osoba na nejpřímější cestě v~síti? --- \emph{Síť
  špionů:} přes koho nejspíš protéká většina informací? \\
Blízkostní & Jak rychle může tato osoba dosáhnout všech v~síti? --- \emph{Síť sexuálních
  vztahů:} jak rychle se do zbytku sítě rozšíří přenosná nemoc? \\
Vlastní & Jak dobře je tato osoba spojena s~jinými dobře propojenými lidmi? --- \emph{Citační
  síť:} který autor je nejvíce citován jinými dobře citovanými autory? \\
\bottomrule
\end{tabular}
\end{center}

\textbf{Množiny klíčových hráčů.} Jednotlivé míry samy o~sobě nestačí. V~příkladu ze slidů je
uzel 10 nejcentrálnější podle stupňové centrality, ale uzly 3 a~5 \emph{společně} dosáhnou
více uzlů. Navíc je vazba mezi 3 a~5 \emph{kritická}: při jejím přerušení se síť rozpadne na
dvě izolované podsítě. Za jinak stejných podmínek jsou tedy uzly 3 a~5 \emph{společně} pro
síť důležitější než uzel 10 samotný. \emph{Uvažovat množiny klíčových hráčů proto pomáhá.}

\begin{defbox}[Strukturní ekvivalence]
Vyjadřuje \emph{podobnost} mezi aktéry v~sociální síti a~je založena na sousedech, které
sdílejí, ekvivalentně na počtu identických vazeb. Dva aktéři jsou \emph{strukturně
ekvivalentní}, lze-li jejich pozice vzájemně vyměnit \emph{bez změny celkové struktury}
sítě. Přibližnou strukturní ekvivalenci lze použít pro hierarchické shlukování sociálních
sítí.
\end{defbox}

\subsection{Koheze: míry na úrovni sítě}

Zatímco centralita hodnotí jednotlivý uzel, následující míry charakterizují síť jako celek.

\begin{defbox}[Reciprocita, hustota, klastrovací koeficient]
\textbf{Reciprocita} (\emph{reciprocity}) měří vzájemnost a~reciproční výměnu v~síti a~souvisí
se sociální kohezí. Je to podíl počtu \emph{reciprokých} vztahů, tedy takových, kde hrana
vede v~obou směrech, na celkovém počtu vztahů v~síti, přičemž dva uzly jsou \uv{ve vztahu},
existuje-li mezi nimi alespoň jedna hrana. \textbf{Má smysl jen v~orientovaných grafech.}

\textbf{Hustota} (\emph{density}) je podíl počtu hran v~síti na celkovém počtu \emph{možných}
hran:
\[ \text{density}=\frac{|E|}{|V|(|V|-1)/2}\quad\text{(neorientovaný graf)},\qquad
   \frac{|E|}{|V|(|V|-1)}\quad\text{(orientovaný)}. \]
Je to běžná míra toho, jak dobře je síť propojená, tedy jak \uv{těsně upletená} je. Perfektně
propojená síť se nazývá \emph{klika} a~má hustotu 1. Orientované grafy mají poloviční hustotu
než jejich neorientované protějšky, protože možných hran je v~nich dvakrát více.

\textbf{Klastrovací koeficient} (\emph{clustering coefficient}) uzlu $i$ měří hustotu jeho
\emph{sousedství}, tedy sítě tvořené pouze uzly přímo spojenými s~$i$. Pro $G=(V,E)$ je
sousedství $N_i=\{v_j: e_{ij}\in E\ \vee\ e_{ji}\in E\}$, $k_i=|N_i|$, a
\[ C_i=\frac{2\,|\{e_{jk}: v_j,v_k\in N_i,\ e_{jk}\in E\}|}{k_i(k_i-1)}
   \ \text{(neorientovaný)},\qquad
   C_i=\frac{|\{e_{jk}\}|}{k_i(k_i-1)}\ \text{(orientovaný)}. \]
\emph{Klastrovací koeficient sítě} $C$ je průměr koeficientů všech uzlů, tedy
$C=\frac{1}{|V|}\sum_{i=1}^{|V|}C_i$. Indikuje přítomnost (sub)komunit a~shlukovací
algoritmy se snaží maximalizovat počet hran padnoucích do téhož klastru.
\end{defbox}

\begin{defbox}[Průměr, průměrná vzdálenost, excentricita, malé světy]
\textbf{Průměr} sítě (\emph{diameter}) je nejdelší nejkratší cesta, tedy vzdálenost, mezi
libovolnými dvěma uzly. Je to užitečná míra \emph{dosahu} sítě, protože udává, jak dlouho
nejdéle trvá dosáhnout libovolného uzlu. Řidší sítě mají obecně větší průměr.

\textbf{Průměrná vzdálenost} je průměr všech nejkratších cest v~síti a~udává, jak daleko od
sebe budou dva uzly v~průměru.

\textbf{Excentricita} uzlu je největší geodetická vzdálenost mezi daným uzlem
a~libovolným jiným uzlem grafu.

\textbf{Malý svět} (\emph{small world}) je síť, která vypadá téměř náhodně, ale má
\emph{významně vysoký klastrovací koeficient} a~\emph{relativně krátkou průměrnou délku
cesty}: uzly se lokálně klastrují a~zároveň jsou dosažitelné v~několika krocích.
V~sociálních sítích je to velmi častý jev, a~to kvůli \emph{tranzitivitě} silných vazeb
a~\emph{schopnosti slabých vazeb dosáhnout přes klastry}. Sítě malého světa mají mnoho
klastrů, ale i~mnoho mostů mezi nimi, které zkracují průměrnou vzdálenost.
\end{defbox}

\subsection{Model reálných komplexních sítí a~jejich vlastnosti}

Popsané míry se nyní spojí do obrazu toho, čím se reálné sítě liší od grafů, které bychom
vygenerovali náhodně.

\begin{defbox}[Šest charakteristických vlastností reálných komplexních sítí]
Reálné komplexní sítě jsou \emph{nenáhodné a~neregulární} grafy s~unikátními rysy; patří mezi
ně sociální sítě, informační a~znalostní sítě, technologické sítě a~biologické sítě.
Opakovaně se u~nich pozoruje šest vlastností:
\begin{enumerate}
\item \textbf{Efekt malého světa.} Na existenci malých světů v~reálných sociálních sítích
      poukázal jako první Stanley Milgram experimentem z~60.~let, jehož se zúčastnilo asi 300
      lidí a~v~němž byla mediánová délka úspěšných cest 6. Vychází z~konceptu \uv{šesti stupňů
      odloučení}: každý je od každého v~průměru šest kroků. Důsledky se týkají rychlosti
      šíření informací i~nemocí.
\item \textbf{Tranzitivita / klastrování.} V~reálných, a~zejména v~sociálních sítích je
      vysoká pravděpodobnost, že uvnitř větších sítí najdeme \emph{úplné podsítě}, v~nichž
      každý zná každého.
\item \textbf{Mocninné rozdělení stupňů.} \emph{Rozdělení stupňů} je rozdělení
      pravděpodobnosti stupňů uzlů přes celou síť, tedy histogram podílu uzlů se stupněm $k$.
      Reálné sítě mají rozdělení stupňů zcela odlišné od náhodných grafů.
\item \textbf{Odolnost sítě} (\emph{resilience}) popisuje, jak se změní propojenost sítě při
      odstranění jednoho či více uzlů. Většina sítí je \emph{robustní} proti náhodnému
      odstranění, ale podstatně méně už proti \emph{cílenému} odstranění uzlů s~nejvyšším
      stupněm. Odstraní-li se \uv{vrátní} (\emph{gatekeepers}), výrazně se změní komunikační
      schopnost sítě, protože se některé uzly odpojí. Odstranění jediného uzlu obvykle
      strukturu neohrozí, a~proto je vhodnější testovat odolnost odstraněním určitého
      \emph{procenta} uzlů.
\item \textbf{Vzorce míchání} (\emph{mixing patterns}). Koexistují-li v~síti různé typy uzlů,
      bývá běžná jistá \emph{selektivita} ve tvorbě spojení. Selektivnímu propojování se říká
      \emph{asortativní míchání} (\emph{assortative mixing}) čili homofilie a~klasickým
      příkladem je míchání podle rasy. Reálné sítě vykazují vyšší tendenci k~asortativnímu
      míchání.
\item \textbf{Komunitní struktura.} Většina reálných sociálních sítí vykazuje komunitní
      strukturu, která vzniká jako důsledek \emph{globální i~lokální heterogenity} rozložení
      hran: v~některých oblastech grafu nacházíme vysoké koncentrace hran (\emph{komunity})
      a~mezi těmito oblastmi koncentrace nízké.
\end{enumerate}
\end{defbox}

\begin{defbox}[Bezškálové sítě a~preferenční připojování]
\textbf{Bezškálové sítě} (Barabási, Albert, 1999) jsou sítě, jejichž rozdělení stupňů
následuje \emph{mocninný zákon}; patří sem WWW, citační sítě, biologické sítě a~některé
sociální sítě. Mocninná rozdělení vznikají typicky tehdy, \emph{závisí-li množství, které
něčeho dostanete, na množství, které už máte}. Topologicky se to projeví jako síť s~několika
silně propojenými uzly (\emph{huby}) a~velkým počtem uzlů s~nízkým stupněm.

Tentýž mechanismus dostal postupně několik jmen. Nejprve se mu říkalo
\emph{rich-get-richer / poor-get-poorer} strategie, ekvivalentně \emph{Matoušův efekt}; Price
jej nazval \emph{kumulativní výhoda} a~Barabási s~Albertem \textbf{preferenční připojování}
(\emph{preferential attachment}). Vysvětlení je prosté: nové uzly vstupující do sítě se
připojují k~dobře propojeným uzlům, které bývají spojeny s~centrálními a~prestižními pozicemi
(více statusu, popularity, znalostí, peněz).

Dva hlavní mechanismy vzniku SF sítě jsou tedy \textbf{růst} a~\textbf{preferenční
připojování}. Výsledná síť má \emph{dlouhý chvost} rozdělení stupňů a~\emph{strukturu malého
světa}. Tranzitivita a~vlastnosti silných a~slabých vazeb jsou běžné a~mohou ke struktuře
malého světa vést, k~jejímu vysvětlení ale \emph{nutné} nejsou.

\textbf{Důvody preferenčního připojování} se dělí do tří skupin. \emph{Popularita}
(\uv{bohatí bohatnou}) znamená, že chceme být spojeni s~populárními lidmi, idejemi a~věcmi,
což jejich popularitu dále zvyšuje, a~to bez ohledu na objektivní měřitelné charakteristiky.
\emph{Kvalita} (\uv{dobří se zlepšují}) předpokládá, že hodnotíme podle objektivních kritérií
kvality, takže kvalitnější uzly přirozeně přitáhnou pozornost rychleji. \emph{Smíšený model}
shrnuje obojí do věty \uv{nelze předpovědět, kdo se stane hvězdou, i~když kvalita hraje
roli}: mezi uzly s~podobnými atributy se \uv{hvězdami} stanou ty, které první dosáhnou
kritického množství (\emph{halo efekt}).
\end{defbox}

\begin{defbox}[Struktura jádro--periferie a~centralizace]
\textbf{Centralizace} (\emph{centralization}) měří, do jaké míry je síť centralizovaná či
decentralizovaná. Vychází z~\emph{rozdílů ve stupních} mezi uzly a~obvykle se normalizuje do
$[0,1]$. Síť, která silně závisí na 1--2 vysoce propojených uzlech, například jako důsledek
preferenčního připojování, vykazuje větší rozdíly ve stupňové centralitě. Centralizované
struktury mohou být lepší v~některých úlohách vyžadujících koordinaci, třeba při týmovém
řešení problémů, ale jsou \emph{náchylnější k~selhání}, odpojí-li se klíčoví hráči.

Mnoho velkých skupin a~komunit má \emph{jádro} hustě propojených uživatelů, které je kritické
pro propojení mnohem větší \emph{periferie}. Jádra lze identifikovat vizuálně, nebo zkoumáním
polohy uzlů s~vysokým stupněm a~jejich společných rozdělení stupňů; k~rozlišení jádra
a~periferních elementů slouží i~\emph{bow-tie analýza}, použitá pro analýzu struktury webu.
\end{defbox}

\subsection{Modelování vlivu}

Zbývá popsat, jak se sítí něco šíří. Modely vlivu se v~detailech liší, sdílejí ale několik
předpokladů.

\begin{defbox}[Společné vlastnosti metod modelování vlivu]
\begin{itemize}
\item Sociální síť je reprezentována \emph{orientovaným} grafem a~každý aktér je jeden uzel.
\item Každý uzel začíná buď jako \emph{aktivní}, nebo jako \emph{neaktivní}.
\item Aktivovaný uzel aktivuje své sousedy.
\item \textbf{Jednou aktivovaný uzel už nelze deaktivovat.} Šíření je tedy nevratné a~množina
      aktivních uzlů může jen růst.
\end{itemize}
\end{defbox}

\begin{defbox}[Model lineárního prahu (LTM)]
Aktér (uzel) podnikne akci, pokud počet jeho přátel, kteří akci podnikli, překročí jistý
\emph{prah}.
\begin{itemize}
\item Každý uzel $v$ si zvolí prah $\theta_v$ \emph{náhodně} z~uniformního rozdělení na
      $[0,1]$.
\item V~každém diskrétním kroku zůstávají aktivní všechny uzly, které byly aktivní v~kroku
      předchozím.
\item Nově se aktivují uzly splňující
      \[ \sum_{u\in N_v^{\text{aktivní}}} w_{u,v}\ \ge\ \theta_v
         \sum_{u\in N_v} w_{u,v}. \]
\end{itemize}
LTM se soustředí na \emph{příjemce} vlivu.
\end{defbox}

\begin{defbox}[Model nezávislých kaskád (ICM)]
Soustředí se naopak na \emph{vysílající} uzly.
\begin{itemize}
\item Uzel $v$ aktivovaný v~kroku $t$ má \emph{jedinou} příležitost aktivovat každého ze svých
      sousedů, a~to náhodně.
\item Pro sousední uzel $u$ uspěje aktivace s~pravděpodobností $p_{v,u}$ (např.\ $p=0{,}5$).
\item Uspěje-li, $u$ se stane aktivním v~kroku $t+1$.
\item V~dalších kolech se už $v$ o~aktivaci $u$ \emph{nepokouší}.
\item Difuze začíná počáteční aktivovanou množinou uzlů a~pokračuje, dokud není možná žádná
      další aktivace.
\end{itemize}
\end{defbox}

\begin{defbox}[Maximalizace vlivu]
\textbf{Úloha:} pro danou síť a~parametr $k$ vybrat $k$ uzlů do aktivační množiny $B$ tak, aby
se maximalizoval vliv ve smyslu počtu aktivních uzlů na konci. Označme $\sigma(B)$ očekávaný
počet uzlů ovlivněných množinou $B$; pak hledáme
\[ \max_{|B|\le k}\ \sigma(B). \]
Úloha je \textbf{NP-těžká}, lze však dokázat, že \emph{hladový} přístup dává řešení na úrovni
\textbf{63~\%} optima.

\textbf{Hladový přístup} začne s~$B=\emptyset$, vyhodnotí $\sigma(\{v\})$ pro každý uzel
a~vybere uzel s~maximálním $\sigma$ jako první prvek $B$. Poté opakovaně vybírá uzel, který
$\sigma(B)$ zvýší nejvíce, tj.\ $\arg\max_{v\in V\setminus B}\sigma(B\cup\{v\})$.
\end{defbox}

\subsection{Vliv versus korelace}

Atributy a~chování uživatelů bývají po síti rozprostřeny nenáhodně: často korelují s~tím, kdo
je s~kým spojen. Uvažujme u~každého uzlu binární atribut, například zda je ten člověk kuřák.
Koreluje-li atribut se sítí, budou aktéři sdílet stejné hodnoty atributu pozitivně korelovaně
se sociálními spoji, tedy kuřáci budou interagovat spíše s~kuřáky.

Zbývá otázka, jak takovou korelaci z~dat vůbec poznat. Postup je přímočarý: spočítá se, kolik
hran mezi odlišnými hodnotami atributu by vzniklo pouhou náhodou, a~toto očekávání se porovná
se skutečným grafem.

\begin{defbox}[Test korelace]
Jestliže podíl hran spojujících uzly s~\emph{různými} hodnotami atributu je významně
\emph{nižší} než očekávaná pravděpodobnost, jde o~důkaz korelace.

\emph{Příklad:} předpokládejme nejprve, že spoje jsou na kuřáctví nezávislé. Podíl kuřáků
v~síti je $p=4/9$, podíl nekuřáků $1-p=5/9$. Jedna hrana pak spojuje dva kuřáky
s~pravděpodobností $p\cdot p$, dva nekuřáky s~pravděpodobností $(1-p)(1-p)$ a~kuřáka
s~nekuřákem s~pravděpodobností $2p(1-p)=2\cdot\frac49\cdot\frac59=49\,\%$. Ve skutečném grafu
je ovšem tento podíl jen $2/14=14\,\%<49\,\%$, takže síť \textbf{vykazuje určitý stupeň
korelace} vzhledem ke kuřáctví.
\end{defbox}

Nalezená korelace ale sama o~sobě neříká nic o~tom, proč vznikla. Vysvětlení jsou tři
a~rozdíl mezi nimi je podstatný.

\begin{defbox}[Tři sociální procesy vysvětlující korelaci]
\begin{itemize}
\item \textbf{Homofilie} je tendence propojovat se s~těmi, kdo s~námi sdílejí určitou
      podobnost; lidově řečeno \uv{vrána k~vráně sedá}.
\item \textbf{Zmatení} (\emph{confounding}) připomíná, že korelaci mohou vytvořit i~\emph{vnější}
      vlivy prostředí: \uv{jednotlivci žijící ve stejném městě se spíše stanou přáteli než
      náhodní jednotlivci}.
\item \textbf{Vliv} (\emph{influence}) je naopak proces, který korelaci chování mezi sousedními
      aktéry teprve způsobuje: \uv{přejde-li většina mých přátel k~jinému mobilnímu operátorovi,
      mohu být jimi ovlivněn a~přejít také}.
\end{itemize}
\textbf{Proč vliv rozlišovat od homofilie a~zmatení?} Ze dvou důvodů. Při \emph{analýze} nám
rozlišení dovolí předpovědět dynamiku systému, tedy zda se nová norma chování, technologie nebo
idea může šířit jako epidemie. Při \emph{návrhu} umožňuje zkonstruovat systém, který konkrétní
chování vyvolá; sem patří vakcinační strategie (náhodné, cílené na demografickou skupinu,
náhodní známí) i~virální marketingové kampaně.
\end{defbox}

Vliv se dá od pouhé korelace oddělit i~kvantitativně, máme-li k~dispozici časové značky
jednotlivých aktivací.

\begin{defbox}[Logistický model vlivu a~shuffle test]
Ve studiích modelování vlivu se vliv určuje \emph{časovými značkami}. Pro hranu $(u,v)$ může
$u$ ovlivnit $v$ v~diskrétních krocích $t=0,1,2,\dots$, přičemž v~každém kroku se každý dosud
neaktivní uzel stane aktivním s~pravděpodobností $p(a)$, kde $a$ je počet aktivních kontaktů.
Pravděpodobnost aktivace se modeluje jako \emph{logistická} funkce počtu aktivních přátel:
\[ p(a)=\frac{e^{\alpha\ln(a+1)+\beta}}{1+e^{\alpha\ln(a+1)+\beta}},\qquad\text{ekvivalentně}\quad
   \ln\frac{p(a)}{1-p(a)}=\alpha\ln(a+1)+\beta, \]
kde $a$ je počet aktivních přátel, $\alpha$ \emph{koeficient sociální korelace}, $\beta$
konstanta vysvětlující vrozenou tendenci k~aktivaci a~$\ln(a+1)$ vysvětlující proměnná.

\textbf{Věrohodnost aktivace.} Nechť v~čase $t$ se $y_{a,t}$ uživatelů s~$a$ aktivními přáteli
stane aktivními a~$n_{a,t}$ uživatelů s~$a$ aktivními přáteli zůstane neaktivních. Věrohodnost
v~čase $t$ je pak $\prod_a p(a)^{y_{a,t}}\big(1-p(a)\big)^{n_{a,t}}$ a~pro celou posloupnost
kroků $\prod_t\prod_a p(a)^{y_{a,t}}\big(1-p(a)\big)^{n_{a,t}}$. Z~logu aktivit uživatelů se
tedy dají spočítat hodnoty $\alpha$ a~$\beta$, které tuto věrohodnost maximalizují, a~právě
\textbf{$\alpha$ měří korelaci mezi uzly}.

\textbf{Shuffle test (test na vliv).} Klíčová myšlenka testu je následující: \emph{nehraje-li
vliv roli, musí být časování aktivace nezávislé na časování ostatních aktérů}. Jinými slovy,
i~když časové značky aktivit náhodně zpřeházíme, měli bychom dostat podobnou sociální korelaci
jako předtím. Samotný \emph{test} z~toho plyne bezprostředně: zpřeházej časové značky aktivit
uživatelů, a~je-li nový odhad sociální korelace \emph{významně odlišný} od odhadu z~původního
logu, jde o~důkaz \textbf{vlivu}.

Postup má pět kroků. \emph{(1)}~Pro všechna $a$ spočítej $y_{a,t}$ a~$n_{a,t}$.
\emph{(2)}~Spočítej $\alpha,\beta$ maximalizující věrohodnost. \emph{(3)}~Vyber náhodnou
permutaci $\pi$ časů aktivace a~nastav čas aktivace uzlu $v_i$ na $t'_i=t_{\pi(i)}$; z~nových
časů spočítej $y'_{a,t}$, $n'_{a,t}$. \emph{(4)}~Spočítej $\alpha',\beta'$ maximalizující novou
věrohodnost. \emph{(5)}~Jsou-li $\alpha$ a~$\alpha'$ blízko sebe, model \emph{nevykazuje}
sociální vliv.
\end{defbox}

\subsection{Homofilie, výběr a~sociální vliv}

Z~trojice vysvětlení se nyní podrobněji podíváme na homofilii, protože právě u~ní se dá dobře
ukázat, jak se testuje a~jak souvisí s~tím, zda vazby formují lidi, nebo lidé vazby.

\begin{defbox}[Test homofilie]
Motivace: \uv{vrána k~vráně sedá} --- vaši přátelé jsou vám podobnější věkem, rasou, zájmy
a~názory než náhodně vybraná skupina jednotlivců. \emph{Homofilie} je princip, že máme tendenci
být podobní svým přátelům.

\textbf{Test:} kdyby byl každému uzlu \emph{náhodně} přiřazen atribut podle poměru v~reálné
síti, počet \uv{různorodých} hran by se od reálné sítě významně nelišil. Je-li podíl $p$ mužů
a~$q$ žen, je pravděpodobnost hrany muž--muž $p^2$, žena--žena $q^2$ a~\emph{různopohlavní}
hrany $2pq$.

\textbf{Je-li podíl heterogenních (různopohlavních) hran významně nižší než $2pq$, jde
o~důkaz homofilie.}

\emph{Příklad:} různopohlavních hran je 5 z~18; $p=6/9=2/3$, $q=3/9=1/3$; bez homofilie by
jich mělo být $2pq=4/9$, tj.\ 8 z~18 $\Rightarrow$ \emph{důkaz homofilie}.
\end{defbox}

\textbf{Poznámky k~testu homofilie.} Formulace \uv{významně nižší} si zaslouží upřesnění:
posuzuje se odchylka pod střední hodnotou. Co ale dělat, má-li síť \emph{významně více} než
$2pq$ různorodých hran? Pak nejde o~selhání testu, nýbrž o~\emph{inverzní homofilii}, jejímž
učebnicovým příkladem jsou partnerské vztahy muž--žena.

Test navíc není vázán na pohlaví. Použít se dá na libovolnou charakteristiku, ať už jde o~rasu,
věk, rodný jazyk, nebo preference. Má-li charakteristika více než dvě hodnoty, počítá se
analogicky: počet heterogenních hran se porovná s~tím, jak by vypadal náhodně generovaný graf,
v~němž jako pravděpodobnosti vystupují podíly z~reálných dat.

\begin{defbox}[Výběr vs.\ sociální vliv]
\textbf{Tranzitivní uzávěr} (\emph{triadic closure}) říká, že mají-li dva jednotlivci
společného přítele, je přátelství mezi nimi pravděpodobnější. Homofilie k~tomu dodává, že dva
jednotlivci jsou si podobnější \emph{kvůli} společnému příteli, takže vazba mezi nimi může
vzniknout, i~když ani jeden o~společném příteli neví. Obvykle je proto obtížné přisoudit vznik
vazby jedinému faktoru.

\textbf{Výběr} (\emph{selection}) je tendence vybírat si přátele, kteří jsou nám podobní.
Působí na různých úrovních intencionality: buď si v~malé skupině podobné přátele aktivně
vybíráte, \emph{nebo} je populace vaší školy homogenní vůči celkové populaci, takže vás
prostředí \emph{nutí} vybírat si podobné přátele, aniž byste to jakkoli zamýšleli.

\textbf{Mutabilní vs.\ imutabilní charakteristiky:} \emph{imutabilní} charakteristiky se
nemění (pohlaví, rasa) nebo se mění konzistentně s~populací (věk, generace), zatímco
\emph{mutabilní} se v~čase mění (chování, přesvědčení, zájmy, názory).

\textbf{Sociální vliv} působí \emph{opačně} než výběr. U~\emph{výběru} řídí individuální
charakteristiky tvorbu vazeb; u~\emph{sociálního vlivu} naopak existující vazby formují
mutabilní charakteristiky lidí.
\end{defbox}

Který z~obou směrů v~konkrétních datech převládá, z~jediného snímku sítě rozhodnout nelze,
a~proto jsou nutné \textbf{longitudinální studie} (studie s~dlouhým časovým rozpětím). Ty
sledují sociální spoje i~individuální chování v~čase a~odpovídají na otázku, zda \emph{se mění
chování po změnách v~síti, nebo se mění síť po změnách v~chování}. Použily se například ke
zjištění, zda měl na školní výsledky a~užívání drog u~teenagerů větší vliv výběr, nebo sociální
vliv.

Odpověď má přímý dopad na návrh intervencí. Vykazuje-li užívání drog v~síti homofilii, může
nejlépe fungovat program cílený na \emph{sociální vliv}, kdy přátelé přesvědčí ostatní přestat.
Vznikla-li však homofilie \emph{výběrem}, bývalí uživatelé si jednoduše vyberou nové přátele
a~chování ostatních to silně neovlivní.

Nejznámější studií tohoto typu je práce, v~níž \emph{Christakis a~Fowler} sledovali obezitu
a~sociální síť 12\,000 lidí po 32 let a~našli homofilii podle obezity. Nabízela se tři
vysvětlení: \emph{efekty výběru} (vybírají si lidé přátele s~podobným stavem?), \emph{zmatení}
(korelují se stavem obezity i~jiné faktory?) a~\emph{sociální vliv} (změnil-li se stav přátel,
ovlivnilo to budoucí stav dané osoby?). \textbf{Nalezen byl významný důkaz pro třetí hypotézu:}
obezita je typ \uv{nákazy}, který se může šířit sociálním vlivem.

\subsection{Afiliační sítě a~uzávěrové procesy}

Dosud jsme v~síti viděli jen lidi a~vazby mezi nimi. Reálné vazby ale vznikají v~nějakém
kontextu, a~ten se dá do sítě zanést jako druhý typ uzlu.

\begin{defbox}[Afiliační a~sociálně-afiliační síť]
\textbf{Afiliační síť} vkládá do sítě kontext tím, že zobrazuje spojení k~aktivitám,
společnostem, organizacím, čtvrtím atd. Je to \emph{bipartitní graf}, v~němž každá hrana
spojuje dva uzly patřící do \emph{různých} množin: osoby na jedné straně a~\emph{afiliace}
čili \emph{fokusy} na straně druhé.

\textbf{Sociálně-afiliační síť} postupuje o~krok dál a~má \emph{dva typy hran}. Hrana
\emph{osoba--osoba} vyjadřuje přátelství nebo jiný sociální vztah, hrana \emph{osoba--fokus}
účast v~daném fokusu.
\end{defbox}

Nad takovou sítí lze pojmenovat tři různé mechanismy, jimiž vznikají nové hrany.

\begin{defbox}[Tři uzávěrové procesy]
\begin{itemize}
\item \textbf{Tranzitivní uzávěr} (\emph{triadic closure}) nastane, když se dva lidé spojí,
      protože mají společného přítele.
\item \textbf{Fokální uzávěr} (\emph{focal closure}) nastane, když se dva lidé spojí, protože
      mají společný fokus; jde tedy o~uzávěr \emph{výběrem}.
\item \textbf{Členský uzávěr} (\emph{membership closure}) nastane, když se osoba zapojí do
      fokusu, protože je v~něm zapojen její přítel; jde o~uzávěr \emph{sociálním vlivem}.
\end{itemize}
\end{defbox}

\textbf{Otázky o~uzávěru.} Zajímá nás, zda by Sue byla pravděpodobněji přítelkyní Alice, kdyby
měly více než jednoho společného přítele. Formálněji se ptáme, \emph{jaká je pravděpodobnost,
že dva lidé vytvoří vazbu, jako funkce počtu společných přátel}. Zcela analogicky se ptáme
u~fokálního uzávěru, kde je argumentem počet společných fokusů, a~u~členského uzávěru, kde je
jím počet přátel už v~daném fokusu zapojených.

\begin{defbox}[Metodika měření tranzitivního uzávěru]
\begin{enumerate}
\item Pořiď dva snímky sítě v~časech $t_1$ a~$t_2$.
\item Pro každé $k$ identifikuj všechny \emph{páry} uzlů, které mají v~$t_1$ přesně $k$
      společných přátel, ale \emph{nejsou} přímo spojeny hranou.
\item Definuj $T(k)$ jako \emph{podíl} těchto párů, které do $t_2$ vytvořily hranu; to je
      hledaná pravděpodobnost, že mezi dvěma lidmi s~$k$ společnými přáteli vznikne vazba.
\item Vynes $T(k)$ jako funkci $k$.
\end{enumerate}

\textbf{Bázový (nulový) model:} je-li pravděpodobnost vytvoření vazby s~jedním společným
přítelem $p$ a~jednotlivé společné přátele považujeme za \emph{nezávislé} příležitosti, pak
\[ T_{\text{base}}(k)=1-(1-p)^{k}\qquad\text{(resp.\ }1-(1-p)^{k-1}\text{ podle konvence).} \]
\end{defbox}

\textbf{Empirické výsledky.} První měření provedli na \emph{e-mailové síti} Kossinets a~Watts
(2006). Sledovali e-mailovou komunikaci 22\,000 studentů velké americké univerzity po jeden
rok; vazbu mezi dvěma lidmi předpokládali tehdy, byl-li mezi nimi za posledních 60 dní odeslán
e-mail. Snímky pořizovali den po dni a~$T(k)$ zprůměrovali přes více párů snímků. Výsledek má
tři části: při \emph{žádném} společném příteli se e-maily téměř nevyměňují, při přechodu z~1 na
2 společné přátele nastává \emph{významný nárůst} a~každý další nárůst je sice stále významný,
ale týká se výrazně menší subpopulace. Uzávěr tedy vykazuje \emph{klesající výnosy}.

Tíž autoři změřili i~\emph{fokální uzávěr}. Z~rozvrhů týchž 22\,000 studentů vytvořili
sociálně-afiliační síť, v~níž fokusy jsou vyučované předměty. Ukázalo se, že \emph{sdílení
předmětu má téměř stejný efekt jako sdílení přítele}, a~to opět s~klesajícími výnosy.

Pro \emph{členský uzávěr} sestrojili Backstrom et al.\ (2006) sociálně-afiliační síť
LiveJournalu, kde přátelství čerpali z~profilů a~fokusy odpovídají členství v~uživatelsky
definovaných komunitách. I~zde nastal významný nárůst od 1 ke 2 a~poté klesající výnosy. Totéž
zopakovali Crandall et al.\ (2008) pro Wikipedii, kde uzly jsou editoři s~diskusní stránkou,
vazba znamená, že spolu přes diskusní stránku komunikovali, a~fokusy jsou editované články;
vzorec vyšel stejný.

\textbf{Výběr a~sociální vliv na datech Wikipedie.} Podobnost dvou editorů se měřila jako
\[ \frac{\text{počet článků editovaných oběma editory}}
   {\text{počet článků editovaných alespoň jedním editorem}} \]
Jde o~obdobu překryvu sousedství pro hrany. Zjištění bylo jednoznačné: \emph{páry editorů,
které komunikovaly, jsou si v~chování významně podobnější než ty, které nikdy nekomunikovaly.}
Otevřená ovšem zůstává otázka, zda homofilie vzniká proto, že editoři mluví s~editory téhož
článku (\emph{výběr}), nebo proto, že je k~článkům dovedou ti, s~nimiž mluví (\emph{sociální
vliv}). Analýza proto zprůměrovala historie interakcí, přičemž 0 označuje okamžik vzniku vazby,
započítali se jen editoři s~alespoň 100 akcemi před i~po první interakci a~baseline tvoří
podobnost náhodného páru editorů. Podobnost roste \emph{už před} vznikem vazby, což svědčí pro
výběr, a~pokračuje v~růstu \emph{i~po} něm, což svědčí pro vliv.

\subsection{Schellingův model segregace}

Homofilie působí lokálně, na úrovni jednotlivce a~jeho sousedů. Následující model ukazuje, jak
překvapivě silné globální důsledky z~ní mohou vzejít.

\begin{defbox}[Schellingův model (T.~Schelling, 70.~léta)]
Jednoduchý model, který ukazuje, jak mohou \emph{globální} vzorce dobrovolné segregace vznikat
z~\emph{lokální} homofilie.
\begin{itemize}
\item Populaci tvoří dva typy jednotlivců, agenti $X$ a~$O$; oba typy reprezentují
      \emph{imutabilní} charakteristiku (rasa, etnicita, země původu, rodný jazyk).
\item Agenti jsou náhodně rozmístěni v~mřížce.
\item Agent je se svou polohou \emph{spokojen}, je-li alespoň $t$ jeho sousedů agentem téhož
      typu, kde $t$ je předem nastavený prah.
\item V~každém kole se nejprve identifikují všichni \emph{nespokojení} agenti a~poté se všichni
      nespokojení přesunou na místo, kde budou spokojeni.
\end{itemize}
Někdy se stane, že agent žádné nové uspokojující místo nenajde; pak buď zůstane, kde je, nebo
se přesune na náhodné místo. Přesuny navíc mohou způsobit, že \emph{dříve spokojený agent
přestane být spokojen}, a~právě proto proces pokračuje dál.

\textbf{Výsledek: mřížka je segregovanější než na začátku.}

\textbf{Varianty modelu:} agenti se mohou přesouvat v~náhodném pořadí nebo \uv{zametacím}
pohybem, mohou mířit na nejbližší uspokojující místo nebo na náhodné, prah může být procentní
nebo se lišit mezi skupinami i~v~rámci nich a~svět se může \uv{obtočit}, takže levý okraj
navazuje na pravý. Variant existuje ještě celá řada, ale \emph{výsledek zůstává stejný:
samovolně vzniklá segregace.}

\textbf{Závěr:} Schellingův model je příkladem, jak se \emph{imutabilní} charakteristiky mohou
stát silně korelovanými s~\emph{mutabilními}. Volba místa k~životu je mutabilní a~v~čase se
přizpůsobí typu agenta, který je imutabilní, čímž vzniká segregace, tedy homofilie.
\end{defbox}

\subsection{Pozitivní a~negativní vztahy: strukturní balance}

Dosud měl každý vztah v~síti pozitivní konotaci. Existují ale sítě, v~nichž lze vyjádřit
i~vztahy negativní, tedy antagonistické či nepřátelské. \emph{Pozitivní} vztahy jsou přátelské,
\emph{negativní} nepřátelské, a~přirozeně vyvstane otázka, \textbf{kdy je síť v~balancovaném
stavu.}

Uvidíme přitom podobný jev jako u~Schellingova modelu: lokální vlastnosti vynucují vlastnosti
globální. Pro začátek si úlohu zjednodušíme a~budeme uvažovat sítě, v~nichž jsou \emph{všechny}
páry uzlů spojeny pozitivním nebo negativním vztahem, tedy úplné grafy čili kliky. Takový
předpoklad se hodí pro malé skupiny lidí nebo pro státy a~později jej uvolníme.

\begin{defbox}[Vlastnost strukturní balance]
Pro tři uzly existují (až na symetrii) čtyři konfigurace znamének:
\begin{center}
\small
\begin{tabular}{@{}p{3.0cm}p{9.5cm}p{2.4cm}@{}}
\toprule
\textbf{Znaménka} & \textbf{Interpretace} & \textbf{Stav} \\
\midrule
$+,+,+$ & $A$, $B$, $C$ jsou vzájemní přátelé & \emph{balancovaný} \\
$+,-,-$ & $A$ a~$B$ jsou přátelé se společným nepřítelem $C$ & \emph{balancovaný} \\
$+,+,-$ & $B$ a~$C$ jsou přátelé $A$, ale spolu nevycházejí & nebalancovaný \\
$-,-,-$ & $A$, $B$, $C$ jsou vzájemní nepřátelé & nebalancovaný \\
\bottomrule
\end{tabular}
\end{center}
\textbf{Vlastnost strukturní balance:} pro každý trojúhelník jsou všechny tři hrany pozitivní,
nebo je pozitivní \emph{právě jedna}.
\end{defbox}

Tato čistě lokální podmínka, formulovaná pro jednotlivé trojúhelníky, má nečekaně silný
globální důsledek.

\begin{defbox}[Věta o~balanci]
\textbf{Věta.} Je-li označkovaný \emph{úplný} graf balancovaný, pak buď
\begin{itemize}
\item jsou všechny páry uzlů přátelé, \emph{nebo}
\item lze uzly rozdělit do dvou skupin $X$ a~$Y$ tak, že každý pár uzlů v~$X$ se má rád,
      každý pár v~$Y$ se má rád, a~každý v~$X$ je nepřítelem každého v~$Y$.
\end{itemize}

\textbf{Důkaz.} Nechť $G$ je balancovaná síť. Jsou-li všechny hrany pozitivní, jsme hotovi.
Existuje-li alespoň jedna negativní hrana, vyber uzel $A$ a~polož
$X=$ všichni přátelé $A$ (včetně $A$), $Y=$ všichni nepřátelé $A$. Zbývá ověřit tři věci.
\emph{(1)}~Každé dva uzly v~$X$ jsou přátelé, protože jinak bychom měli trojúhelník
$A$ + dva jeho přátelé s~jedinou negativní hranou, tedy dvě kladné a~jednu zápornou, což je
nebalancované. \emph{(2)}~Každé dva uzly v~$Y$ jsou přátelé, protože jinak bychom měli
trojúhelník $A$ + dva jeho nepřátelé se třemi negativními hranami, což je rovněž
nebalancované. \emph{(3)}~Každý uzel v~$X$ je nepřítelem každého uzlu v~$Y$, jinak by
trojúhelník $A$, přítel z~$X$, nepřítel z~$Y$ měl dvě kladné a~jednu zápornou hranu.
$\square$
\end{defbox}

\textbf{Aplikace strukturní balance.} Klasickou ilustrací je \emph{vývoj aliancí v~Evropě
1872--1907}, kde plné tmavé hrany značí přátelství a~tečkované červené nepřátelství. Sledujeme
Ligu tří císařů 1872--81, Trojspolek 1882, přerušení německo-ruské smlouvy 1890,
francouzsko-ruskou alianci 1891--94, Entente Cordiale 1904 a~britsko-ruskou alianci 1907. Síť
při tom postupně \emph{sklouzává k~balancovanému označkování}, a~tím do první světové války.

Druhým příkladem je \emph{konflikt o~odtržení Bangladéše od Pákistánu (1972)}. Poněkud
překvapivá podpora Pákistánu ze strany USA je mnohem méně překvapivá, uvědomíme-li si, že SSSR
byl nepřítelem Číny, Čína nepřítelem Indie a~Indie měla tradičně špatné vztahy s~Pákistánem.
USA tehdy zlepšovaly vztahy s~Čínou, a~tak podpořily nepřátele nepřátel Číny.

Předpoklad úplného grafu i~samotná definice balance jsou pro reálné sítě příliš přísné, takže
se obojí dá oslabit.

\begin{defbox}[Slabá strukturní balance a~generalizace]
\textbf{Vlastnost slabé strukturní balance:} neexistuje trojice uzlů, jejíž hrany tvoří
\emph{právě dvě pozitivní a~jednu negativní}. Konfigurace \uv{tři negativní} je tedy nově
dovolena, protože odpovídá situaci, kdy pár vytvoří alianci proti třetímu, zatímco \uv{dvě
pozitivní a~jedna negativní} je silnější nebalancovanost, neboť $B$ a~$C$ se snaží svou
animozitu vyřešit.

\textbf{Věta.} Je-li označkovaný úplný graf slabě balancovaný, lze jeho uzly rozdělit do
\emph{skupin vzájemných přátel} tak, že každé dva uzly patřící do různých skupin jsou
nepřátelé.

\textbf{Důkaz.} \emph{(1)}~Nechť $A$ je uzel a~$X$ množina přátel $A$ včetně $A$.
\emph{(2)}~Všichni přátelé $A$ jsou přátelé mezi sebou. \emph{(3)}~$A$ a~všichni jeho přátelé
jsou nepřáteli všech ostatních. \emph{(4)}~Odstraň $X$ ze sítě a~pokračuj krokem 1. $\square$

\textbf{Generalizace na neúplné grafy.} \emph{(a)}~Nejsou-li všechny páry v~síti spojeny,
prohlásíme (neúplný) graf za \emph{balancovaný} tehdy, lze-li jej \uv{doplnit} přidáním hran
na označkovaný úplný graf, který je balancovaný. \emph{(b)}~Nemusí-li být balancované všechny
trojúhelníky, stačí, aby jich byla balancovaná \emph{většina}; pak lze síť \emph{přibližně}
rozdělit na dvě frakce.
\end{defbox}

\subsection{Analýza vazeb: huby, autority a~algoritmus HITS}

Rok 1998 byl pro modely analýzy vazeb na webu zásadní, protože v~něm byly publikovány
\emph{oba} algoritmy, PageRank i~HITS. Souvislosti mezi nimi jsou nápadné, přesto se
dominantním modelem stal PageRank, a~to díky své \emph{nezávislosti na dotazu}, schopnosti
bránit se spamu a~obchodnímu úspěchu Googlu.

\begin{defbox}[Autority a~huby]
\textbf{Autorita} (\emph{authority}) je hrubě řečeno stránka s~mnoha vstupními odkazy. Myšlenka
za tím je jednoduchá: stránka může mít k~nějakému tématu dobrý či autoritativní obsah, a~proto
jí mnoho lidí věří a~odkazuje na ni.

\textbf{Hub} je naproti tomu stránka s~mnoha \emph{výstupními} odkazy. Slouží jako organizátor
informací k~danému tématu a~odkazuje na mnoho dobrých autorit.

\textbf{Klíčová myšlenka HITS} obě role svazuje dohromady: \emph{dobrý hub odkazuje na mnoho
dobrých autorit} a~\emph{dobrá autorita je odkazována mnoha dobrými huby.} Autority a~huby jsou
tedy ve vztahu \emph{vzájemného zesilování}, který se v~grafu projeví jako bipartitní podgraf
hustě propojených hubů a~autorit.

\textbf{Proč huby oddělovat od autorit?} Protože konkurenční firmy na sebe neodkazují, takže je
nelze chápat tak, že se přímo vzájemně podporují. V~PageRanku naopak podpora přechází přímo
z~jedné prominentní stránky na druhou, neboť stránka je důležitá, je-li citována jinými
důležitými stránkami. To je dominantní způsob podpory u~akademických a~vládních stránek, mezi
blogery, u~osobních stránek a~ve vědecké literatuře.
\end{defbox}

\begin{defbox}[Algoritmus HITS (\emph{Hypertext Induced Topic Search})]
Zatímco PageRank je \emph{statický}, HITS je \emph{závislý na dotazu}. Po zadání dotazu proto
nejprve rozšíří seznam relevantních stránek vrácených vyhledávačem a~teprve pak nad rozšířenou
množinou vytvoří \emph{dvě} pořadí: pořadí autorit a~pořadí hubů.

\textbf{Sběr stránek.} Pro obecný dotaz $q$ se postupuje takto: pošli $q$ vyhledávači a~posbírej
$t$ nejvýše hodnocených stránek (v~původním článku $t=200$); této množině se říká
\emph{kořenová množina} $S$. Poté $S$ rozšiř o~každou stránku, na kterou odkazuje nějaká
stránka z~$S$, a~o~každou stránku, která odkazuje na stránku z~$S$; vznikne tak větší
\emph{bázová množina} $B$.

\textbf{Skóre.} Nechť $G=(V,E)$ je hypertextový graf $B$ o~$n$ stránkách s~maticí sousednosti
$L$, kde $L_{ij}=1$, existuje-li hrana $(i,j)$, a~0 jinak. Označme $a(i)$ skóre autority
a~$h(i)$ skóre hubu stránky $i$. Vzájemné zesilování se pak zapíše jako
\[ a(i)=\sum_{(j,i)\in E}h(j),\qquad h(i)=\sum_{(i,j)\in E}a(j), \]
maticově, s~$a=(a(1),\dots,a(n))^{\!\top}$ a~$h=(h(1),\dots,h(n))^{\!\top}$:
\[ a=L^{\!\top}h,\qquad h=La. \]

\textbf{Výpočet mocninnou iterací:}
\begin{algorithmic}[1]
\State $a_0\gets h_0\gets (1,1,\dots,1)^{\!\top}$; $k\gets 1$
\Repeat
  \State $a_k\gets L^{\!\top}h_{k-1}$; \quad $h_k\gets La_k$
  \State $a_k\gets a_k/\|a_k\|_1$; \quad $h_k\gets h_k/\|h_k\|_1$ \Comment{normalizace}
  \State $k\gets k+1$
\Until{$\|a_k-a_{k-1}\|_1<\varepsilon$ \textbf{a} $\|h_k-h_{k-1}\|_1<\varepsilon$}
\State \Return $a$ a~$h$
\end{algorithmic}
\textbf{Normalizace je nezbytná}, protože velikost hodnot autorit a~hubů s~každou aktualizací
roste; konvergence nastane jen s~normalizací, kdy se dělí součtem všech skóre.
\end{defbox}

Že iterace ke~správným hodnotám opravdu konverguje, se ukáže standardní úvahou o~vlastních
číslech.

\begin{defbox}[Spektrální analýza HITS]
Po $k$ krocích platí (spektrum matice $=$ množina jejích vlastních čísel)
\[ a_k=\big(L^{\!\top}L\big)^{k-1}a_0,\qquad h_k=\big(LL^{\!\top}\big)^{k}h_0. \]
Vektor $a_k$ tedy musí konvergovat k~\emph{vlastnímu vektoru} matice $L^{\!\top}L$. Existují
totiž normalizační konstanty $c_k$ takové, že $a_k/c_k$ konverguje k~limitě, a~směr se při
násobení maticí nezmění právě tehdy, je-li vektor vlastním vektorem.

Matice $L^{\!\top}L$ je \emph{symetrická}, má tedy $n$ vlastních vektorů $z_1,\dots,z_n$
s~vlastními čísly $c_1,\dots,c_n$, kde $|c_1|\ge|c_2|\ge\cdots\ge|c_n|$. Rozvineme-li
$a_0=x_1z_1+\cdots+x_nz_n$, dostaneme
\[ a_k=c_1^{k}x_1z_1+\cdots+c_n^{k}x_nz_n,\qquad
   \frac{a_k}{c_1^{k}}=x_1z_1+x_2\Big(\frac{c_2}{c_1}\Big)^{\!k}z_2+\cdots \]
Za předpokladu $c_1>c_2\ge\cdots$ jde pro $k\to\infty$ každý člen kromě prvního k~nule, takže
$a_k/c_1^k$ konverguje k~$x_1z_1$, tedy k~\emph{hlavnímu} vlastnímu vektoru. (Konvergenci
nezávisle na počátečním vektoru a~uvolnění předpokladu $c_1>c_2$ zde nedokazujeme.)
\end{defbox}

Obě matice, které se v~HITS objevily, mají v~bibliometrii dlouho zavedený význam.

\begin{defbox}[Vztah ke ko-citaci a~bibliografickému párování]
\textbf{Ko-citace} stránek $i$ a~$j$ udává, kolik článků cituje \emph{oba} články $i$ a~$j$:
\[ C_{ij}=\sum_{k=1}^{n}L_{ki}L_{kj}=\big(L^{\!\top}L\big)_{ij}. \]
Matice autorit $L^{\!\top}L$ algoritmu HITS je tedy \emph{ko-citační matice}.

\textbf{Bibliografické párování} stránek $i$ a~$j$ naproti tomu udává, kolik článků je citováno
\emph{oběma} články $i$ a~$j$:
\[ B_{ij}=\sum_{k=1}^{n}L_{ik}L_{jk}=\big(LL^{\!\top}\big)_{ij}. \]
Matice hubů $LL^{\!\top}$ je tedy matice bibliografického párování.
\end{defbox}

\textbf{Silné a~slabé stránky HITS.} \emph{Silnou stránkou} je schopnost řadit stránky
\emph{podle tématu dotazu}, což může poskytnout relevantnější autority a~huby. \emph{Slabé
stránky} jsou tři. \emph{(1)}~HITS lze snadno spamovat, protože přidat výstupní odkazy na
vlastní stránku je velmi snadné. \emph{(2)}~Nastává \emph{drift tématu} (\emph{topic drift}),
neboť mnoho stránek v~rozšířené množině nemusí být k~tématu. \emph{(3)}~Algoritmus je
neefektivní v~době dotazu, jelikož sběr kořenové množiny, její rozšíření i~výpočet vlastních
vektorů jsou drahé operace.

\subsection{PageRank}
\label{ssec:pagerank}

Než dojdeme k~PageRanku samotnému, hodí se pojmenovat myšlenku, na které stojí: hlasy nemají
stejnou váhu.

\begin{defbox}[Prestiž pořadí (\emph{rank prestige})]
U~měr prestiže se uvažuje důležitý faktor, totiž \emph{prominence jednotlivých aktérů, kteří
\uv{hlasují}}. V~reálném světě je osoba vybraná důležitou osobou prestižnější než osoba vybraná
méně důležitou, protože hlas generálního ředitele váží více než hlas zaměstnance. Je-li váš
okruh vlivu plný prestižních aktérů, je i~vaše vlastní prestiž vysoká.

Prestiž pořadí $P(i)$ je proto definována jako lineární kombinace pořadí těch, kteří odkazují
na $i$:
\[ P(i)=A_{1i}P(1)+A_{2i}P(2)+\cdots+A_{ni}P(n),\qquad\text{maticově}\quad P=A^{\!\top}P, \]
kde $A_{ji}=1$, odkazuje-li $j$ na $i$, a~0 jinak. To je \emph{charakteristická rovnice}
vlastního systému matice $A^{\!\top}$, takže $P$ je vlastní vektor $A^{\!\top}$.
\end{defbox}

\begin{defbox}[Základní PageRank]
PageRank interpretuje hyperodkaz ze stránky $u$ na stránku $v$ jako \emph{hlas} stránky $u$ pro
stránku $v$. Nezajímá jej ovšem pouhý počet hlasů: analyzuje i~stránku, která hlas dává, neboť
\emph{hlasy důležitých stránek váží více a~pomáhají činit ostatní stránky důležitějšími}.
A~protože stránka může odkazovat na mnoho jiných, musí své skóre \emph{rozdělit}.

Web chápeme jako orientovaný graf $G=(V,E)$ s~$n=|V|$. PageRank stránky $i$ je pak
\[ P(i)=\sum_{(j,i)\in E}\frac{P(j)}{O_j}, \]
kde $O_j$ je počet výstupních odkazů $j$. Maticově: nechť $A_{ij}=1/O_i$ pro $(i,j)\in E$
a~0 jinak; pak $P=A^{\!\top}P$.

\textbf{Intuice:} PageRank se chová jako \uv{kapalina}, která cirkuluje sítí, přechází od uzlu
k~uzlu po hranách a~shromažďuje se v~nejdůležitějších uzlech. \emph{Celkový PageRank v~síti
zůstává konstantní}, protože každá stránka svůj PageRank rozdělí a~pošle po odkazech dál.
PageRank se nikdy nevytváří ani nezaniká, jen se přemísťuje, a~proto na rozdíl od HITS
\emph{není nutná normalizace}.
\end{defbox}

\textbf{Problém základní definice.} V~mnoha sítích skončí veškerý PageRank ve \uv{špatných}
uzlech. Odkazují-li si $F$ a~$G$ navzájem místo na $A$, PageRank, který k~nim přiteče z~$C$, se
už nikdy nemůže vrátit; pro velká $k$ dostaneme $\frac12$ pro každý z~$F$, $G$ a~0 pro všechny
ostatní. Tomuto jevu se říká \uv{pomalý únik} (\emph{slow leak}) a~způsobují jej malé množiny
uzlů dosažitelné z~ostatního grafu, ale bez cest zpět. \textbf{Řešení:} posílat malý zlomek
PageRanku \emph{všem} uzlům.

\begin{defbox}[Škálovaný PageRank a~tlumicí faktor]
Zvol \emph{škálovací faktor} $d$ mezi 0 a~1 a~nahraď základní pravidlo aktualizace takto:
\begin{enumerate}
\item nejprve aplikuj základní pravidlo;
\item všechny hodnoty PageRanku sniž faktorem $d$, čímž se celkový PageRank v~síti zmenší
      z~1 na $d$;
\item zbývajících $1-d$ jednotek rozděl \emph{rovnoměrně} mezi všechny uzly, tj.\ $(1-d)/n$
      každému.
\end{enumerate}
\textbf{Proč to funguje:} je to \uv{vodní cyklus}, který v~každém kroku $1-d$ jednotek
PageRanku \uv{odpaří} a~rovnoměrně je \uv{sprchne} na všechny uzly.

Opakovaná aplikace škálovaného pravidla \emph{konverguje} k~limitním hodnotám PageRanku
a~ekvilibrium je \emph{jednoznačné}; hodnoty ovšem závisejí na volbě $d$. V~praxi se používá
$d$ mezi 0,8 a~0,9 ($d=0{,}85$ v~původním článku) a~$d$ se nazývá \emph{tlumicí faktor}
(\emph{damping factor}).
\end{defbox}

Celý mechanismus se dá popsat i~jazykem teorie pravděpodobnosti, což zároveň vysvětlí, odkud se
berou podmínky konvergence.

\begin{defbox}[PageRank jako markovský řetězec]
Web se modeluje jako \emph{stochastický proces}, v~němž každá stránka je \emph{stav} a~každý
hyperodkaz \emph{přechod} s~určitou pravděpodobností. Při \emph{náhodném surfování} je každá
přechodová pravděpodobnost $1/O_i$, předpokládáme-li, že surfař klikne na odkazy uniformně
náhodně a~nepoužívá tlačítko \uv{zpět} ani nezadává URL.

Nechť $A$ je matice přechodových pravděpodobností a~$p_0$ počáteční rozdělení
pravděpodobnosti. Pak $p_1=A^{\!\top}p_0$ a~obecně $p_k=A^{\!\top}p_{k-1}$. \emph{Věta
o~markovských řetězcích} nyní říká, že konečný markovský řetězec definovaný stochastickou
maticí $A$ má \textbf{jednoznačné stacionární rozdělení}, je-li $A$ \emph{ireducibilní}
a~\emph{aperiodická}. Stacionární rozdělení znamená, že $p_k$ konverguje ke stálému vektoru
$\pi$ nezávisle na volbě $p_0$. V~ustáleném stavu platí $\pi=A^{\!\top}\pi$, takže $\pi$ je
hlavní vlastní vektor $A^{\!\top}$ s~vlastním číslem 1, a~právě $\pi$ se použije jako vektor
PageRanku. Je to rozumné a~intuitivní, protože $\pi$ odráží \emph{dlouhodobé pravděpodobnosti},
že náhodný surfař stránku navštíví, a~stránka má vysokou prestiž, je-li pravděpodobnost její
návštěvy vysoká.

\textbf{Ani jedna podmínka však pro webový graf neplatí:}
\begin{itemize}
\item \emph{$A$ není stochastická}, protože mnoho stránek nemá výstupní odkazy, což se
      v~matici projeví řádky samých nul (\emph{visící stránky}, \emph{dangling pages}). Nabízejí
      se dvě řešení: \emph{(1)}~takové stránky během výpočtu odstranit, neboť neovlivňují přímo
      pořadí žádné jiné stránky; \emph{(2)}~přidat z~každé takové stránky úplnou množinu odkazů
      na všechny stránky, tj.\ položit $A_{ij}=1/n$.
\item \emph{$A$ není ireducibilní}, neboť ireducibilita znamená, že webový graf je \emph{silně
      souvislý}, tedy že pro každý pár $i,j$ existuje cesta z~$i$ do $j$; obecný webový graf
      silně souvislý není.
\item \emph{$A$ není aperiodická.} Stav $i$ je \emph{periodický} s~periodou $k>1$, je-li $k$
      nejmenší číslo takové, že všechny cesty ze stavu $i$ zpět do $i$ mají délku, která je
      násobkem $k$; není-li stav periodický ($k=1$), je aperiodický, a~řetězec je aperiodický,
      jsou-li aperiodické všechny stavy. Příkladem periodického řetězce s~$k=3$ je situace, kdy
      se z~1 lze vrátit do 1 jen cestou $1\to2\to3\to1$, takže každý návrat trvá 3 přechody.
\end{itemize}
\textbf{Obě posledně jmenované vlastnosti se řeší jedinou strategií:} přidá se z~každé stránky
odkaz na každou stránku a~každému takovému odkazu se dá malá přechodová pravděpodobnost řízená
parametrem $d$. Takto rozšířená přechodová matice je zjevně ireducibilní i~aperiodická.
\end{defbox}

\begin{defbox}[Finální PageRank]
Po augmentaci má náhodný surfař na stránce dvě možnosti: s~pravděpodobností $d$ zvolí náhodně
odkaz, který bude následovat, a~s~pravděpodobností $1-d$ skočí na náhodnou stránku. Odtud
\[ P=\Big(\frac{1-d}{n}\Big)\,\mathbf{E}+d\,A^{\!\top}P, \]
kde $\mathbf E=\mathbf{ee}^{\!\top}$ a~$\mathbf e$ je sloupcový vektor jedniček, takže
$\mathbf E$ je matice $n\times n$ samých jedniček. Matice
$\big(\frac{1-d}{n}\big)\mathbf E+dA^{\!\top}$ je stochastická (transponovaná), ireducibilní
i~aperiodická. Přeškálujeme-li ji tak, aby $\mathbf e^{\!\top}P=n$, dostaneme
\[ P=(1-d)\,\mathbf e+d\,A^{\!\top}P,\qquad\text{tj.\ pro každou stránku } i: \quad
   P(i)=(1-d)+d\sum_{j=1}^{n}A_{ji}P(j), \]
což je ekvivalentní formuli z~původního článku o~PageRanku
\[ P(i)=(1-d)+d\sum_{(j,i)\in E}\frac{P(j)}{O_j}. \]

\textbf{Výpočet mocninnou iterací:}
\begin{algorithmic}[1]
\State $P_0\gets \mathbf e/n$; $k\gets 1$
\Repeat
  \State $P_k\gets (1-d)\mathbf e+dA^{\!\top}P_{k-1}$
  \State $k\gets k+1$
\Until{$\|P_k-P_{k-1}\|_1<\varepsilon$}
\State \Return $P_k$
\end{algorithmic}
\end{defbox}

\textbf{Výhody PageRanku.} První výhodou je \emph{boj se spamem}: stránka je důležitá, jsou-li
důležité stránky, které na ni odkazují, a~protože pro majitele stránky není snadné přidat
\emph{vstupní} odkazy z~jiných důležitých stránek, není snadné PageRank ovlivnit. Druhou
výhodou je \emph{globální a~dotazově nezávislá míra}: hodnoty PageRanku všech stránek se
počítají a~ukládají \emph{offline}, nikoli v~době dotazu. \textbf{Kritika} míří
právě na tuto dotazovou nezávislost, protože PageRank neumí rozlišit stránky autoritativní
\emph{obecně} od stránek autoritativních \emph{k~tématu dotazu}.

\textbf{Souhrn.} Analýza vazeb nám ukázala principy \emph{centrality} a~\emph{prestiže}, dvojici
měr \emph{ko-citace} a~\emph{bibliografického párování} a~algoritmy \emph{HITS} a~\emph{PageRank},
z~nichž druhý stojí za Googlem. Yahoo! a~MSN mají vlastní nepublikované algoritmy založené na
odkazech. \textbf{Důležité je nezapomenout}, že hyperodkazové řazení není jediný algoritmus
používaný ve vyhledávačích: kombinuje se s~mnoha faktory založenými na obsahu, aby vzniklo
finální pořadí prezentované uživateli.

\subsection{Tematické komunity a~jejich objevování}

Odkazy se dají využít i~k~něčemu jinému než k~hodnocení důležitosti jednotlivých stránek. Lze
z~nich vyčíst \emph{tematické komunity}, tedy skupiny tvůrců obsahu nebo lidí, kteří sdílejí
nějaký společný zájem. Podle toho, v~jakých datech se hledají, se mluví o~webových komunitách,
o~e-mailových komunitách nebo o~komunitách pojmenovaných entit.

Blízko k~tomu má \emph{cílené procházení} (\emph{focused crawling}), které kombinuje obsah
stránek a~odkazy mezi nimi, aby prošlo jen tu část webu, jež se týká konkrétního tématu.
Prohledávač následuje odkazy a~u~každé navštívené stránky pomocí učení a~klasifikace rozhodne,
zda k~tématu patří.

\begin{defbox}[Tematická komunita]
Mějme danou konečnou množinu entit $E=\{e_1,\dots,e_n\}$ \emph{téhož typu}. \emph{Tematická
komunita} je pár $K=(T,M)$, kde $T$ je \emph{téma} komunity a~$M\subseteq E$ množina všech entit
z~$E$, které téma $T$ sdílejí. Je-li $e\in M$, říkáme, že $e$ je \emph{členem} komunity $K$.
Z~takto postavené definice plyne několik vlastností:
\begin{itemize}
\item \emph{Téma definuje komunitu.} Pro dané téma $T$ je množina členů určena jednoznačně,
      takže dvě komunity jsou si rovny, mají-li totéž téma.
\item \emph{Téma lze definovat libovolně.} Může jím být událost, třeba sportovní událost nebo
      skandál, stejně dobře ale i~celý koncept, například dolování z~webu.
\item \emph{Entita může být v~libovolném počtu komunit.} Komunity se proto mohou
      \emph{překrývat} a~více komunit může sdílet členy.
\item \emph{Entity v~$E$ jsou téhož typu.} Lidé a~organizace tedy nemohou být v~téže komunitě.
\end{itemize}
Komunity mohou tvořit \emph{hierarchickou strukturu}: tematická komunita $(T,M)$ může obsahovat
množinu tematických \emph{subkomunit} $\{(T_1,M_1),\dots,(T_k,M_k)\}$, kde $T_i$ je podtéma
$T$ a~$M_i\subseteq M$. Původní komunita $(T,M)$ pak vystupuje jako \emph{superkomunita}
a~každou subkomunitu lze dekomponovat dál.

\textbf{Cílem objevování} je najít v~datové sadě obsahující entity skryté komunity, a~to
u~každé z~nich jak její \emph{téma} (obvykle reprezentované množinou klíčových slov), tak její
\emph{členy}.
\end{defbox}

\textbf{Jak se komunita projeví v~datech.} Členové komunity jsou spolu nějak propojeni
a~přidružené texty obsahují slova indikativní pro téma komunity. Konkrétní podoba obojího se
ale liší podle typu dat. Na \emph{webových stránkách} jsou členové komunity spíše vzájemně
propojeni \emph{hyperodkazy} a~jejich stránky obsahují \emph{slova} vztažená k~tématu.
V~\emph{e-mailech} spolu členové spíše komunikují a~slova vztažená k~tématu nesou samotné
zprávy. V~\emph{textových dokumentech} se členové spíše objevují \emph{spolu} v~téže větě či
v~témže dokumentu a~téma indikují slova v~těchto větách.

\begin{defbox}[Bipartitní jádra komunit]
$(i,j)$-\emph{bipartitní jádro} je úplný bipartitní podgraf s~alespoň $i$ \emph{vějířovými}
uzly (\emph{fans}) a~alespoň $j$ \emph{středovými} uzly (\emph{centers}), obsahuje tedy všechny
možné hrany mezi $i$ vějířovými a~$j$ středovými uzly. Obě strany jádra mají svůj protějšek
ve slovníku hubů a~autorit: vějíře reprezentují \emph{specializované huby}, zatímco středy
odpovídají \emph{autoritám}.

\textbf{Postup hledání bipartitních jader} má dva kroky:
\begin{itemize}
\item \emph{Krok 1: prořezání.} Nejprve se prořezává podle vstupního stupně, tedy odstraní se
      všechny stránky, které jsou velmi silně odkazované (počet vstupních odkazů $>k$, např.\
      $k=50$). Potom se iterativně prořezávají vějíře a~středy, aby zbyli kandidáti na
      $(i,j)$-jádra: odstraní se všechny vějíře s~výstupním stupněm menším než $j$, odstraní se
      všechny středy se vstupním stupněm menším než $i$ a~z~grafu se odstraní hrany spojené
      s~odstraněnými uzly.
\item \emph{Krok 2: generování všech $(i,j)$-jader.} Nejdřív se najdou všechna $(1,j)$-jádra,
      tj.\ množina všech uzlů s~výstupním stupněm alespoň $j$. Všechna $(2,j)$-jádra se pak
      zkonstruují kontrolou každého vějíře, který odkazuje také na některý střed
      v~$(1,j)$-jádru, a~stejným způsobem se pokračuje až do $i$.
\end{itemize}
\textbf{Omezení.} Algoritmus najde jen \emph{jádrové} stránky komunit, ne všechny členské
stránky. Nenajde ani \emph{témata} komunit, ani jejich hierarchickou strukturu.
\end{defbox}

\begin{defbox}[Překrývající se komunity pojmenovaných entit]
\textbf{Cílem} je najít komunity entit z~textového korpusu. \textbf{Hlavní idea} spočívá
v~tom, že dvě entity jsou propojeny, \emph{vyskytnou-li se v~téže větě}: jsou-li dva lidé
zmíněni v~jedné větě, obvykle mezi nimi nějaký vztah existuje. Jedna entita se přitom může
objevit ve více komunitách a~téma komunity reprezentují nejvýznamnější klíčová slova. Metoda
přinesla slibné výsledky při hledání komunit politických osobností a~celebrit z~webových
dokumentů.

\textbf{Algoritmus} probíhá ve čtyřech krocích:
\begin{enumerate}
\item \emph{Postav graf vazeb.} Každý dokument se parsuje a~v~každé větě se identifikují
      obsažené pojmenované entity. Má-li věta více než jednu entitu, propojí se tyto entity
      \emph{po párech} a~klíčová slova z~věty se k~propojeným párům připojí jako textový obsah.
      Ostatní věty se zahodí.
\item \emph{Najdi všechny trojúhelníky.} Trojúhelník tvoří tři entity svázané dohromady
      a~reprezentuje základní stavební blok komunit. Vychází se z~pozorování, že komunita má
      tendenci se rozšiřovat trojúhelníky sdílejícími společnou hranu.
\item \emph{Najdi jádra komunit.} Jádro je skupina těsně svázaných trojúhelníků, tedy
      \emph{uvolněný} úplný podgraf (klika). Intuitivně jde o~množinu těsně propojených členů
      komunity.
\item \emph{Klastruj okolo jader.} Trojúhelníky a~páry entit, které neleží v~žádném jádru, se
      přiřadí k~jádrům podle \emph{podobnosti textového obsahu} s~nalezenými jádry.
\end{enumerate}
\end{defbox}

\subsection{Detekce komunit}

Předchozí metody hledaly komunity i~s~jejich tématy, tedy současně v~obsahu a~ve struktuře.
Detekce komunit v~užším smyslu si klade skromnější otázku: jak se síť rozpadá na skupiny, když
se díváme jen na její topologii.

\begin{defbox}[Komunita a~detekce komunit]
\emph{Komunity} jsou tvořeny jednotlivci tak, že ti \emph{uvnitř} skupiny interagují mezi sebou
\emph{častěji} než s~těmi vně skupiny. V~různých kontextech se jim říká také \emph{skupina},
\emph{klastr}, \emph{kohezní podgrupa} nebo \emph{modul}.

\emph{Detekce komunit} je objevování skupin v~síti, kde příslušnost jednotlivců ke skupinám
\emph{není explicitně dána}.

\textbf{Definice komunity je subjektivní.} Komunitou může být hustě propojený podgraf, stejně
tak jí ale může být každá komponenta souvislosti.

\textbf{V~sociálních médiích se rozlišují dva typy skupin:} \emph{explicitní} skupiny tvořené
předplacením uživatelů a~\emph{implicitní} skupiny tvořené sociálními interakcemi.

Nabízí se otázka, proč skupiny vůbec extrahovat z~topologie, když některé weby umožňují do
skupin přímo vstoupit. Důvody jsou tři: ne všechny weby platformu pro komunity poskytují, ne
všichni lidé jsou ochotni vynaložit úsilí na vstup do skupiny a~skupiny se mohou \emph{dynamicky
měnit}. Interakce v~síti navíc samy o~sobě poskytují informaci o~vztazích mezi uživateli.
Doplňují ostatní druhy informací, pomáhají vizualizaci a~navigaci a~poskytují základ pro další
úlohy.
\end{defbox}

\begin{defbox}[Metody minimálního řezu]
Metoda \textbf{minimálního řezu} (Ahuja et al., 1993) vychází z~prosté úvahy: většina interakcí
nastává \emph{uvnitř} skupiny, kdežto mezi skupinami jsou interakce vzácné $\Rightarrow$
detekce komunit $\approx$ úloha minimálního řezu.

\emph{Řez} je rozdělení uzlů grafu do dvou disjunktních množin. \emph{Úloha minimálního
řezu} hledá takové rozdělení, které minimalizuje počet hran mezi oběma množinami, u~vážených
sítí součet jejich vah; tomuto počtu se říká \emph{velikost řezu}.

Chceme-li $k>2$ skupin, implementuje se strategie \emph{iterativního bisekce}: ve druhé
a~dalších iteracích se už nalezené skupiny postupně dělí na dvě podskupiny.

\textbf{Nevýhodou je, že metoda produkuje skupiny nevyvážených velikostí} a~jedna z~množin
často bývá jednoprvková.

\textbf{Složitost.} Jednoduché řešení pomocí $s$-$t$ řezu založeného na maximálním toku má
složitost $O(n^5)$; implementace Kargerovým algoritmem klesne na $O(n^4)$, ale nezaručuje vždy
nalezení minimálního řezu. Pro jednotkové váhy $w_{ij}=1$ a~bez balančních omezení je úloha
2-cestného řezu polynomiálně řešitelná v~$O(n^2\log n)$ (Kargerův randomizovaný minimální řez).
Jakmile se sejdou obecné váhy hran \emph{a} balanční omezení, stává se úloha \textbf{NP-těžkou}.
\end{defbox}

\begin{defbox}[Poměrový a~normalizovaný řez]
Minimální řez často vrací nevyvážené rozdělení, kde je jedna množina jednoprvková $\Rightarrow$
změníme účelovou funkci tak, aby brala v~úvahu \emph{velikost komunity}. Pro rozdělení
$\pi=\{C_1,\dots,C_k\}$ se definují dvě míry:
\[ \mathrm{RatioCut}(\pi)=\sum_{i=1}^{k}\frac{\mathrm{cut}(C_i,\bar C_i)}{|C_i|},\qquad
   \mathrm{NCut}(\pi)=\sum_{i=1}^{k}\frac{\mathrm{cut}(C_i,\bar C_i)}{\mathrm{vol}(C_i)}, \]
kde $|C_i|$ je počet uzlů v~$C_i$, $\mathrm{vol}(C_i)$ součet stupňů v~$C_i$
a~$\mathrm{cut}(A,B)=\sum_{i\in A,\,j\in B}a_{ij}$. \textbf{Obě míry preferují vyvážené
rozdělení.}
\end{defbox}

\begin{defbox}[Girvanové--Newmanův algoritmus (2002)]
Jde o~\emph{divizivní hierarchický} algoritmus, tedy o~přístup shora dolů: dekomponuje počáteční
úplný graf na postupně menší souvislé části, dokud nezůstanou žádné hrany k~odstranění a~každý
uzel nereprezentuje sám komunitu.

\textbf{Idea} je identifikovat hrany spojující uzly z~\emph{různých} komunit, tzv.\ mosty,
a~iterativně je z~grafu odstraňovat. Slouží k~tomu \textbf{kritérium} \emph{mezilehlosti hrany}
(\emph{edge betweenness}), což je podíl nejkratších cest procházejících danou hranou. Hrany
s~vysokou mezilehlostí leží na velkém počtu nejkratších cest mezi uzly, a~jsou tedy mosty.

\begin{algorithmic}[1]
\State spočítej mezilehlost všech hran v~síti
\State odstraň hranu s~\emph{nejvyšší} mezilehlostí
\State opakuj předchozí kroky, dokud v~grafu zbývají hrany k~odstranění
\end{algorithmic}
\textbf{Vstupem} je celý graf, \textbf{výstupem} hierarchická struktura, reprezentovaná např.\
dendrogramem.

\textbf{Volba počtu komunit.} Algoritmus vrátí \emph{množinu} možných řešení, ale neurčí, které
z~nich je nejlepší. Nejlepší rozdělení se proto vybere tak, že se pro každé dělení sítě spočítá
\emph{modularita} a~vezme se rozdělení s~nejvyšší modularitou.

\textbf{Nevýhodou} jsou vysoké výpočetní náklady, kvůli nimž je algoritmus vhodný pouze pro sítě
\emph{střední velikosti}: $O(m^2n)$ pro graf o~$m$ hranách a~$n$ uzlech, resp.\ $O(n^3)$ pro
řídké sítě.
\end{defbox}

\emph{Příklad výpočtu mezilehlosti hrany:} mezilehlost hrany $e(1,2)$ je 4, protože všechny
nejkratší cesty z~2 do $\{4,5,6,7,8,9\}$ musí procházet buď $e(1,2)$, nebo $e(2,3)$, a~$e(1,2)$
je nejkratší cesta mezi 1 a~2. Po odstranění $e(4,5)$ se mezilehlost $e(4,6)$ zvýší na 20, což
je nejvyšší hodnota; po odstranění $e(4,6)$ má nejvyšší mezilehlost 4 hrana $e(7,9)$.

\begin{defbox}[Modularita $Q$]
Modularita kvantifikuje kvalitu daného rozdělení grafu na komunity. \textbf{Základní idea} je
ta, že síť má smysluplnou komunitní strukturu, je-li počet hran \emph{mezi} komunitami menší,
než by se očekávalo při náhodné volbě. Očekávaný počet hran mezi uzly $i$ a~$j$ se stupni $k_i$,
$k_j$ je $k_ik_j/2m$, odkud
\[ Q=\frac{1}{2m}\sum_{i,j}\Big(A_{ij}-\frac{k_ik_j}{2m}\Big)\delta(c_i,c_j), \]
kde $A_{ij}$ je prvek matice sousednosti (počet hran mezi $i$ a~$j$), $m$ počet hran, $k_i$
stupeň uzlu $i$, $c_i$ skupina, do níž $i$ patří, a~$\delta$ Kroneckerovo delta.

\textbf{Hodnoty.} Modularita nabývá hodnot $Q\in[-1,1]$ a~může být pozitivní i~negativní.
Je-li pozitivní, existuje možnost nalézt v~síti komunitní strukturu; jsou-li hodnoty velké
a~pozitivní, může odpovídající rozdělení odrážet skutečnou komunitní strukturu. Pro smysluplné
komunity by mělo být $Q\ge 0{,}3$ (Clauset et al., 2004).
\end{defbox}

\begin{defbox}[Louvainova metoda (Blondel et al., 2008)]
Louvainova metoda je \emph{hladová} optimalizační metoda pro \emph{aglomerativní hierarchickou}
maximalizaci modularity a~běží ve dvou fázích.

Ve \textbf{fázi 1} se každý uzel považuje za samostatnou komunitu a~postupuje se takto:
\begin{enumerate}
\item Spočítej modularitu $Q$.
\item Přesuň izolovaný uzel z~jeho komunity do sousední komunity.
\item Spočítej zisk či ztrátu modularity, kterou tato změna přiřazení způsobila.
\item Zvýší-li se modularita, tedy jde-li o~zisk, ponech uzel v~\uv{nové} komunitě.
\item Opakuj, dokud není možné žádné další zlepšení.
\end{enumerate}
\emph{Výstupem} fáze 1 je rozdělení $k$-té úrovně, kde $k$ je číslo iterace.

Ve \textbf{fázi 2} se vytvoří nová síť (\emph{supergraf}) odvozená z~původní. Každý nový uzel
(\emph{superuzel}) je agregací uzlů přiřazených dané komunitě nalezené ve fázi 1 a~dva
superuzly jsou spojeny hranou, existuje-li alespoň jedna hrana spojující dva uzly uvnitř
odpovídajících komunit. Síť se tím zmenší a~lze na ni znovu pustit fázi 1.

\textbf{Iterace.} Obě fáze se střídají tak dlouho, dokud není dosaženo maximální modularity.

\textbf{Výhodou} je dobrý výkon i~na velkých síťových grafech při nízkých výpočetních nákladech
$O(n\log n)$, kde $n$ je počet uzlů. \textbf{Nevýhodou} je, že výsledky jsou \emph{citlivé
na pořadí} zpracování.

\textbf{Další nevýhody.} Metoda může produkovat \emph{nesouvislé komunity}: přesune-li se uzel
do jiné komunity, může se původní komunita stát \emph{vnitřně nesouvislou}, a~přesto zůstanou
její uzly lokálně optimálně přiřazeny, takže v~ní zůstanou. Metoda dále \emph{uvázne v~lokálních
optimech} a~\emph{selhává v~upřesnění komunit} po počátečním slučování.
\end{defbox}

\begin{defbox}[Leidenův algoritmus (Traag, Waltman, van Eck, 2019)]
Leidenův algoritmus je upřesněním Louvainovy metody pro aglomerativní hierarchickou maximalizaci
modularity. Probíhá ve třech fázích: \emph{(1)}~lokální přesuny uzlů, obdobně jako u~Louvaina;
\emph{(2)}~\textbf{upřesnění rozdělení}, což je právě ten krok, který Louvainovi chybí;
\emph{(3)}~agregace sítě, opět obdobně jako u~Louvaina.

\textbf{Výhodou} je vyšší efektivita než u~Louvainovy metody. U~obrovských grafů ale i~tak může
docházet k~dlouhým dobám zpracování a~rychlost pak mohou zvýšit paralelní vícejádrové
implementace. \textbf{Nevýhodou} je \emph{tvrdé rozdělení}, kdy uzly mohou patřit jen do jedné
komunity.
\end{defbox}

\textbf{Výzvy detekce komunit.} Přes velké množství algoritmů zůstává detekce komunit náročným
problémem, a~to ze dvou hlavních důvodů. \emph{(1)}~Počet komunit obvykle není znám a~je třeba
jej určit. \emph{(2)}~Komunity jsou obvykle \emph{heterogenní} co do velikosti a~hustoty.

\subsubsection{Taxonomie metod detekce komunit}

Metody lze rozdělit do čtyř kategorií podle toho, na jak velký kus sítě se jejich kritérium
dívá; samotná kritéria se přitom mohou podle úlohy lišit. \emph{Uzlově orientované} metody
žádají, aby určité vlastnosti splňoval každý uzel ve skupině. \emph{Skupinově orientované}
metody se dívají na spojení uvnitř skupiny jako celku a~nezoomují až na úroveň jednotlivých
uzlů. \emph{Síťově orientované} metody pracují s~celou sítí a~rozdělují ji do několika
disjunktních množin. A~konečně \emph{hierarchicky orientované} metody konstruují hierarchickou
strukturu komunit.

\begin{defbox}[Uzlově orientovaná detekce komunit]
Po uzlech se může požadovat, aby splňovaly různé vlastnosti používané v~tradiční SNA. Následující
čtyři rodiny definic se liší právě tím, kterou vlastnost si vyberou.

\textbf{Úplná vzájemnost: kliky.} \emph{Klika} je maximální úplný podgraf tří a~více uzlů,
z~nichž všechny jsou navzájem sousední. Nalezení maximální kliky je \textbf{NP-těžké}
a~přímočará implementace vychází časově velmi draho. Pomůže \emph{idea rekurzivního
prořezávání}: v~klice velikosti $k$ má každý uzel stupeň $\ge k-1$, takže uzly se stupněm
$<k-1$ v~maximální klice nebudou. Opakovaně se tedy vzorkuje podsíť, v~ní se najde klika (např.\
hladově), a~je-li velikosti $k$, odstraní se pro nalezení větší kliky všechny uzly se stupněm
$\le k-1$; postup se opakuje, dokud není síť dost malá. V~sociálních médiích se takto prořeže
mnoho uzlů, protože stupně uzlů následují mocninné rozdělení. Kliky se běžně používají jako
\emph{jádro} či \emph{semeno} pro prozkoumání větších komunit.

\textbf{Metoda perkolace klik} (\emph{clique percolation method}, CPM) reaguje na to, že klika je
striktní definice a~v~praxi nestabilní. Zároveň je to příklad metody, která umí najít
\emph{překrývající se} komunity.

\emph{Vstupem} je parametr $k$ a~síť. \emph{Postup} začne nalezením všech klik velikosti $k$,
z~nichž se zkonstruuje \emph{graf klik}: dvě kliky jsou v~něm sousední, sdílejí-li $k-1$ uzlů.
Každá komponenta souvislosti grafu klik pak tvoří jednu komunitu.
\emph{Příklad:} kliky velikosti 3 jsou $\{1,2,3\}$, $\{1,3,4\}$, $\{4,5,6\}$,
$\{5,6,7\}$, $\{5,6,8\}$, $\{5,7,8\}$, $\{6,7,8\}$ $\Rightarrow$ komunity $\{1,2,3,4\}$
a~$\{4,5,6,7,8\}$, přičemž uzel 4 je v~obou.

\textbf{Dosažitelnost: $k$-klika, $k$-klub, $k$-klan.} Hlavní ideou je, že každý uzel ve skupině
má být dosažitelný v~$k$ skocích. \emph{$k$-klika} je maximální podgraf, v~němž je největší
geodetická vzdálenost mezi libovolnými dvěma uzly $\le k$. Pozor na to, že $k$-klika může mít
\emph{uvnitř podgrafu} průměr $\ge k$: v~jiném příkladovém grafu ze slidů má 2-klika
$\{12,4,10,1,6\}$ $d(1,6)=3$. \emph{$k$-klub} je podstruktura s~průměrem $\le k$
a~\emph{$k$-klan} je $k$-klika s~průměrem $\le k$. \emph{Příklad:} kliky $\{1,2,3\}$; 2-kliky
$\{1,2,3,4,5\}$, $\{2,3,4,5,6\}$; 2-kluby $\{1,2,3,4\}$, $\{1,2,3,5\}$, $\{2,3,4,5,6\}$.

\textbf{Stupně uzlů: $k$-jádro, $k$-plex.} Zde má mít každý uzel určitý počet spojení na uzly
\emph{uvnitř} skupiny. \emph{$k$-jádro} ($k$-core) je podstruktura, kde je každý uzel spojen
s~alespoň $k$ dalšími členy skupiny. U~\emph{$k$-plexu} má být ve skupině o~$n_s$ uzlech každý
uzel sousední s~ne méně než $n_s-k$ uzly ve skupině. Obě definice jsou komplementární, protože
$k$-jádro je $(n_s-k)$-plex. Množina všech $k$-jader pro dané $k$ však \emph{není} identická
s~množinou všech $(n_s-k)$-plexů, protože velikost skupiny $n_s$ se mezi $k$-jádry může lišit.

\textbf{Vazby uvnitř vs.\ vně: LS množiny.} \emph{LS množina} (Luccio--Sami) je definována tak,
že libovolná její vlastní podmnožina má více vazeb na ostatní uzly \emph{ve} skupině než
\emph{mimo} ni. Cílem bylo najít kohezní skupiny uzlů pro návrh obvodů na čipech; pro síťovou
analýzu je taková definice \emph{příliš striktní}. Jejím uvolněním je \emph{$k$-komponenta},
tedy maximální podgraf, který zůstává souvislý i~po odstranění méně než $k$ uzlů či hran.
Vyžaduje ovšem výpočet hranové souvislosti mezi libovolným párem uzlů algoritmem minimálního
řezu či maximálního toku.

\textbf{Omezení uzlově orientovaných definic.} Jsou příliš striktní, byť je lze použít jako
jádro komunity. Nejsou škálovatelné, takže se běžně používají pro sítě malé velikosti. A~někdy
nejsou konzistentní s~vlastnostmi rozsáhlých sítí, například se stupni uzlů v~bezškálových
sítích.
\end{defbox}

\begin{defbox}[Skupinově orientovaná detekce: kvazikliky]
Skupinově orientovaná kritéria vyžadují, aby určitou podmínku splňovala \emph{celá skupina},
například aby hustota skupiny byla dostatečně velká, tedy $\ge$ daný prah. Podgraf
$G_S=(V_S,E_S)$ je \emph{$\gamma$-hustá kvaziklika}, je-li
\[ |E_S|\ \ge\ \gamma\cdot\frac{|V_S|\,(|V_S|-1)}{2}. \]
Použít lze strategii podobnou té u~klik. \emph{Lokální prohledávání} vzorkuje podgraf a~najde
v~něm maximální $\gamma$-hustou kvazikliku velikosti $k$, \emph{heuristické prořezávání} pak
odstraní uzly se stupněm menším než průměrný stupeň, tedy $<\gamma(k-1)$.
\end{defbox}

\begin{defbox}[Síťově orientovaná detekce komunit]
Kritéria uvažují spojení v~síti \emph{globálně} a~cílem je rozdělit uzly sítě do
\emph{disjunktních} množin. Přístupů je pět: shlukování podle podobnosti uzlů, modely
latentního prostoru, aproximace blokovým modelem, spektrální shlukování a~maximalizace
modularity.

\textbf{Shlukování podle podobnosti uzlů.} Na uzly se aplikuje $k$-means nebo shlukování
založené na podobnosti, přičemž podobnost uzlů je definována podobností jejich
\emph{sousedství}. \emph{Strukturní ekvivalence} říká, že dva uzly jsou strukturně ekvivalentní
právě tehdy, jsou-li spojeny s~\emph{touž} množinou aktérů; pro praktické použití je to příliš
restriktivní. Prakticky se proto spojení chápou jako příznaky a~měří se
\[ \mathrm{Jaccard}(v_i,v_j)=\frac{|N_i\cap N_j|}{|N_i\cup N_j|},\qquad
   \cos(v_i,v_j)=\frac{|N_i\cap N_j|}{\sqrt{|N_i|\,|N_j|}}. \]
Pro rozsáhlé sítě se tedy spočítá podobnost uzlů kosinovou nebo Jaccardovou mírou a~aplikuje
se klasický $k$-means, v~němž je každý klastr asociován s~centroidem a~každý uzel je přiřazen ke
klastru s~nejbližším centroidem.

\textbf{Modely latentního prostoru.} Zobrazí uzly do \emph{nízkorozměrného} prostoru tak, aby
byla zachována blízkost uzlů podle propojenosti sítě, a~teprve pak aplikují $k$-means.
\emph{Vícerozměrné škálování} (MDS) pro danou síť zkonstruuje matici blízkosti $P$
reprezentující párové vzdálenosti uzlů, například geodetické. Nechť $S\in\mathbb R^{n\times l}$
jsou souřadnice uzlů v~nízkorozměrném prostoru; řešením je $S=V\Lambda^{1/2}$, kde $V$ je $l$
hlavních vlastních vektorů (a~$\Lambda$ diagonální matice odpovídajících vlastních čísel)
dvojitě centrované matice blízkosti.

\textbf{Blokové modely.} Matice sousednosti se aproximuje jako $A\approx S\,\Sigma\,S^{\!\top}$,
kde $\Sigma$ je \emph{míchací matice} a~$S$ \emph{indikátorová matice komunit}. Uvolní-li se
$S$ na reálné hodnoty, odpovídá optimální řešení hlavním vlastním vektorům $A$.

\textbf{Spektrální shlukování.} Poměrový i~normalizovaný řez lze přeformulovat jako minimalizaci
$\mathrm{tr}(S^{\!\top}LS)$, kde $L=D-A$ je \emph{grafový laplacián} (pro poměrový řez), resp.\
$L=D^{-1/2}(D-A)D^{-1/2}$ \emph{normalizovaný} grafový laplacián (pro normalizovaný řez), a~$D$
je diagonální matice stupňů. Po \emph{spektrální relaxaci} je optimálním řešením matice tvořená
vlastními vektory s~\emph{nejmenšími} vlastními čísly. První vlastní vektor je přitom
nepoužitelný, protože říká, že všechny uzly patří do téhož klastru.

\textbf{Maximalizace modularity.} Modularita měří sílu rozdělení tím, že bere v~úvahu rozdělení
stupňů. \emph{Síla komunity} je $\sum_{i,j\in C}\big(A_{ij}-\frac{d_id_j}{2m}\big)$
a~modularita $Q=\frac{1}{2m}\sum_{C}\sum_{i,j\in C}\big(A_{ij}-\frac{d_id_j}{2m}\big)$.
Analogicky ke spektrálnímu shlukování lze maximalizaci modularity přeformulovat pomocí
\emph{matice modularity} $B=A-\frac{dd^{\!\top}}{2m}$. Optimálním řešením jsou pak \emph{hlavní}
vlastní vektory matice modularity, na které se jako postprocessing aplikuje $k$-means.

\textbf{Jednotný pohled: maticová faktorizace.} Modely latentního prostoru, blokové modely,
spektrální shlukování a~maximalizaci modularity lze všechny zapsat jako faktorizaci
$X\approx S\,\Sigma\,S^{\!\top}$, kde
\[ X=\begin{cases}
\text{(dvojitě centrovaná) matice blízkosti} & \text{modely latentního prostoru},\\
\text{sociomatice }A & \text{aproximace blokovým modelem},\\
\text{grafový laplacián} & \text{minimalizace řezu},\\
\text{matice modularity} & \text{maximalizace modularity}.
\end{cases} \]
Liší se tedy jen tím, co se za $X$ dosadí. \textbf{Omezení} mají všechny stejné: vyžadují, aby
uživatel zadal \emph{počet komunit} předem.
\end{defbox}

\begin{defbox}[Hierarchicky orientovaná detekce komunit]
\textbf{Cílem} je vybudovat hierarchickou strukturu komunit na základě topologie sítě, což
umožňuje analyzovat síť v~různých \emph{rozlišeních}. Reprezentativní přístupy jsou dva.

\textbf{Divizivní hierarchické shlukování}, jehož zástupcem je Girvanové--Newman ($O(m^2n)$,
resp.\ $O(n^3)$ pro řídké sítě), rozdělí uzly do několika množin a~každou množinu dále rozdělí
na menší; na rozdělení lze aplikovat síťově orientované metody. Jedním konkrétním příkladem je
rekurzivní odstraňování \uv{nejslabší} vazby: najde se hrana s~nejmenší silou, odstraní se
a~aktualizuje se síla každé hrany, a~to rekurzivně, dokud se síť nerozpadne na požadovaný počet
komponent souvislosti. Každá komponenta pak tvoří komunitu. Sílu vazby lze měřit
\emph{mezilehlostí hrany}, protože hrana s~vyšší mezilehlostí bývá mostem mezi dvěma komunitami.

\textbf{Aglomerativní hierarchické shlukování}, jehož zástupcem je Louvain s~$O(n\log n)$,
postupuje opačně: každý uzel se inicializuje jako komunita a~komunity se pak postupně slučují do
větších podle konkrétního kritéria, například na základě \emph{přírůstku modularity}.
\end{defbox}

\subsection{Predikce vazeb}

Detekce komunit se ptá, jak síť vypadá teď. Predikce vazeb se ptá, jak bude vypadat dál, a~opírá
se přitom o~dva principy: o~homofilii a~o~tranzitivní uzávěr.

\begin{defbox}[Úloha predikce vazeb]
\textbf{Cílem} je predikovat \emph{budoucí} vazby mezi páry uzlů v~síti. V~komerčních sociálních
sítích jako Facebook nebo LinkedIn se tak uživatelům doporučují potenciální přátelé. Opřít se
lze o~podobnost obsahu i~o~strukturu sítě:
\begin{itemize}
\item \textbf{Míry založené na obsahu} využívají princip \emph{homofilie}, tedy že uzly
      s~podobným obsahem se spíše propojí. Predikci vazeb zlepšují, jsou ale citlivé na
      konkrétní síť, jak ukazují krátké a~zašuměné tweety s~nestandardními akronymy na Twitteru.
      Platí přitom, že \emph{obsahová homofilie nutně neimplikuje strukturní propojenost!}
\item \textbf{Strukturní míry} využívají princip \emph{tranzitivního uzávěru}, tedy že dva uzly,
      které sdílejí sousedství, se spíše propojí. Tranzitivní uzávěr je velmi častý napříč
      doménami, a~proto jsou tyto míry účinné v~různých typech sítí a~přímo aplikovatelné.
      Patří k~nim \emph{sousedské} míry, \emph{Katzova} míra a~míry založené na \emph{náhodných
      procházkách}. Opačným směrem už implikace platí: \emph{strukturní propojenost obvykle
      implikuje i~obsahovou homofilii.}
\end{itemize}
\end{defbox}

\begin{defbox}[Sousedské míry predikce vazeb]
Sousedské míry kvantifikují pravděpodobnost budoucí vazby počtem \emph{společných sousedů} mezi
párem uzlů $i$ a~$j$. Nechť $S_i$ je množina sousedů uzlu $i$.

\textbf{Míra společných sousedů} je prostě $\mathrm{CN}(i,j)=|S_i\cap S_j|$. Její
\emph{slabinou} je, že nezohledňuje \emph{relativní} počet společných sousedů vůči počtu
ostatních spojení: kdyby byly Alice a~Bob velmi populární veřejné osobnosti spojené s~velkým
počtem aktérů, mohli by mít mnoho společných sousedů \emph{jen náhodou}.

\textbf{Jaccardova míra} $\mathrm{Jaccard}(i,j)=\dfrac{|S_i\cap S_j|}{|S_i\cup S_j|}$ měří
podobnost sousedství \emph{normalizovanou} vůči stupňům uzlů, které se kvůli mocninnému
rozdělení silně liší. Větší stupně Alice nebo Boba tedy jejich Jaccardův koeficient sníží.
\emph{Slabinou} je, že se míra dobře nepřizpůsobuje \emph{prostředním} sousedům: jsou-li
společní sousedé Alice a~Boba velmi populární, tedy mají vysoký stupeň, pak se vyskytují jako
společní sousedé i~u~mnoha jiných párů uzlů, a~jsou proto pro predikci \emph{méně důležití}.

\textbf{Adamicova--Adarova míra} je vážená verze míry společných sousedů, navržená tak, aby
zohlednila různou důležitost různých společných sousedů. Váha společného souseda je
\emph{klesající} funkcí jeho stupně, např.\ $1/\log|S_k|$:
\[ \mathrm{AA}(i,j)=\sum_{k\in S_i\cap S_j}\frac{1}{\log|S_k|}. \]
\emph{Příklad:} pro čtyři společné sousedy se stupni 4, 2, 2, 4 vyjde
$\frac{1}{\log 4}+\frac{1}{\log 2}+\frac{1}{\log 2}+\frac{1}{\log 4}=\frac{3}{\log 2}$.
\end{defbox}

\begin{defbox}[Katzova míra]
Sousedské míry \emph{nejsou účinné}, je-li počet sousedů sdílených párem uzlů malý. Sdílejí-li
Alice a~Bob jediného společného souseda a~Alice a~Jim také jediného, sousedské míry mezi těmito
případy nerozliší, přestože existuje \emph{nepřímá} propojenost přes delší cesty. V~takových
případech jsou vhodnější míry založené na \emph{procházkách}, jako je Katzova míra.

Nechť $W_{ij}^{(t)}$ je počet procházek délky $t$ mezi uzly $i$ a~$j$. Pro předem zvolený
parametr $\beta<1$ se definuje
\[ \mathrm{Katz}(i,j)=\sum_{t=1}^{\infty}\beta^{t}\,W_{ij}^{(t)}, \]
kde $\beta$ je \emph{diskontní faktor}, který potlačuje význam delších procházek; pro malá
$\beta$ nekonečná suma konverguje. Je-li $A$ symetrická matice sousednosti neorientované sítě,
lze matici Katzových koeficientů $K$ rozměru $n\times n$ spočítat jako
\[ K=\sum_{t=1}^{\infty}\beta^{t}A^{t}=(I-\beta A)^{-1}-I, \]
neboť vlastní čísla $A^t$ jsou $t$-té mocniny vlastních čísel $A$. Silně centrální uzly mají
tendenci tvořit vazby s~mnoha uzly $\Rightarrow$ součet Katzových koeficientů uzlu $i$ vůči
ostatním uzlům se nazývá jeho \textbf{Katzova centralita}.
\end{defbox}

\begin{defbox}[Predikce vazeb jako klasifikační úloha]
Přítomnost či absence vazby mezi párem uzlů se dá chápat jako \emph{binární} indikátor třídy.
Pro každý pár uzlů se pak zkonstruuje vícerozměrný datový záznam:
\begin{itemize}
\item \emph{příznaky} zahrnují všechny různé podobnosti mezi uzly, tedy sousedské, Katzovy
      i~založené na procházkách, plus řadu příznaků preferenčního připojování, například stupně
      obou uzlů v~páru;
\item výsledkem je \emph{pozitivně-neoznačkovaná} klasifikační úloha, protože páry uzlů
      s~hranou jsou pozitivní příklady a~zbývající páry jsou neoznačkované;
\item neoznačkované příklady lze při trénování přibližně považovat za negativní;
\item negativních párů je ale v~rozsáhlých a~řídkých sítích příliš mnoho, a~proto se použije jen
      jejich \emph{vzorek}.
\end{itemize}

\textbf{Algoritmus.} V~\emph{trénovací fázi} se vygenerují data s~jedním záznamem pro každý pár
uzlů s~hranou a~se vzorkem záznamů z~párů bez hrany; příznaky odpovídají extrahovaným
podobnostním a~strukturním rysům a~označením třídy je přítomnost či absence hrany. Nad těmito
daty se postaví a~natrénuje model klasifikátoru. V~\emph{testovací fázi} se každý testovaný pár
uzlů převede na vícerozměrný záznam a~natrénovaný klasifikátor predikuje jeho označení. Běžnou
volbou bázového klasifikátoru je \emph{logistická regrese}; kvůli nevyváženosti dat se často
používají \emph{cenově citlivé} verze klasifikátorů.

\textbf{Výhody supervizovaného přístupu.} Obsahové příznaky lze extrahovat a~použít snadno
a~klasifikátor se \emph{sám} naučí jejich relevanci. Pro orientované sítě lze navíc extrahovat
příznaky i~\emph{asymetricky}, například vstupní a~výstupní stupně místo obyčejných stupňů nebo
různé míry založené na náhodných procházkách jako personalizovaný PageRank. Celkově přístup
nabízí vyšší flexibilitu díky schopnosti učit se vztahy mezi vazbami a~příznaky různých typů.
\end{defbox}

\textbf{Souhrn predikce vazeb.} Jednotlivé míry se v~účinnosti liší podle datové sady.
\emph{Sousedské} míry lze efektivně spočítat i~pro rozsáhlá data a~fungují téměř tak dobře jako
ostatní nesupervizované míry. Míry založené na \emph{náhodných procházkách} a~\emph{Katzovy}
míry jsou užitečné pro \emph{velmi řídké} sítě, v~nichž nelze počet společných sousedů robustně
změřit. \emph{Supervize} poskytuje lepší přesnost a~adaptabilitu napříč doménami sociálních
sítí, je ale výpočetně drahá. V~posledních letech se k~vylepšení predikce používá i~obsah, avšak
strukturní míry jsou \emph{výrazně silnější}, což lze přičíst tranzitivním vlastnostem reálných
sítí. Obsahové míry jsou založeny na \uv{obrácené homofilii}, kdy se pro predikci vazeb využívá
podobný nebo s~vazbami korelovaný obsah.

\subsection{Shrnutí ke~zkoušce}

\begin{itemize}
\item \textbf{SNA} vyrostla ze sociálních věd, z~analýzy sítí a~z~teorie grafů a~jejím
      předmětem jsou \emph{vztahy}, nikoli entity a~jejich atributy. Síť se dá zapsat obrázkem,
      seznamem hran nebo maticí sousednosti, která je u~orientovaných grafů nesymetrická.
      Rozlišuje se \emph{ego síť} a~\uv{celá} síť, jenže žádná studovaná síť ve skutečnosti
      celá není: tomu se říká \emph{problém specifikace hranic}.
\item \textbf{Síla vazeb} dělí hrany na silné a~slabé; klasickým výsledkem je \emph{síla slabých
      vazeb} (Granovetter 1973). Se silou vazeb souvisejí homofilie, tranzitivita a~přemostění,
      přičemž homofilie spolu s~tranzitivitou vedou ke~vzniku \emph{klik}. Most je hrana, jejíž
      odstranění rozpojí koncové uzly, a~o~\emph{zkratkový most} jde tehdy, když po jejím
      odstranění vzroste vzdálenost nad 2. Sílu vazby kvantifikuje \textbf{překryv sousedství}
      $O(i,j)=\frac{|N_i\cap N_j|}{|N_i\cup N_j|-2}$: čím větší, tím silnější vazba.
\item \textbf{Čtyři míry centrality} se ptají každá na něco jiného. \emph{Stupňová} říká, kolik
      lidí uzel přímo dosáhne. \emph{Mezilehlá} počítá, na kolika nejkratších cestách uzel leží,
      a~tím ukazuje, kde se síť rozpadne. \emph{Blízkostní} je převrácená průměrná délka
      nejkratších cest a~odpovídá rychlosti šíření. \emph{Vlastní} se počítá z~$Ax=\lambda x$,
      kde nezáporný vektor dává jen \emph{největší} vlastní číslo, a~měří, jak dobře je uzel
      spojen s~dobře propojenými. U~zkoušky se chce interpretační tabulka a~vysvětlení, proč
      jedna míra nestačí (uzly 3 a~5 \emph{společně} $>$ uzel 10). \textbf{Strukturní
      ekvivalence} znamená, že si dva uzly mohou vyměnit pozice, aniž se změní struktura.
\item \textbf{Kohezi} popisuje několik ukazatelů. Reciprocita má smysl jen u~orientovaných
      grafů. Hustota je $|E|/\binom{|V|}{2}$, klika má hodnotu 1 a~orientované grafy mají
      hustotu poloviční. \textbf{Klastrovací koeficient} uzlu je
      $C_i=\frac{2|\{e_{jk}\}|}{k_i(k_i-1)}$ a~koeficient celé sítě je jejich průměr.
      Vzdálenosti popisuje průměr sítě (\emph{diameter}), tedy nejdelší nejkratší cesta, dále
      průměrná vzdálenost a~excentricita. O~\textbf{malý svět} jde tam, kde je vysoké $C$
      a~zároveň krátká průměrná cesta.
\item Reálné komplexní sítě mají \textbf{šest vlastností}: efekt malého světa (Milgram naměřil
      medián 6), tranzitivitu neboli klastrování, mocninné rozdělení stupňů, odolnost (síť je
      robustní k~náhodnému, ale zranitelná cíleným odstraněním hubů), vzorce míchání
      (asortativní míchání odpovídá homofilii) a~komunitní strukturu. \textbf{Bezškálové sítě}
      vznikají růstem a~\emph{preferenčním připojováním}, kterému se také říká rich-get-richer,
      Matoušův efekt nebo kumulativní výhoda; vysvětluje se popularitou, kvalitou a~smíšeným
      modelem (halo efekt). Celkový tvar sítě pak doplňuje \textbf{centralizace} a~struktura
      \emph{jádro--periferie} (bow-tie analýza).
\item \textbf{Modelování vlivu} stojí na společných rysech: pracuje se s~orientovaným grafem,
      uzly jsou aktivní, nebo neaktivní, a~\emph{aktivovaný uzel nelze deaktivovat}. \textbf{LTM}
      losuje prah $\theta_v$ uniformně z~$[0,1]$ a~uzel aktivuje, jakmile
      $\sum_{\text{aktivní}}w_{u,v}\ge\theta_v\sum w_{u,v}$; soustředí se tedy na
      \emph{příjemce}. \textbf{ICM} naopak dává uzlu \emph{jedinou} příležitost aktivovat
      každého souseda, a~to s~pravděpodobností $p_{v,u}$; těžiště je zde na \emph{vysílajícím}.
      \textbf{Maximalizace vlivu} je NP-těžká, hladový přístup ale zaručuje \textbf{63~\%}
      optima.
\item Odlišit \textbf{vliv od korelace} dá několik kroků. \emph{Test korelace}
      porovná podíl hran mezi různými hodnotami atributu s~hodnotou $2p(1-p)$; u~kuřáků vyšlo
      $14\,\%<49\,\%$. Za korelací mohou stát tři procesy: homofilie, \emph{zmatení} a~vliv.
      Logistický model $\ln\frac{p(a)}{1-p(a)}=\alpha\ln(a+1)+\beta$ má v~$\alpha$ koeficient
      sociální korelace a~\textbf{shuffle test} zpřehází časové značky; liší-li se $\alpha'$
      významně od $\alpha$, je to důkaz \emph{vlivu}.
\item \textbf{Test homofilie} porovnává počet heterogenních hran s~hodnotou $2pq$. Je-li jich
      významně méně, síť vykazuje homofilii; je-li jich významně \emph{více}, jde
      o~\emph{inverzní} homofilii, jakou znají třeba partnerské vztahy. U~zkoušky se počítá
      příklad: 5 z~18 proti $2pq=8$ z~18.
\item \textbf{Výběr a~sociální vliv} působí opačným směrem: při výběru řídí charakteristiky
      tvorbu vazeb, kdežto vliv znamená, že vazby formují \emph{mutabilní} charakteristiky.
      \emph{Imutabilní} je pohlaví, rasa a~věk, \emph{mutabilní} je chování a~názory. Oba směry
      lze rozlišit jedině \emph{longitudinálními} studiemi; studie obezity (Christakis \&
      Fowler, 12\,000 lidí, 32 let) přinesla důkaz pro \emph{sociální vliv}.
\item \textbf{Afiliační sítě} jsou bipartitní a~spojují osoby $\times$ fokusy, zatímco
      sociálně-afiliační síť má rovnou dva typy hran. Uzávěry jsou \textbf{tři}: tranzitivní
      přes společného přítele, \emph{fokální} přes společný fokus, což je \emph{výběr},
      a~\emph{členský}, kdy přítel už ve fokusu je, což je \emph{sociální vliv}. Měří se tak, že
      se $T(k)$ odečte ze dvou snímků sítě a~porovná s~bázovým modelem
      $T_{\text{base}}(k)=1-(1-p)^k$; empiricky vychází velký nárůst z~1 na 2 a~poté
      \emph{klesající výnosy}.
\item \textbf{Schellingův model} rozmístí do mřížky agenty dvou typů. Agent je spokojený,
      jakmile má alespoň $t$ sousedů téhož typu, a~nespokojení se přesouvají, takže vzniká
      \emph{samovolná segregace}. Model ukazuje dvě věci: jak se \emph{imutabilní}
      charakteristika stane silně korelovanou s~\emph{mutabilní}, tady s~volbou místa, a~jak
      lokální homofilie vytvoří globální vzorec.
\item \textbf{Strukturní balance} pokládá za balancované trojúhelníky $+{+}{+}$ a~$+{-}{-}$,
      tedy ty, kde jsou všechny tři vazby pozitivní, nebo právě jedna. \textbf{Věta o~balanci}
      říká, že v~úplném balancovaném grafu jsou buď všichni přátelé, nebo se graf rozdělí na dvě
      skupiny vnitřních přátel, které jsou si navzájem nepřáteli; u~zkoušky se žádá i~důkaz, kde
      $X$ jsou přátelé $A$, $Y$ nepřátelé $A$ a~postupuje se ve~třech krocích. \textbf{Slabá
      balance} zakazuje jedinou konfiguraci, totiž \emph{dvě pozitivní a~jednu negativní},
      a~připouští proto rozdělení na \emph{libovolný počet} skupin. Aplikovala se na evropské
      aliance 1872--1907 a~na Bangladéš 1972.
\item \textbf{HITS} je dotazově \emph{závislý}: z~kořenové množiny $S$ ($t=200$) se vytvoří
      bázová množina $B$ a~pak se mocninnou iterací střídá $a=L^{\!\top}h$ a~$h=La$, přičemž
      normalizace je zde \textbf{nutná}. Posloupnost $a_k=(L^{\!\top}L)^{k-1}a_0$ konverguje
      k~hlavnímu vlastnímu vektoru matice $L^{\!\top}L$, což je \emph{ko-citační} matice,
      kdežto $LL^{\!\top}$ je bibliografické párování. Slabinou je snadné spamování, drift
      tématu a~neefektivnost v~době dotazu.
\item \textbf{PageRank} je naproti tomu dotazově \emph{nezávislý} a~počítá se jako
      $P(i)=\sum_{(j,i)\in E}P(j)/O_j$. Normalizaci \emph{nepotřebuje}, protože celkový
      PageRank zůstává konstantní, jako by šlo o~\uv{kapalinu}. Problém \uv{pomalého úniku} řeší
      \emph{škálované} pravidlo s~tlumicím faktorem $d$, který se volí mezi 0,8--0,9 a~v~článku
      má hodnotu $0{,}85$; obrazně se tomu říká \uv{vodní cyklus}. Markovský řetězec potřebuje
      matici \emph{stochastickou, ireducibilní a~aperiodickou} a~webový graf nesplňuje ani jednu
      z~těchto podmínek: má visící stránky, není silně souvislý a~vykazuje periodicitu. Všechny
      tři nedostatky přitom napraví \emph{jediná} strategie, totiž odkaz z~každé stránky na
      každou s~malou pravděpodobností. Finální tvar je
      $P(i)=(1-d)+d\sum_{(j,i)\in E}P(j)/O_j$ a~řeší se mocninnou iterací. Výhodou je odolnost
      proti spamu, protože vstupní odkazy nelze snadno získat, a~globální offline míra;
      kritizuje se dotazová nezávislost.
\item \textbf{Tematické komunity} se zapisují jako $K=(T,M)$: komunitu určuje téma, komunity se
      mohou \emph{překrývat} a~tvoří hierarchii subkomunit. \emph{Bipartitní jádra} pracují
      s~$(i,j)$-jádrem, kde vějíře jsou huby a~středy autority; postupuje se prořezáním
      a~generováním jader, jenže výsledkem jsou pouze jádra, ne témata ani hierarchie.
      \emph{Komunity pojmenovaných entit} se hledají tak, že se ze společných výskytů ve větě
      postaví graf, v~něm se najdou trojúhelníky a~z~nich jádra, která se nakonec klastrují
      podle textové podobnosti.
\item \textbf{Detekce komunit} naráží hned na začátku na to, že \textbf{definice komunity je
      subjektivní}; rozlišují se skupiny explicitní a~implicitní. \emph{Minimální řez} dává
      nevyvážené skupiny a~stojí $O(n^5)$, s~Kargerovým algoritmem $O(n^4)$; s~obecnými vahami
      a~balančními omezeními je \textbf{NP-těžký}, a~proto se používá \emph{poměrový}
      a~\emph{normalizovaný řez}, které dělí velikostí $|C_i|$, resp.\ objemem
      $\mathrm{vol}(C_i)$.
\item \textbf{Girvanové--Newman} je divizivní algoritmus: postupně odstraňuje hrany s~nejvyšší
      \emph{mezilehlostí}, což stojí $O(m^2n)$, tedy $O(n^3)$. Počet komunit se vybere podle
      nejvyšší \textbf{modularity}
      $Q=\frac{1}{2m}\sum_{i,j}(A_{ij}-\frac{k_ik_j}{2m})\delta(c_i,c_j)$, která nabývá hodnot
      $Q\in[-1,1]$; za smysluplné se považují komunity s~$Q\ge 0{,}3$.
\item \textbf{Louvain} maximalizuje modularitu hladově a~aglomerativně ve dvou fázích: nejdřív
      lokální přesuny, pak stažení do supergrafu, celkem za $O(n\log n)$. Nevýhodou je citlivost
      na pořadí zpracování, možné \emph{nesouvislé komunity} a~uváznutí v~lokálních optimech.
      \textbf{Leiden} proto přidává fázi \emph{upřesnění rozdělení}, pořád ale dává jen tvrdé
      rozdělení. Otevřenou výzvou zůstává neznámý počet komunit a~heterogenní velikosti
      a~hustoty.
\item Metody se dělí do \textbf{čtyř kategorií}. \emph{Uzlově orientované} stojí na klikách:
      hledání kliky je NP-těžké, pomáhá rekurzivní prořezávání uzlů se stupněm $<k-1$
      a~\textbf{CPM} najde i~\emph{překrývající se} komunity přes graf klik, které sdílejí $k-1$
      uzlů. Sem patří $k$-klika, $k$-klub a~$k$-klan, dále $k$-jádro a~$k$-plex, jež jsou
      komplementární, a~také LS množiny a~$k$-komponenty. \emph{Skupinově orientované} metody
      hledají $\gamma$-husté kvazikliky. \emph{Síťově orientované} vycházejí z~podobnosti
      uzlů: Jaccard nebo kosinus v~kombinaci s~$k$-means, dále MDS, blokové modely,
      spektrální shlukování přes grafový laplacián s~\emph{nejmenšími} vlastními čísly
      a~maximalizace modularity přes matici modularity s~\emph{největšími}. Všechny se dají
      sjednotit jako maticová faktorizace $X\approx S\Sigma S^{\!\top}$, která ovšem vyžaduje
      zadat počet komunit. \emph{Hierarchicky orientované} jsou divizivní Girvanové--Newman
      a~aglomerativní Louvain.
\item \textbf{Predikce vazeb} používá buď obsahové míry opřené o~homofilii, nebo strukturní
      míry vycházející z~tranzitivního uzávěru, které jsou silnější. Ze sousedských měr počítají
      \emph{společní sousedé} $|S_i\cap S_j|$, ti ale nezohledňují stupně; \emph{Jaccard}
      $\frac{|S_i\cap S_j|}{|S_i\cup S_j|}$ normalizuje stupně obou uzlů, nikoli však
      prostředních sousedů; \emph{Adamic--Adar} $\sum_{k}\frac{1}{\log|S_k|}$ dává společnému
      sousedovi tím menší váhu, čím vyšší má stupeň. \textbf{Katzova míra}
      $\sum_t\beta^tW_{ij}^{(t)}=(I-\beta A)^{-1}-I$ se hodí pro \emph{řídké} sítě s~málo
      společnými sousedy a~její součet přes ostatní uzly dává \emph{Katzovu centralitu}.
      V~\textbf{klasifikačním pojetí} se úloha formuluje jako pozitivně-neoznačkovaná: nasbírá
      se vzorek negativ, bázovým klasifikátorem bývá logistická regrese a~existují cenově
      citlivé varianty. Výhoda je dvojí, protože klasifikátor sám najde relevanci příznaků
      a~u~orientovaných sítí lze použít \emph{asymetrické} příznaky.
\end{itemize}

%=====================================================================
\newpage
\section{Praktické využití technik}
\label{sec:praxe}
%=====================================================================

Poslední věta okruhu žádá spojení obou polovin dohromady: k~čemu se probrané techniky
\emph{skutečně} používají. Odpověď má u~zkoušky pokaždé stejnou stavbu, takže se vyplatí ji mít
nacvičenou. Ke každé aplikaci se hodí umět říct, \emph{jakou úlohu} řeší
(kap.~\ref{sec:paradigmata}), \emph{jakou metodou} se na ni jde (kap.~\ref{sec:metody})
a~\emph{jak se} její výsledek \emph{vyhodnocuje} (kap.~\ref{sec:reprez}). Aplikace bez těchto
tří údajů zůstává jen historkou.

\subsection{Dobývání znalostí}

Následující přehled spojuje typické aplikace s~úlohou, kterou řeší, a~s~metodou, která se na ni
nasazuje. Třetí sloupec vždy doplňuje, co je na té které aplikaci v~praxi podstatné.

\begin{center}
\small
\begin{tabular}{@{}p{3.5cm}p{4.4cm}p{6.9cm}@{}}
\toprule
\textbf{Aplikace} & \textbf{Typ úlohy / metoda} & \textbf{Poznámka} \\
\midrule
Analýza nákupního košíku & asociační pravidla, Apriori / FP-Growth &
  plánování a~rozvržení prodejny, kupony a~časově omezené slevy, \uv{balíčkování} produktů;
  srovnání nových a~existujících prodejen pomocí \emph{virtuálních položek} \\
Detekce podvodů & klasifikace, detekce odlehlých hodnot &
  20 mil.\ výkonů zdravotní pojišťovny $\to$ 214 tis.\ podivných $\to$ 14 tis.\ předaných
  k~ruční kontrole $\to$ 8 tis.\ potvrzených podvodů; stroj \emph{zaostřuje} pozornost
  experta \\
Kreditní skórování & rozhodovací stromy, naivní Bayes, logistická regrese &
  průběžný příklad v~celém textu; požaduje se \emph{srozumitelnost} (regulace) a~zacházení
  s~chybějícími hodnotami \\
Segmentace zákazníků & klastrová analýza ($k$-means, hierarchické) &
  reprezentace segmentu \emph{centroidem} (\uv{typický zákazník}) je pro byznys velmi
  intuitivní \\
Predikce odchodu (churn) & klasifikace na nevyvážených datech &
  accuracy nestačí $\Rightarrow$ precision/recall, $F_1$, ROC/AUC; \emph{lift analýza} pro
  rozhodnutí, kolik zákazníků kontaktovat \\
Cílené marketingové kampaně & skórování a~řazení, lift &
  ekonomická úvaha nad lift křivkou (příklad s~hodinkami: 1.~přihrádka $+\$550$,
  10.~přihrádka $-\$475$) \\
Klasifikace textu & SVM s~kernely &
  patří k~nejlepším klasifikátorům pro text (vysokorozměrná data); srozumitelnost se
  neočekává \\
Predikce časových řad & regrese, neuronové sítě &
  řada se transformuje posuvným oknem na tabulku vstupů a~výstupů \\
Klasifikace chemických látek & klasifikace strukturních dat &
  reprezentace řetězci SMILES nebo fakty v~Prologu \\
Web usage mining & sekvenční vzory (GSP) &
  navigační vzory uživatelů podle posloupností návštěv stránek \\
Doporučovací systémy \mimo{} & klastrová analýza, klasifikace, predikce vazeb &
  kolaborativní filtrování (podobnost uživatelů či položek) vs.\ obsahové filtrování;
  studený start řeší obsah \\
\bottomrule
\end{tabular}
\end{center}

\textbf{Poznámka k~motivaci rozhodovacích stromů.} Přednáška uvádí případ banky Barclays, která
zavedla telefonní systém rozdělující volající podle úvěrového hodnocení. Klienti s~vysokými
úsporami nebo s~dostatečným úvěrovým limitem byli spojováni s~britským operátorem, všichni
ostatní putovali do indického call centra. Cílem bylo \uv{odfiltrovat} zákazníky, kteří nemají
peníze na nákup produktů banky.

Kritika systému mířila na dvě věci. Zaprvé zvýhodňuje bohatší klienty, zadruhé si vůbec neklade
za cíl poskytnout volajícímu nejlepší radu. Je to dobrá ilustrace toho, že i~\emph{technicky
správný} klasifikační model může mít problematické důsledky, a~zároveň argument pro to, aby fáze
\emph{vyhodnocení} v~CRISP-DM zahrnovala i~pohled mimo přesnost modelu.

\subsection{Analýza sociálních sítí}

Aplikace analýzy sociálních sítí sahají od podnikové komunikace přes bezpečnost až po medicínu.
Přehled je proto řazený podle oblasti, ve které se metoda nasazuje.

\begin{center}
\small
\begin{tabular}{@{}p{3.7cm}p{11.2cm}@{}}
\toprule
\textbf{Oblast} & \textbf{Využití} \\
\midrule
Podniky & analýza a~zlepšení komunikačních toků v~organizaci, s~partnery či zákazníky;
  marketingové kampaně založené na sociálních sítích (vizualizace komunikačních toků před
  a~po zavedení systému pro správu obsahu) \\
Bezpečnostní složky a~armáda & identifikace kriminálních a~teroristických sítí ze zachycené
  komunikace; identifikace \emph{klíčových hráčů} v~nich (míry centrality,
  odd.~\ref{sec:sna}) \\
Sociální sítě (Facebook) & identifikace a~doporučení potenciálních přátel na základě
  \uv{přátel přátel} --- přímo úloha \emph{predikce vazeb} \\
Občanská společnost & odhalování konfliktů zájmů ve skrytých vazbách mezi státními orgány,
  lobbisty a~podniky \\
Síťoví operátoři & optimalizace struktury a~kapacity komunikačních sítí \\
Medicína & nakažlivé nemoci a~jejich prevence; sexuálně přenosné nemoci --- epidemie
  syfilis v~Rockdale County (Atlanta, 1996): skupina bílých dívek okolo 16 let, skupina
  bílých chlapců 17--21 let a~skupina afroamerických chlapců 17--21 let; syfilis
  diagnostikována u~6 bílých žen, 2 bílých mužů a~2 afroamerických mužů \\
Bezškálové sítě & \emph{výpočetní technika:} sítě s~SF architekturou (WWW) jsou extrémně
  odolné proti náhodným poruchám, ale velmi zranitelné koordinovaným útokem --- vymýcení
  internetových virů je v~podstatě nemožné; \emph{medicína:} vakcinační kampaně cílené na
  huby (lidé s~mnoha kontakty --- jejich identifikace je ale obtížná), nové léky mířící na
  klíčové molekuly (huby) dané nemoci s~minimalizací vedlejších účinků pomocí nalezených
  vazebných sítí v~buňkách; \emph{business:} porozumění vazbám mezi firmami, odvětvím
  a~ekonomikou pro monitoring a~eliminaci kaskádových finančních krachů; marketing ---
  zkoumání šíření \uv{nákazy} v~SF sítích a~efektivnější reklama nových produktů \\
Detekce komunit & \emph{Belgický mobilní operátor} (Louvainova metoda): cca 2 miliony
  zákazníků; velikost uzlu odpovídá počtu jednotlivců v~komunitě, barva hlavnímu jazyku
  komunity (červená francouzština, zelená nizozemština); zobrazeny jen komunity nad
  100 zákazníků. Zoom ve vyšším rozlišení odhalí \emph{přechodovou} komunitu tvořenou
  subkomunitami s~méně zřetelným jazykovým oddělením \\
\bottomrule
\end{tabular}
\end{center}

\textbf{Kdy a~proč SNA použít.} Důvodů je několik a~u~zkoušky se hodí mít je srovnané. Nejprostší
z~nich je studium samotné sociální sítě, ať už offline, nebo online, případně snaha o~zvýšení
její efektivity. Druhým je vizualizace: data se zobrazí tak, aby z~nich vystoupily vzorce ve
vztazích či interakcích. Třetím pak sledování cest, kterými v~sociálních sítích proudí informace.

Zvláštní místo má síťová perspektiva v~kvantitativním výzkumu, cenná je ovšem i~pro výzkum
kvalitativní. Rozsah akcí a~příležitostí dostupných jednotlivci je totiž často funkcí jeho
\emph{pozice} v~sociální síti, a~odhalení těchto pozic (např.\ jako otců, matek, učitelů,
pracovníků) může přinést zajímavé a~někdy překvapivé výsledky. Kvantitativní analýza sítě k~tomu
pomáhá identifikovat různé typy aktérů, především klíčové hráče. Poslední skupinou důvodů je
práce se sociálními médii obecně: hledání příčin, proč je komunita nebo síť dysfunkční, a~naopak
podpora sociální koheze a~růstu online komunity.

\begin{poznbox}[Úvahy o~návrhu (\emph{thoughts on design})]
Teorie SNA stála u~zrodu některých z~prvních sociálních sítí, například SixDegrees. Při
návrhových rozhodnutích se přesto dodnes používá spíš zřídka, a~tak zůstává otevřená řada otázek:
\begin{itemize}
\item Jak může online platforma, případně její správci, využít metody a~poznatky SNA?
\item Jak může platforma povzbudit uživatele, aby jednali (\emph{lokálně}) tak, že tím
      rozšíří komunitu?
\item Jaká opatření může online komunita přijmout k~optimalizaci své struktury? Nežádoucí mohou
      být třeba \emph{kliky}, protože odhánějí nově příchozí.
\item Jaké struktury jsou žádoucí pro různé typy online platforem?
\item Jak mohou online komunity identifikovat a~využít klíčové hráče ke svému prospěchu?
\end{itemize}
\end{poznbox}

\subsection{Provozní a~etické aspekty nasazení}

\mimo{} Fází \emph{Deployment} v~CRISP-DM (odd.~\ref{ssec:metodologie}) projekt nekončí. Model
nasazený do provozu žije dál a~přináší si s~sebou vlastní třídu problémů, technických
i~etických:
\begin{itemize}
\item \textbf{Monitoring a~drift.} Rozdělení dat se v~čase mění (\emph{data drift}) a~stejně tak
      se mění i~vztah mezi vstupy a~cílem (\emph{concept drift}). Model proto stárne: je nutné
      pravidelně sledovat pokles jeho kvality na produkčních datech a~přeučovat ho.
\item \textbf{Zpětná vazba modelu do dat.} Nasazený model ovlivňuje tutéž realitu, kterou měří.
      Skórovací model pro schvalování úvěrů vidí jen ty žadatele, kterým úvěr schválil, takže
      další trénovací data jsou zkreslená jím samým.
\item \textbf{Soukromí.} Analýza sociálních sítí pracuje s~daty \emph{o~vztazích}, které
      uživatel často nemohl ovlivnit, a~i~anonymizovaný graf lze podle jeho struktury
      deanonymizovat. Predikce vazeb a~modelování vlivu jsou v~podstatě nástroje pro
      předpovídání a~ovlivňování chování konkrétních lidí.
\item \textbf{Spravedlnost a~diskriminace.} Model naučený na historických datech reprodukuje
      historickou nerovnost, jak ukazuje případ Barclays výše. Vypustit chráněný atribut ze
      vstupů přitom nestačí, protože jej lze často \emph{rekonstruovat} z~ostatních atributů.
\item \textbf{Právní rámec.} GDPR ve svém čl.~22 zakotvuje právo na vysvětlení automatizovaného
      rozhodnutí a~AI Act klade požadavky na vysoce rizikové systémy. V~regulovaných doménách je
      tedy interpretovatelnost právním požadavkem, nikoli pohodlím.
\end{itemize}

\subsection{Shrnutí ke~zkoušce}

\begin{itemize}
\item U~každé aplikace je dobré umět říct \textbf{trojici} údajů: o~jaký typ úlohy jde
      (deskriptivní, nebo prediktivní; s~učitelem, nebo bez učitele), jakou metodou se řeší
      a~jak se výsledek vyhodnocuje.
\item Mezi \textbf{vlajkové aplikace dobývání znalostí} patří analýza nákupního košíku
      (asociační pravidla), detekce podvodů (příklad zdravotní pojišťovny, kde stroj
      \emph{zaostřuje} pozornost experta, ale nenahrazuje jej), kreditní skórování (stromy
      a~Bayes s~požadavkem srozumitelnosti), segmentace zákazníků (klastrování, centroid $=$
      typický zákazník), churn a~cílený marketing (nevyvážená data $\to$ ROC/AUC a~\emph{lift}),
      klasifikace textu (SVM), predikce časových řad (posuvné okno) a~web usage mining
      (sekvenční vzory).
\item Za \textbf{vlajkové aplikace SNA} se dají označit identifikace klíčových hráčů
      v~kriminálních a~teroristických sítích (centralita), doporučování přátel (predikce vazeb),
      optimalizace komunikačních sítí, epidemiologie (Rockdale County 1996), vakcinace cílená na
      \emph{huby}, marketing a~kaskádové finanční krachy v~bezškálových sítích a~segmentace
      zákazníků operátora podle jazykových komunit (Louvain, 2~mil.\ zákazníků).
\item \textbf{Klíčová vlastnost SF sítí pro praxi} je jejich dvojaká odolnost: náhodnému výpadku
      odolají až do 80~\% uzlů, ale cílený útok na 5--15~\% hubů je rozloží. Z~toho plyne jak
      strategie vakcinace, tak nemožnost vymýtit internetové viry.
\item \textbf{Nasazením projekt nekončí.} Model je nutné monitorovat kvůli driftu, počítat se
      zpětnou vazbou jeho rozhodnutí do dat, hlídat soukromí (i~anonymizovaný graf lze
      deanonymizovat) i~spravedlnost (chráněný atribut lze rekonstruovat z~ostatních)
      a~respektovat GDPR čl.~22 a~AI Act.
\end{itemize}

\end{document}
